% !TEX program = LuaLaTeX+se
\pagestyle{empty}

\hspace{0pt}
\vfill

\begin{center}
\uppercase{\textbf{\Large Prayer Book}} \par
\vspace{1em}
\uppercase{\textbf{\normalsize of the Order \\
of St. Paul \\
the First Hermit}} \par
\end{center}

\vfill
\hspace{0pt}

\newpage\null\thispagestyle{empty}\newpage

\hspace{0pt}
\vfill

\begin{center}
\uppercase{\textbf{\huge Prayer Book}} \par
\vspace{1em}
\uppercase{\textbf{\large of the Order\\
of St. Paul \\
the First Hermit}} \par
\end{center}

\begin{center}
\includegraphics[width=0.5\textwidth]{coat-of-arms.png}
\end{center}

\vfill

\begin{center}
Australia 2018
\end{center}

\hspace{0pt}
\pagebreak

\begin{center}
The Prayers \\ of the Order of St. Paul the First Hermit,\\
compiled by the Editorial Team \\ and directed by the decision
of the General Superior,\\ The Most Reverend Arnold Chrapkowski, OSPPE
\end{center}

\vspace{4em}

Editorial Team:
\begin{itemize}
\item Fr. Piotr Polek, OSPPE
\item Fr. Jozef Platek, OSPPE
\item Fr. Stefan Rozej, OSPPE
\item Fr. Ireneusz Pozarowszczyk, OSPPE
\item Fr. Konstancjusz Kunc, OSPPE
\end{itemize}
\medbreak

English Translation for the 1999 edition:
\begin{itemize}
\item Fr. Marian Zalecki, OSPPE
\end{itemize}
\medbreak

Revised English Translation for the 2015 edition:
\begin{itemize}
\item Br. Casimir Zielinski, OSPPE
\end{itemize}
\medbreak

Editor:
\begin{itemize}
\item Br. Joseph Maria Buckley, OSPPE
\end{itemize}

\vfill

\pagebreak
\includeimage{madonna.pdf}\pagestyle{empty}
\hspace{0pt}
\vfill

\textit{``From the very beginning, our Order has enjoyed
the special assistance of Mary. She has always been present
in the life and activity of the Pauline Community. Our Order
has always honoured her as Mother, Lady and Queen and
has striven to propagate (art. 60) special devotion to the
Mother of God; a devotion which expresses itself principally
in the imitation of Mary, in the action of one's personal
sanctification and in the carrying out of apostolic
works.'' (art. 3)}
\medbreak
Almighty God, look kindly upon the prayers your
servants offer you in honour of the Most Holy Virgin Mary,
the Mother of God, and in your mercy grant us your help.
Through Christ Our Lord. Amen.
\vfill
\textit{The Miraculous Image of the Mother of God of Jasna Gora}
\hspace{0pt}
\newpage
\newpageheader{Introduction}{Introduction}

\section*{Introduction}

\setlength{\parskip}{1em}

Prayer forms the basis of the Religious life, yet even more so, for our order, which is called to remain Alone with God Alone in prayer. It is with joy that I write this preface to this new English translation of the \textit{Modlitwenik Zakonu Paulinskiego}. As more vocations are cultivated in Australia the Province desires ever more so to authentically express the Pauline Charism among Australians.

\textit{Lex Orandi, Lex Credendi, Lex Vivendi}, as one prayers, so too does one believe and also live. Therefore, I encourage every member of our province, as well as other English speaking houses who choose to use it, to pray with this prayer book and make it apart of their way of living out our Pauline form of life. \textit{Perfectae Caritatis 6.} reminds us that: ``Drawing therefore upon the authentic sources of Christian spirituality, members of religious communities should resolutely cultivate both the spirit and practice of prayer.''  Also so do we, the Sons of St Paul of Thebes, need to draw on authentic sources of Pauline prayer in our daily lives.

The more closely we can follow Christ in prayer, the more closely can we live out the special gifts and charisms particular to our order. (Cf. Can. 577) In turn, God will bless and make ever more fruitful our mission in the Australian Church, so as to see one day the vision of Fr Augustine Lazur OSPPE flourish with the Order of St. Paul the First Hermit taking deep roots upon the land of Our Lady Help of Christians and under the Southern Cross.

\begin{flushright}
    \includegraphics[width=2in,keepaspectratio]{albert.png}

    Very Reverend Father Albert Wasinowski OSPPE \\
    Provincial of Province of Our Lady Queen of Hermits
\end{flushright}

\section*{Introduction to the Polish Edition}

{\setstretch{1}``Prayer is the fundamental characteristic of our spirituality. In prayer we familiarise ourselves with God's plan for ourselves and for the world, we plunge into the depths of the saving mystery of Christ and unite ourselves with the Father in the Holy Spirit renewing the face of the world.'' (The constitutions and directory of the Order of St Paul the First Hermit, no. 26, art. 26)\par}

{\setstretch{1}The Congregation for Institutes of Consecrated Life and Societies of Apostolic Life in the year 2002, reminded all anew, who through setting foot on the path of consecrated life, determined to give the gift of themselves to the Triune God, begin their religious call, the first enchantment with Christ, which one can never loose. In the then contemporary document of St. John Paul II, we read the following words, which allude to the thought of the saint:\par}

\begin{addmargin}[2em]{2em}
\textit{Every vocation to consecrated life is born in contemplation, from moments of intense communion and from a deep relationship of friendship with Christ, from the beauty and light which was seen shining on his face. From there the desire to always be with the Lord -- and to follow him -- matures: ``how good it is for us to be here'' (Mt 17:4).}
\end{addmargin}

\begin{flushright}
No. 25 starting afresh from Christ: \\ a renewed commitment to consecrated life in the  \\ Third Millennium 2002
\end{flushright}


{\setstretch{1}Gifted with the grace of a vocation in the order of St Paul the First Hermit with greater zeal we desire each day to abide with our Master and Lord -- Jesus Christ. During prayer, let us desire to gaze into his most holy countenance, because the religious life renews itself through the unceasing contemplation of the saviour. In him we, most fully trusting in the intercession of St Paul the First Hermit and in Blessed Eusebius, go out into the world, that through our apostate we may carry the joy of the Gospel to every man.\par}


{\setstretch{1}Just as His disciples gathered in the upper room, let us abide with Mary, the Mother of the risen Lord, in prayer! Let her obtain for us all the grace of unceasing spiritual and apostolic renewal.\par}


\begin{flushright}
    \includegraphics[width=2in,keepaspectratio]{arnold.png}

    Very Reverend Father Arnold Chrapkowski \\
    General of the Order of St. Paul the Frist Hermit
\end{flushright}

\setlength{\parskip}{0em}

\newpage
\newpageheader{Daily Prayers}{Morning Prayers Ordinary Time}

\chapter{Daily Prayers}

\section{Morning Prayers}

\subsection{Ordinary Time}

\begin{samepage}
\prayer{L.}{\textbf{In the name of the Father, and of the Son, and of the Holy Spirit.}}
\prayer{All}{Amen.}
\end{samepage}
\medbreak
\begin{samepage}
\prayer{L.}{\textbf{Eternal Father}}
\prayer{All}{In union with the intentions of your Son, Jesus Christ, animated by the power of the Holy Spirit, calling upon the intercession of the Blessed Virgin Mary, St. Paul, Patriarch of our Order and all the Saints, we offer you all our thoughts, words and actions of this day for your greater glory, for our sanctification and for the salvation of the world.}
\end{samepage}
\medbreak

\label{morningordinaryhtmlangelus}\medbreak
\begin{samepage}
\prayer{L.}{\textbf{The Angel of the Lord declared unto Mary,}}
\prayer{All}{And she conceived of the Holy Spirit}
\end{samepage}
\medbreak

\begin{addmargin}[6em]{0em}\noindent\textit{Hail Mary, full of grace, the Lord is with thee; blessed art thou among women, and blessed is the fruit of thy womb, Jesus. }\end{addmargin}
\medbreak

\begin{addmargin}[6em]{0em}\noindent\textit{Holy Mary, Mother of God, pray for us sinners, now and at the hour of our death. Amen. }\end{addmargin}
\medbreak
\begin{samepage}
\prayer{L.}{\textbf{Behold the handmaid of the Lord,}}
\prayer{All}{Be it done unto me according to thy word.}
\end{samepage}
\medbreak

\begin{addmargin}[6em]{0em}\noindent\textit{Hail Mary\ldots }\end{addmargin}
\medbreak
\begin{samepage}
\prayer{L.}{\textbf{And the Word was made flesh,}}
\prayer{All}{And dwelt among us.}
\end{samepage}
\medbreak

\begin{addmargin}[6em]{0em}\noindent\textit{Hail Mary\ldots }\end{addmargin}
\medbreak
\begin{samepage}
\prayer{L.}{\textbf{Pray for us, O Holy Mother of God.}}
\prayer{All}{That we may be made worthy of the promises of Christ.}
\end{samepage}
\medbreak
\begin{samepage}
\prayer{L.}{\textbf{Let us pray:}\\Pour forth, we beseech you, O Lord, Your grace into our hearts, that we, to whom the Incarnation of Christ, Your Son, was made known by the message of an Angel; may by His Passion and Cross be brought to the glory of His Resurrection. Through the same Christ our Lord.}
\prayer{All}{Amen.}
\end{samepage}
\medbreak
\begin{samepage}
\prayer{L.}{\textbf{Glory be to the Father, and to the Son, and to the Holy Spirit.}}
\prayer{All}{As it was in the beginning, is now, and ever shall be world without end. Amen.}
\end{samepage}
\medbreak
\begin{samepage}
\prayer{L.}{\textbf{Eternal rest grant unto them, O Lord.}}
\prayer{All}{And let perpetual light shine upon them.}
\smallbreak
\prayer{L.}{\textbf{May they rest in peace.}}
\prayer{All}{Amen.}
\end{samepage}
\medbreak

\label{morningordinaryhtmlmaryordinary}\medbreak
\begin{samepage}
\prayer{L.}{\textbf{Rejoice, O Mother of God, Immaculate Virgin}}
\begin{litany}\prayer{All}{Rejoice, you who received glad tidings from an Angel,\par Rejoice, you who gave birth to the Brightness of the Eternal Light,\par Rejoice, O Most Holy Mother of God.\par You alone are inviolate.\par Every creature praises you, O Mother of Light!\par Be for us, we beseech you, therefore,\par our perpetual Intercessor to our Lord, Jesus Christ.}\end{litany}
\end{samepage}
\medbreak
\begin{samepage}
\prayer{L.}{\textbf{In all our tribulations and difficulties.}}
\prayer{All}{Help us, O Most Holy Virgin Mary.}
\end{samepage}
\medbreak
\begin{samepage}
\prayer{L.}{\textbf{Let us pray:}\\Almighty Lord and God, protect us by your power throughout the course of this day, even as you have enabled us to begin it: do not let us turn aside to any sin, but let our every thought, word and deed aim at doing what is pleasing in your sight. Through Christ our Lord.}
\prayer{All}{Amen.}
\end{samepage}
\medbreak
\begin{samepage}
\prayer{L.}{\textbf{Angel of God}}
\prayer{All}{My Guardian dear, to whom God's love commits me here, ever this day be at my side, to enlighten, to guard, to rule and to guide. Amen.}
\end{samepage}
\medbreak

\subsection*{Prayers Before Meditation}

\begin{samepage}
\prayer{L.}{\textbf{Come, Holy Spirit,}}
\prayer{All}{Fill the hearts of your faithful and enkindle in them the fire of your love.}
\smallbreak
\prayer{L.}{\textbf{Send forth your Spirit, and they shall be created.}}
\prayer{All}{And you shall renew the face of the earth.}
\end{samepage}
\medbreak
\begin{samepage}
\prayer{L.}{\textbf{Let us pray:}\\ O God, who did instruct the hearts of your faithful by the light of the Holy Spirit, grant us in the same Spirit to be truly wise and ever to rejoice in his consolation. Through Christ our Lord.}
\prayer{All}{Amen.}
\end{samepage}
\medbreak
\begin{samepage}
\prayer{L.}{\textbf{Queen of Hermits,}}
\prayer{All}{Pray for us.}
\smallbreak
\prayer{L.}{\textbf{Mother of our Order,}}
\prayer{All}{Pray for us.}
\end{samepage}
\medbreak

\instruction{or:}
\medbreak
 \newpage 

\subsubsection*{Prayer of St. Augustine}

\begin{samepage}
\prayer{L.}{\textbf{God, Holy Spirit,}}
\begin{litany}\prayer{All}{Breathe in me,\par That my thoughts may all be holy.\par Act in me, O Holy Spirit,\par That my work, too, may be holy.\par Draw my heart, O Holy Spirit,\par That I may love all that is holy.\par Strengthen me, O Holy Spirit,\par To defend all that is holy.\par Guard me, then, O Holy Spirit,\par That I always may be holy. Amen}\end{litany}
\end{samepage}
\medbreak
\begin{samepage}
\prayer{L.}{\textbf{Queen of Hermits,}}
\prayer{All}{Pray for us.}
\smallbreak
\prayer{L.}{\textbf{Mother of our Order,}}
\prayer{All}{Pray for us.}
\end{samepage}
\medbreak

\subsection*{Prayers After Meditation}

\begin{samepage}
\prayer{L.}{\textbf{We fly to your patronage,}}
\prayer{All}{O Holy Mother of God, despise not our petitions in our necessities, but deliver us always from all dangers. O glorious and blessed Virgin. We humbly ask you to be our Lady, our Intercessor, and our Advocate. Reconcile us to your Son. Recommend us to your Son. Bring us to your Son.}
\end{samepage}
\medbreak

\instruction{or:}
\medbreak
\begin{samepage}
\prayer{L.}{\textbf{O, Mary, Mother of Grace,}}
\prayer{All}{Mother of Mercy, defend us from the enemy and welcome us in the hour of our death. To you, O Jesus, born of the purest Mother and Virgin be glory and honour with the Father and the Holy Spirit. Amen.}
\end{samepage}
\medbreak

\instruction{or:}
\medbreak
 \newpage \medbreak
\begin{samepage}
\prayer{L.}{\textbf{Maria, Mater Gratiae,}}
\begin{litany}\prayer{All}{Mater, Misericordiae,\par Tu nos ab hoste protege,\par et mortis hora suscipe.\par \medbreak Jesu, Tibi sit gloria,\par Qui natus es de Virgine.\par Cum Patre et Almo Spiritu\par in sempiterna saecula. Amen.}\end{litany}
\end{samepage}
\medbreak
\newpage
\newpageheader{Daily Prayers}{Morning Prayers Lent}

\subsection{Lent}

\begin{samepage}
\prayer{L.}{\textbf{In the name of the Father, and of the Son, and of the Holy Spirit.}}
\prayer{All}{Amen.}
\end{samepage}
\medbreak
\begin{samepage}
\prayer{L.}{\textbf{Eternal Father}}
\prayer{All}{In union with the intentions of your Son, Jesus Christ, animated by the power of the Holy Spirit, calling upon the intercession of the Blessed Virgin Mary, St. Paul, Patriarch of our Order and all the Saints, we offer you all our thoughts, words and actions of this day for your greater glory, for our sanctification and for the salvation of the world.}
\end{samepage}
\medbreak

\label{morninglenthtmlangelus}\medbreak
\begin{samepage}
\prayer{L.}{\textbf{The Angel of the Lord declared unto Mary,}}
\prayer{All}{And she conceived of the Holy Spirit}
\end{samepage}
\medbreak

\begin{addmargin}[6em]{0em}\noindent\textit{Hail Mary, full of grace, the Lord is with thee; blessed art thou among women, and blessed is the fruit of thy womb, Jesus. }\end{addmargin}
\medbreak

\begin{addmargin}[6em]{0em}\noindent\textit{Holy Mary, Mother of God, pray for us sinners, now and at the hour of our death. Amen. }\end{addmargin}
\medbreak
\begin{samepage}
\prayer{L.}{\textbf{Behold the handmaid of the Lord,}}
\prayer{All}{Be it done unto me according to thy word.}
\end{samepage}
\medbreak

\begin{addmargin}[6em]{0em}\noindent\textit{Hail Mary\ldots }\end{addmargin}
\medbreak
\begin{samepage}
\prayer{L.}{\textbf{And the Word was made flesh,}}
\prayer{All}{And dwelt among us.}
\end{samepage}
\medbreak

\begin{addmargin}[6em]{0em}\noindent\textit{Hail Mary\ldots }\end{addmargin}
\medbreak
\begin{samepage}
\prayer{L.}{\textbf{Pray for us, O Holy Mother of God.}}
\prayer{All}{That we may be made worthy of the promises of Christ.}
\end{samepage}
\medbreak
\begin{samepage}
\prayer{L.}{\textbf{Let us pray:}\\Pour forth, we beseech you, O Lord, Your grace into our hearts, that we, to whom the Incarnation of Christ, Your Son, was made known by the message of an Angel; may by His Passion and Cross be brought to the glory of His Resurrection. Through the same Christ our Lord.}
\prayer{All}{Amen.}
\end{samepage}
\medbreak
\begin{samepage}
\prayer{L.}{\textbf{Glory be to the Father, and to the Son, and to the Holy Spirit.}}
\prayer{All}{As it was in the beginning, is now, and ever shall be world without end. Amen.}
\end{samepage}
\medbreak
\begin{samepage}
\prayer{L.}{\textbf{Eternal rest grant unto them, O Lord.}}
\prayer{All}{And let perpetual light shine upon them.}
\smallbreak
\prayer{L.}{\textbf{May they rest in peace.}}
\prayer{All}{Amen.}
\end{samepage}
\medbreak

\label{morninglenthtmlmarylent}\medbreak
\begin{samepage}
\prayer{L.}{\textbf{Hail Mary, Mother of God!}}
\begin{litany}\prayer{All}{You, in your innocent womb, have embraced the Infinite.\par Thanks to you, we worship the Most Holy Trinity.\par Thanks to you, the Angels and Archangels rejoice.\par Thanks to you, the powers of evil take flight.\par Thanks to you, the people turn to penance.\par Thanks to you, the only Begotten Son of God has shown the Light to those who walk in darkness and in the shadow of death,\par Thanks to you, the Apostles announced salvation to all nations.\par Thanks to you, the kings reign by the power of the Holy Trinity.}\end{litany}
\end{samepage}
\medbreak
\begin{samepage}
\prayer{L.}{\textbf{In all our tribulations and difficulties.}}
\prayer{All}{Help us, O Most Holy Virgin Mary.}
\end{samepage}
\medbreak
\begin{samepage}
\prayer{L.}{\textbf{Let us pray:}\\Almighty Lord and God, protect us by your power throughout the course of this day, even as you have enabled us to begin it: do not let us turn aside to any sin, but let our every thought, word and deed aim at doing what is pleasing in your sight. Through Christ our Lord.}
\prayer{All}{Amen.}
\end{samepage}
\medbreak
\begin{samepage}
\prayer{L.}{\textbf{Angel of God}}
\prayer{All}{My Guardian dear, to whom God's love commits me here, ever this day be at my side, to enlighten, to guard, to rule and to guide. Amen.}
\end{samepage}
\medbreak

\subsection*{Prayers Before Meditation}

\begin{samepage}
\prayer{L.}{\textbf{Come, Holy Spirit,}}
\prayer{All}{Fill the hearts of your faithful and enkindle in them the fire of your love.}
\smallbreak
\prayer{L.}{\textbf{Send forth your Spirit, and they shall be created.}}
\prayer{All}{And you shall renew the face of the earth.}
\end{samepage}
\medbreak
\begin{samepage}
\prayer{L.}{\textbf{Let us pray:}\\ O God, who did instruct the hearts of your faithful by the light of the Holy Spirit, grant us in the same Spirit to be truly wise and ever to rejoice in his consolation. Through Christ our Lord.}
\prayer{All}{Amen.}
\end{samepage}
\medbreak
\begin{samepage}
\prayer{L.}{\textbf{Queen of Hermits,}}
\prayer{All}{Pray for us.}
\smallbreak
\prayer{L.}{\textbf{Mother of our Order,}}
\prayer{All}{Pray for us.}
\end{samepage}
\medbreak

\instruction{or:}
\medbreak

\subsubsection*{Prayer of St. Augustine}

\begin{samepage}
\prayer{L.}{\textbf{God, Holy Spirit,}}
\begin{litany}\prayer{All}{Breathe in me,\par That my thoughts may all be holy.\par Act in me, O Holy Spirit,\par That my work, too, may be holy.\par Draw my heart, O Holy Spirit,\par That I may love all that is holy.\par Strengthen me, O Holy Spirit,\par To defend all that is holy.\par Guard me, then, O Holy Spirit,\par That I always may be holy. Amen}\end{litany}
\end{samepage}
\medbreak
\begin{samepage}
\prayer{L.}{\textbf{Queen of Hermits,}}
\prayer{All}{Pray for us.}
\smallbreak
\prayer{L.}{\textbf{Mother of our Order,}}
\prayer{All}{Pray for us.}
\end{samepage}
\medbreak

\subsection*{Prayers After Meditation}

\begin{samepage}
\prayer{L.}{\textbf{We fly to your patronage,}}
\prayer{All}{O Holy Mother of God, despise not our petitions in our necessities, but deliver us always from all dangers. O glorious and blessed Virgin. We humbly ask you to be our Lady, our Intercessor, and our Advocate. Reconcile us to your Son. Recommend us to your Son. Bring us to your Son.}
\end{samepage}
\medbreak

\instruction{or:}
\medbreak
\begin{samepage}
\prayer{L.}{\textbf{O, Mary, Mother of Grace,}}
\prayer{All}{Mother of Mercy, defend us from the enemy and welcome us in the hour of our death. To you, O Jesus, born of the purest Mother and Virgin be glory and honour with the Father and the Holy Spirit. Amen.}
\end{samepage}
\medbreak

\instruction{or:}
\medbreak
\begin{samepage}
\prayer{L.}{\textbf{Maria, Mater Gratiae,}}
\begin{litany}\prayer{All}{Mater, Misericordiae,\par Tu nos ab hoste protege,\par et mortis hora suscipe.\par \medbreak Jesu, Tibi sit gloria,\par Qui natus es de Virgine.\par Cum Patre et Almo Spiritu\par in sempiterna saecula. Amen.}\end{litany}
\end{samepage}
\medbreak
\newpage
\newpageheader{Daily Prayers}{Morning Prayers Easter}

\subsection{Easter}

\begin{samepage}
\prayer{L.}{\textbf{In the name of the Father, and of the Son, and of the Holy Spirit.}}
\prayer{All}{Amen.}
\end{samepage}
\medbreak
\begin{samepage}
\prayer{L.}{\textbf{Eternal Father}}
\prayer{All}{In union with the intentions of your Son, Jesus Christ, animated by the power of the Holy Spirit, calling upon the intercession of the Blessed Virgin Mary, St. Paul, Patriarch of our Order and all the Saints, we offer you all our thoughts, words and actions of this day for your greater glory, for our sanctification and for the salvation of the world.}
\end{samepage}
\medbreak

\label{morningeasterhtmlreginacaeli}\medbreak

\instruction{From Easter Sunday to the solemnity of Pentecost instead of the "Angelus" we recite the following antiphon:}
\medbreak
\begin{samepage}
\prayer{L.}{\textbf{Queen of heaven, rejoice, alleluia.}}
\begin{litany}\prayer{All}{For he, whom you were worthy to bear, alleluia,\par has risen, as He said, alleluia.\par Pray for us to God, alleluia.}\end{litany}
\end{samepage}
\medbreak
\begin{samepage}
\prayer{L.}{\textbf{Rejoice and be glad, O Virgin Mary, alleluia.}}
\prayer{All}{For the Lord has truly risen, alleluia.}
\end{samepage}
\medbreak
\prayer{L.}{\textbf{Let us pray:}\\God our Father, you have given joy to the world by the resurrection of your Son Jesus Christ. Through the prayers of his Mother, the Virgin Mary, bring us to the happiness of eternal life. Through Christ our Lord.}
\prayer{All}{Amen.}
\medbreak
\begin{samepage}
\prayer{L.}{\textbf{Glory be to the Father, and to the Son, and to the Holy Spirit.}}
\prayer{All}{As it was in the beginning, is now, and ever shall be world without end. Amen.}
\end{samepage}
\medbreak
\begin{samepage}
\prayer{L.}{\textbf{Eternal rest grant unto them, O Lord.}}
\prayer{All}{And let perpetual light shine upon them.}
\smallbreak
\prayer{L.}{\textbf{May they rest in peace.}}
\prayer{All}{Amen.}
\end{samepage}
\medbreak

\label{morningeasterhtmlmaryordinary}\medbreak
\begin{samepage}
\prayer{L.}{\textbf{Rejoice, O Mother of God, Immaculate Virgin}}
\begin{litany}\prayer{All}{Rejoice, you who received glad tidings from an Angel,\par Rejoice, you who gave birth to the Brightness of the Eternal Light,\par Rejoice, O Most Holy Mother of God.\par You alone are inviolate.\par Every creature praises you, O Mother of Light!\par Be for us, we beseech you, therefore,\par our perpetual Intercessor to our Lord, Jesus Christ.}\end{litany}
\end{samepage}
\medbreak
\begin{samepage}
\prayer{L.}{\textbf{In all our tribulations and difficulties.}}
\prayer{All}{Help us, O Most Holy Virgin Mary.}
\end{samepage}
\medbreak
\begin{samepage}
\prayer{L.}{\textbf{Let us pray:}\\Almighty Lord and God, protect us by your power throughout the course of this day, even as you have enabled us to begin it: do not let us turn aside to any sin, but let our every thought, word and deed aim at doing what is pleasing in your sight. Through Christ our Lord.}
\prayer{All}{Amen.}
\end{samepage}
\medbreak
\begin{samepage}
\prayer{L.}{\textbf{Angel of God}}
\prayer{All}{My Guardian dear, to whom God's love commits me here, ever this day be at my side, to enlighten, to guard, to rule and to guide. Amen.}
\end{samepage}
\medbreak

\subsection*{Prayers Before Meditation}

\begin{samepage}
\prayer{L.}{\textbf{Come, Holy Spirit,}}
\prayer{All}{Fill the hearts of your faithful and enkindle in them the fire of your love.}
\smallbreak
\prayer{L.}{\textbf{Send forth your Spirit, and they shall be created.}}
\prayer{All}{And you shall renew the face of the earth.}
\end{samepage}
\medbreak
\begin{samepage}
\prayer{L.}{\textbf{Let us pray:}\\ O God, who did instruct the hearts of your faithful by the light of the Holy Spirit, grant us in the same Spirit to be truly wise and ever to rejoice in his consolation. Through Christ our Lord.}
\prayer{All}{Amen.}
\end{samepage}
\medbreak
\begin{samepage}
\prayer{L.}{\textbf{Queen of Hermits,}}
\prayer{All}{Pray for us.}
\smallbreak
\prayer{L.}{\textbf{Mother of our Order,}}
\prayer{All}{Pray for us.}
\end{samepage}
\medbreak

\instruction{or:}
\medbreak

\subsubsection*{Prayer of St. Augustine}

\begin{samepage}
\prayer{L.}{\textbf{God, Holy Spirit,}}
\begin{litany}\prayer{All}{Breathe in me,\par That my thoughts may all be holy.\par Act in me, O Holy Spirit,\par That my work, too, may be holy.\par Draw my heart, O Holy Spirit,\par That I may love all that is holy.\par Strengthen me, O Holy Spirit,\par To defend all that is holy.\par Guard me, then, O Holy Spirit,\par That I always may be holy. Amen}\end{litany}
\end{samepage}
\medbreak
\begin{samepage}
\prayer{L.}{\textbf{Queen of Hermits,}}
\prayer{All}{Pray for us.}
\smallbreak
\prayer{L.}{\textbf{Mother of our Order,}}
\prayer{All}{Pray for us.}
\end{samepage}
\medbreak

\subsection*{Prayers After Meditation}

\begin{samepage}
\prayer{L.}{\textbf{We fly to your patronage,}}
\prayer{All}{O Holy Mother of God, despise not our petitions in our necessities, but deliver us always from all dangers. O glorious and blessed Virgin. We humbly ask you to be our Lady, our Intercessor, and our Advocate. Reconcile us to your Son. Recommend us to your Son. Bring us to your Son.}
\end{samepage}
\medbreak

\instruction{or:}
\medbreak
\begin{samepage}
\prayer{L.}{\textbf{O, Mary, Mother of Grace,}}
\prayer{All}{Mother of Mercy, defend us from the enemy and welcome us in the hour of our death. To you, O Jesus, born of the purest Mother and Virgin be glory and honour with the Father and the Holy Spirit. Amen.}
\end{samepage}
\medbreak

\instruction{or:}
\medbreak
\begin{samepage}
\prayer{L.}{\textbf{Maria, Mater Gratiae,}}
\begin{litany}\prayer{All}{Mater, Misericordiae,\par Tu nos ab hoste protege,\par et mortis hora suscipe.\par \medbreak Jesu, Tibi sit gloria,\par Qui natus es de Virgine.\par Cum Patre et Almo Spiritu\par in sempiterna saecula. Amen.}\end{litany}
\end{samepage}
\medbreak
\newpage
\newpageheader{Daily Prayers}{Daily Consecration to Our Lady of Jasna Gora}

\subsection{Daily Consecration to Our Lady of Jasna Gora}

\begin{samepage}
\prayer{L.}{\textbf{Holy Mother of Jasna Gora,}}
\prayer{All}{Mediatrix of all graces, you are full of grace, goodness and mercy. I consecrate to you all my thoughts, words and actions, my soul and body. I ask your help and especially your prayers for my salvation. Today I consecrate myself to you, Holy Mother, totally, with body and soul amid joys and sufferings to obtain for myself and others your help on this earth and eternal life in heaven. Amen.}
\end{samepage}
\medbreak
\newpageheader{Daily Prayers}{Afternoon Prayers}

\section[Afternoon Prayers]{Afternoon Prayers \\ {\small\normalfont Adoration of the Most Blessed Sacrament}}

\instruction{After the meal, we process to the Monastery Chapel and sing a psalm, a song or any appropriate hymn, related to the liturgical season. The Prior starts the singing, or the Subprior, or his substitute, either a cantor or the antiphonarian. The adoration of the Blessed Sacrament should begin with an Eucharistic hymn. The recitation of the ``Angelus'' (if it was not said earlier) follows. After, there is a foreseen time for a particular examination of conscience, a prayerful reflection (about 3 minutes) follows. Further, in accordance with norm 72 of the Directory, we pray for vocations to our Order and for other actual intentions. At the end of adoration it has been proposed to pray a litany, for example:}
\medbreak

\noindent\begin{tabu} to \textwidth { X[-1,l] X[l]}
Monday & Invocations to St. Paul the First Hermit p.~\pageref{saintpaulhtmllitany}.\\
Tuesday & Invocations to the Guardian Angels p.~\pageref{invocationsangelshtmlinvocations}.\\
Wednesday & Litany to St. Joseph p.~\pageref{saintjosephhtmllitany}.\\
Thursday & Litany to the Holy Name of Jesus p.~\pageref{litanieshtmlholynamejesus} or to the Sacred Heart of Jesus p.~\pageref{litanieshtmlsacredheart}.\\
Friday & Litany to the Sacred Blood of Jesus p.~\pageref{litanieshtmlbloodjesus} or Litany to the Five wounds of the Lord Jesus Christ p.~\pageref{litanieshtmlwoundsjesus}.\\
Saturday & Litany to the Mother of God p.~\pageref{loretohtmllitany}.
\end{tabu}\medbreak
\newpage
\newpageheader{Daily Prayers}{Second Meditation}

\section{Second Meditation}

\instruction{At the appropriate time, according to art. 28, a second meditation should occur which can be substituted by another form of prayer or devotion. Before and after the meditation some suitable prayer should be said, for example:}
\medbreak

\instruction{Before meditation:}
\medbreak
\begin{samepage}
\prayer{L.}{\textbf{My Lord and my God!}}
\prayer{All}{I strongly believe that you are present here and that your eyes are turned upon me. I fall on my knees knowing that I am unworthy to stand before your Face but full of confidence in your infinite goodness, I humbly implore your grace that I may make this meditation for your greater glory and my spiritual benefit. Illuminate my mind, move my heart, strengthen my will, that I may know you better, love you more and serve you more faithfully. I beseech this through the intercession of the Virgin Mary, St. Paul -- the Father of our Order -- and my Guardian Angel. Amen.}
\end{samepage}
\medbreak
\begin{samepage}
\prayer{L.}{\textbf{Queen of Hermits,}}
\prayer{All}{Pray for us.}
\smallbreak
\prayer{L.}{\textbf{Mother of our Order,}}
\prayer{All}{Pray for us.}
\end{samepage}
\medbreak

\instruction{After meditation:}
\medbreak
\begin{samepage}
\prayer{L.}{\textbf{O God,}}
\prayer{All}{I thank you for the light and graces granted me during this meditation. I resolve to remain attentive to your will and to grow in your love. Renew me that I may be more aware of your voice and ready to assist my neighbours in their need. O Mary, into your hands I entrust my petition for the grace of perseverance in my resolutions and in love. Amen.}
\end{samepage}
\medbreak
\newpage
\newpageheader{Daily Prayers}{Evening Prayers}

\section{Evening Prayers}

\instruction{Vespers should be recited at the time stated in the daily monastery horarium. Followed by the Invocations to St. Paul. Next, the Superior recalls the names of the deceased Brothers for the following day and concludes with a prayer for the dead. Depending on the daily horarium the ``Angelus'' should be recited at the beginning of the Evening Prayers or after supper.}
\medbreak

\subsection*{Invocations to St. Paul the First Hermit}

\begin{litany}Saint Paul, Patriarch of our Order, \responce{Pray for us.}\par 
Saint Paul, faithful follower of Christ, \par
Saint Paul, man of strong and lively faith, \par
Saint Paul, lover of solitude, \par
Saint Paul, persevering in prayer and penance, \par
Saint Paul, patient in bearing the difficulties of the eremitical life, \par
Saint Paul, embracing with your devoted love the Church of Christ and all of mankind, \par
Saint Paul, model of perseverance in your chosen way of life, \par
Saint Paul, prepared to meet God at the time of death, \par
\end{litany}
\medbreak
\begin{litany}That we may be faithful to our monastic vocation, \responce{Intercede to God for us.}\par 
That through the observance of our religious vows, we may become more Christ-like, \par
That we may fulfil Christ's command to love all men, \par
That we may pray well and unceasingly, \par
That we may participate in the saving love of Christ through a penitential life, \par
That we may not succumb to the temptation of mediocrity, \par
That we may prudently profit from the fruits of the land and from technical and cultural achievements, \par
That we may fearlessly strive to attain that love exemplified in the Gospel, \par
That we may feel responsible for our Order and its apostolic works, \par
That we may persevere in our Order until death, \par
\end{litany}
\medbreak
\begin{litany}Through the intercession of the Blessed Virgin Mary, Queen of Hermits and Patroness of our Order, \responce{We beg you Lord, hear us!}\par 
Through the merits of all the Saints of our Order, \par
That you may graciously grant our religious community an increase in number and an enrichment of virtues, \par
That you may generously provide graces for our families, confraters and benefactors, \par
That you may show your mercy to our brothers, who have left our community, \par
That you may remedy the failures and defects of our love, \par
That you may grant the grace of conversion to the enemies of the Church, \par
That you may grant eternal happiness to our deceased brothers, parents, relatives, confraters and benefactors, \par
\end{litany}
\medbreak
\begin{samepage}
\prayer{L.}{\textbf{Pray for us St. Paul the First Hermit,}}
\prayer{All}{That we may be made worthy of the promises of Christ.}
\end{samepage}
\medbreak
\begin{samepage}
\prayer{L.}{\textbf{Let us pray:}\\ God, you granted that Saint Paul, our Father, achieve sanctity in the desert, grant through his intercession, that we who are striving to develop in ourselves a spirit of prayer and service may come ever closer to you in love. Through Christ our Lord.}
\prayer{All}{Amen.}
\end{samepage}
\medbreak

\instruction{In the prayer for the deceased we use the following form:}
\medbreak
\prayer{L.}{\textbf{Our deceased Brothers whose anniversary of death will occur tomorrow\ldots}}

\label{eveninghtmlnecrology}
\label{eveninghtmlnecrologynotfound}\smallbreak
\prayer{L.}{\textbf{Let us pray for our deceased Brothers}}
\medbreak

\instruction{(a short silence)}
\medbreak
\begin{samepage}
\prayer{L.}{\textbf{O God, giver of forgiveness and lover of human salvation, we implore your goodness, grant that our Brothers, Parents, Confraters and Benefactors, who have departed from this life, may, through the intercession of the Most Holy Virgin Mary and all your Saints, be admitted to eternal happiness. We ask this through Christ our Lord.}}
\prayer{All}{Amen.}
\end{samepage}
\medbreak
\begin{samepage}
\prayer{L.}{\textbf{Eternal rest grant unto them, O Lord.}}
\prayer{All}{And let perpetual light shine upon them.}
\smallbreak
\prayer{L.}{\textbf{May they rest in peace.}}
\prayer{All}{Amen.}
\end{samepage}
\medbreak

\instruction{At the time of the notice of a Brother's death, this prayer is added:}
\medbreak
\begin{samepage}
\prayer{L.}{\textbf{O Lord, God of mercy, we pray that you may welcome your deceased Servant and our Brother \rubric{N. N.}, that he may share in the glory of Your Resurrection with you and your Saints,}}
\prayer{All}{Amen.}
\end{samepage}
\medbreak
\begin{samepage}
\prayer{L.}{\textbf{Eternal rest grant unto them, O Lord.}}
\prayer{All}{And let perpetual light shine upon them.}
\smallbreak
\prayer{L.}{\textbf{May they rest in peace.}}
\prayer{All}{Amen.}
\end{samepage}
\medbreak
\newpage
\newpageheader{Daily Prayers}{Night Prayer}

\section{Night Prayer}

\instruction{This is the last Hour of the Breviary recited before repose. A brief time should be given for an examination of conscience (about 3 minutes) as foreseen by norm 27 of the Directory. The following points should be taken into consideration:}
\medbreak
\begin{itemize} \item Let us place ourselves in the presence of God and thank Him for all graces received, particularly for those that were granted today. \item Let us ask the Lord God for the grace to know our sins and to turn away from them. \item Let us call to mind all sins we have committed today in our thoughts, in our speech, in our actions, by our omissions but especially the sins which we are inclined to commit more often. Let us pay attention to any neglect of the resolutions, made during the morning meditation. \item Let us make a sincere act of contrition for our sins and promise God to avoid them, with the help of His grace. \item Let us live in the state, in which we wish to be at the hour of death. \end{itemize}\newpageheader{Daily Prayers}{Prayer Before Repose}

\section{Prayer Before Repose}

\instruction{Before we retire, following our custom, we are to pray in our private room, on our knees with hands raised in the form of the cross, while reciting three times, the ``Hail Mary''. After each ``Hail Mary'' we are to repeat the following invocation to the Most Holy Virgin Mary:}
\medbreak

Through your Holy and Immaculate Conception, O Most Holy Virgin Mary, make my heart pure and my soul holy.
\medbreak

\instruction{Then to honour St. Paul we pray this additional prayer}
\medbreak

O Lord God Almighty, who deigned to detach St. Paul the First Hermit from the noisy world and in an admirable way joined him to the Angelic choirs, I ask you to purify me of all earthly profanity and permeate me with heavenly gifts. Through Christ, Our Lord. Amen.
\medbreak
\newpage
\newpageheader{Daily Prayers}{Prayers in the Refectory}

\section{Prayers in the Refectory}

\subsection{Prayer before meals}

\begin{samepage}
\prayer{L.}{\textbf{In the name of the Father and of the Son and of the Holy Spirit.}}
\prayer{All}{Amen}
\end{samepage}
\medbreak
\begin{samepage}
\prayer{L.}{\textbf{We adore you, Lord who call us to the feast of the Lamb.}}
\prayer{All}{You gather us around your table and generously endow us with your gifts.}
\end{samepage}
\medbreak
\begin{samepage}
\prayer{L.}{\textbf{Our eyes turn trustingly towards you.}}
\prayer{All}{And you give us food at the due time.}
\end{samepage}
\medbreak
\begin{samepage}
\prayer{L.}{\textbf{Our Father \ldots}}
\end{samepage}
\medbreak

\paragraph*{Sunday}

\begin{samepage}
\prayer{L.}{O Holy Trinity, One God, bless \cross{\GreDagger} us and these gifts and those who have prepared them. Teach us to share our bread and joy with all. You who live and reign for ever and ever.}
\prayer{All}{Amen}
\end{samepage}
\medbreak

\paragraph*{Monday}

\begin{samepage}
\prayer{L.}{O God, Almighty Father, bless \cross{\GreDagger} us and this food to strengthen our bodies to do your will. Through Christ our Lord.}
\prayer{All}{Amen}
\end{samepage}
\medbreak

\paragraph*{Tuesday}

\begin{samepage}
\prayer{L.}{O Lord, our God, bless \cross{\GreDagger} us and these gifts which we will enjoy from your holy Providence. Through Christ our Lord.}
\prayer{All}{Amen}
\end{samepage}
\medbreak

\paragraph*{Wednesday}

\begin{samepage}
\prayer{L.}{O Lord, our God, bless \cross{\GreDagger} us and this food and grant that we may participate in the heavenly banquet. Through Christ our Lord.}
\prayer{All}{Amen}
\end{samepage}
\medbreak

\paragraph*{Thursday}

\begin{samepage}
\prayer{L.}{O Lord, who has revealed yourself to your Disciples in the breaking of bread, bless \cross{\GreDagger} us and these gifts and sanctify us with your presence who live and reign for ever and ever.}
\prayer{All}{Amen}
\end{samepage}
\medbreak

\paragraph*{Friday}

\begin{samepage}
\prayer{L.}{O God, our Father, bless \cross{\GreDagger} us and these gifts and those who have prepared them. Grant that we may be brothers at table and in daily life. Through Christ our Lord.}
\prayer{All}{Amen}
\end{samepage}
\medbreak

\paragraph*{Saturday}

\begin{samepage}
\prayer{L.}{O God, Father of all goodness, bless \cross{\GreDagger} us and these gifts with which you sustain our strength and life. Grant that this communal meal will strengthen us in mutual love. Through Christ our Lord.}
\prayer{All}{Amen}
\end{samepage}
\medbreak

\paragraph*{Solemnity}

\begin{samepage}
\prayer{L.}{O Lord, our God, bless \cross{\GreDagger} us and this food and those who have prepared it. Teach us to share our bread with others and the joy that flows from the mysteries of the feast we celebrate today. Through Christ our Lord.}
\prayer{All}{Amen}
\end{samepage}
\medbreak

\paragraph*{Advent Season}

\begin{samepage}
\prayer{L.}{O God, our Heavenly Father, bless \cross{\GreDagger} us and these gifts that are the signs of your love for us and announce the birth of Jesus Christ --- the Light that illuminates every man, who lives and reigns for ever and ever.}
\prayer{All}{Amen}
\end{samepage}
\medbreak

\paragraph*{Christmas Vigil Supper}

\instruction{All the Community gathers in the refectory. At the beginning, the Superior or a Cantor sings an appropriate Christmas carol. During the singing all lights should be lit. After the song, a part of the Holy Scripture should be read that relates the story of Jesus Christ's birth (e.g. Lk 2:1-20). Next, the Superior says a prayer, as designated below, expresses his wishes to the Community and invites those present to exchange their Christmas wishes individually, while breaking the wafer. During the Christmas vigil meal we sing carols.}
\medbreak
\begin{samepage}
\prayer{L.}{O Lord, Almighty God, bless \cross{\GreDagger} us, this holy wafer and this Christmas meal. Give us the grace to understand your love, your mercy and the goodness you have shown through the Nativity of your Son, Jesus Christ. We commend to you all the people of the world who as a consequence of famine, war, hatred and suffering cannot fully experience your love. Open our eyes and our hearts to all who expect of us our love, especially the poor, the sick, the abandoned and also our parents, relatives and benefactors. May this sacred Christmas Vigil inspire us so that in our behaviour everyone may recognize that you, O God, are the Father of us all. Through Christ our Lord.}
\prayer{All}{Amen}
\end{samepage}
\medbreak

\paragraph*{Christmas Season}

\begin{samepage}
\prayer{L.}{O God, the Almighty Father, who through the Nativity of your Son, Jesus Christ, have sanctified family life, bless \cross{\GreDagger} us and this meal and grant that we may with joy celebrate the presence of the Word that has become Flesh and dwells among us. Who lives and reigns for ever and ever.}
\prayer{All}{Amen}
\end{samepage}
\medbreak

\paragraph*{Lenten Season}

\begin{samepage}
\prayer{L.}{Lord, Jesus Christ, bless \cross{\GreDagger} us and this meal. Grant that we may fast with cheerful hearts. Grant that remembering your passion and death upon the cross we may look, not for perishable food, but for food that lasts for ever. You, who live and reign for ever and ever.}
\prayer{All}{Amen}
\end{samepage}
\medbreak

\paragraph*{Resurrection Sunday}

\instruction{At a proper time (at the Dawn Vigil of the Resurrection of the Lord), the Superior or the Monastery Administrator, wearing a surplice and stole, blesses the Easter food, saying this prayer:}
\medbreak
\begin{samepage}
\prayer{L.}{Lord, Jesus Christ, we joyfully glorify you in your Resurrection and in your revelation to your Disciples through the breaking of the Bread. Now we ask you to bless \cross{\GreDagger} us and these Easter gifts. Grant that receiving them with our Brothers we may become sharers of the heavenly feast in your Kingdom, who lives and reigns for ever and ever.}
\prayer{All}{Amen}
\end{samepage}
\medbreak

\paragraph*{Easter Season}

\begin{samepage}
\prayer{L.}{Lord, our God, you have conquered death and opened the gates of eternal life. Bless \cross{\GreDagger} us and this meal, that celebrating it with gratitude we may become worthy to sit at the table in the Kingdom of your Risen Son, who lives and reigns with you for ever and ever.}
\prayer{All}{Amen}
\end{samepage}
\medbreak

\paragraph*{Pentecost}

\begin{samepage}
\prayer{L.}{O God, Almighty Father, through the mystery of this day you sanctify your Church, we ask you to bless \cross{\GreDagger} us and this meal and send the Gifts of the Holy Spirit to all the earth. Accomplish in our hearts your faithful signs of unity, love and peace which you granted to your disciples gathered in the Upper Room. Through Christ our Lord.}
\prayer{All}{Amen}
\end{samepage}
\medbreak

\subsection{Prayer after meals}

\begin{samepage}
\prayer{L.}{\textbf{Blessed is God in his gifts.}}
\prayer{All}{And Holy in all his works. His are power and glory for ever.}
\end{samepage}
\medbreak
\begin{samepage}
\prayer{L.}{\textbf{Glorify the Lord with me.}}
\prayer{All}{Together, let us praise his Name.}
\end{samepage}
\medbreak
\begin{samepage}
\prayer{L.}{\textbf{Glory be to the Father and to the Son and to the Holy Spirit,}}
\prayer{All}{As it was in the beginning, is now and ever shall be, world without end. Amen}
\end{samepage}
\medbreak

\paragraph*{Sunday}

\begin{samepage}
\prayer{L.}{Heavenly Father, we are grateful for this meal which we have celebrated thanks to your goodness. We ask you to endow with it those who go hungry this day and that you will bestow upon us the desire of the heavenly gifts. Through Christ our Lord.}
\prayer{All}{Amen}
\end{samepage}
\medbreak

\paragraph*{Monday}

\begin{samepage}
\prayer{L.}{Father of infinite goodness be glorified for your kindness. Through Christ our Lord.}
\prayer{All}{Amen}
\end{samepage}
\medbreak

\paragraph*{Tuesday}

\begin{samepage}
\prayer{L.}{We thank you Lord, for the sustenance you have given us. Grant that we may willingly and resolutly serve you and our neighbours. Who lives and reigns for ever and ever.}
\prayer{All}{Amen}
\end{samepage}
\medbreak

\paragraph*{Wednesday}

\begin{samepage}
\prayer{L.}{God, our Father, through your loving providence you generously feed everything that lives, accept our gratitude for today's meal. Through Christ our Lord.}
\prayer{All}{Amen}
\end{samepage}
\medbreak

\paragraph*{Thursday}

\begin{samepage}
\prayer{L.}{O Lord our God, we thank you for the goods which you bestow upon us, may your Blessing remain with us always. Through Christ our Lord.}
\prayer{All}{Amen}
\end{samepage}
\medbreak

\paragraph*{Friday}

\begin{samepage}
\prayer{L.}{O God, source of all goodness, strengthened by this food of your kindness, we go willingly to do your will. You who live and reign for ever and ever.}
\prayer{All}{Amen}
\end{samepage}
\medbreak

\paragraph*{Saturday}

\begin{samepage}
\prayer{L.}{God our Father, who through the assistance of Mary and Joseph, nourished your Son Jesus Christ. Bless all those who prepared this meal and make us sensitive to the needs of others. Through Christ our Lord.}
\prayer{All}{Amen}
\end{samepage}
\medbreak

\paragraph*{Solemnity}

\begin{samepage}
\prayer{L.}{O Lord our God, kindly bless all, who in the spirit of your love hasten to our aid. Through Christ our Lord.}
\prayer{All}{Amen}
\end{samepage}
\medbreak

\paragraph*{Advent Season}

\begin{samepage}
\prayer{L.}{O Lord Jesus, Master and Redeemer of the world, whose coming we are awaiting, we thank you for your presence during this meal. Grant that all people may know you in the mystery of your Nativity and celebrate the approaching solemnity in joy and peace. Who lives for ever and ever.}
\prayer{All}{Amen}
\end{samepage}
\medbreak

\paragraph*{Christmas Vigil Supper}

\begin{samepage}
\prayer{L.}{Lord our God, you granted that this most Holy Night shine for us with a radiance of true light. We beseech you that we may participate in the joys of the heavenly banquet of Christ, whose Mystery of the Incarnation we have come to know on earth. For you live for ever and ever.}
\prayer{All}{Amen}
\end{samepage}
\medbreak

\paragraph*{Christmas Season}

\begin{samepage}
\prayer{L.}{God our Father, grateful for the gift of your Only Begotten Son, who has come into the world to give us eternal life. We also thank you for the gifts with which you sustain our temporal life. Through Christ our Lord.}
\prayer{All}{Amen}
\end{samepage}
\medbreak

\paragraph*{Lenten Season}

\begin{samepage}
\prayer{L.}{Lord our God, we are grateful for your gifts and for the opportunity to participate in the Lenten Fast through our denial in food and drink. May this be a help for us to imitate more closely your Son, Jesus Christ, who lives and reigns for ever and ever.}
\prayer{All}{Amen}
\end{samepage}
\medbreak

\paragraph*{Resurrection Sunday}

\begin{samepage}
\prayer{L.}{Lord our God, Giver of all goodness, fill our hearts with the Paschal Joy and similarly as you have given us food from the earth, grant that this new life, which Jesus Christ has gained for us through his death and resurrection and in whose mercy it has been given to us, will remain in us always. Who lives and reigns for ever and ever.}
\prayer{All}{Amen}
\end{samepage}
\medbreak

\paragraph*{Easter Season}

\begin{samepage}
\prayer{L.}{Almighty and eternal God, you have called us to share in the glory of your Risen Son. Accept our thanksgiving and grant that through the fulfilment of your will we may become participants in the heavenly banquet and to take the places prepared for us from the foundation of the world.}
\prayer{All}{Amen}
\end{samepage}
\medbreak

\paragraph*{Pentecost Sunday}

\begin{samepage}
\prayer{L.}{Lord Jesus Christ, we thank you for your gifts which we have enjoyed. We ask you to send upon the ways of our daily life the Spirit of love, the Spirit of strength and fidelity that as new people we may walk victoriously through life and enjoy the gifts prepared for us in the heavenly Kingdom. For you live and reign for ever and ever.}
\prayer{All}{Amen}
\end{samepage}
\medbreak
\newpage
\includeimage{pauline-praying.pdf}\pagestyle{empty}
\hspace{0pt}
\vfill

\textit{``Through the Word of God, the Eucharistic Sacrifice, the
Sacraments and the Liturgy of Hours, Christ continues his priestly
ministry effecting our salvation and glorifying the Father. This is
the very goal of our own ministry'' (art 20).}
\medbreak
O God our Father, you have sent your Son, the Word of
Truth, and the Spirit, the Sanctifier, to reveal to us the Mystery of
Divine life. Grant that we, who confess our true faith, may also
acknowledge the eternal glory of the Most Holy Trinity and the
unity of the Divine Persons in their power of creation, through
Christ our Lord. Amen.
\vfill
\textit{A Pauline Monk during prayer, an illumination from the missal of Jan Olbracht, from the library of Jasna Góra, from the end of the XV century.}
\hspace{0pt}
\newpage
\newpageheader{Life of Prayer and Sacrifice with Christ}{Prayers of Adoration}

\chapter{Life of Prayer and Sacrifice with Christ}

\instruction{The Second Vatican Council teaches that ``the Eucharistic Sacrifice is the source and summit of the Christian life'' (LG 11) and ``\ldots ~the other sacraments and indeed all ecclesiastical ministries and works of the apostolate are bound up with the Eucharist and are directed towards it'' (\textit{Presbyterorum Ordinis} 5). But for us, however, the Eucharist ``becomes the power to live a life according to the Evangelical Counsels and is an effective help in bringing to reality the mission of our Monastic Community within the Church of God'' (art. 20) and ``to live the Mystery of the Cross'' (art. 34).}
\medbreak

\section{Prayers of Adoration}

\subsection{Litany of the Most Blessed Sacrament}

\begin{litany}Lord, have mercy, \responce{Lord, have mercy.}\par 
\end{litany}
\begin{litany}Christ, have mercy, \responce{Christ, have mercy.}\par 
\end{litany}
\begin{litany}Lord, have mercy, \responce{Lord, have mercy.}\par 
\end{litany}
\begin{litany}Christ, hear us, \responce{Christ, graciously hear us.}\par 
\end{litany}
\begin{litany}God the Father of Heaven, \responce{Have mercy on us.}\par 
God the Son, Redeemer of the World, \par
God the Holy Spirit, \par
Holy Trinity, One God, \par
God, hidden in the Blessed Sacrament, \par
Eternal Word of the Father who became Flesh and dwelt among us, \par
Jesus, present in the Blessed Sacrament, true God and true Man, \par
Unceasing Sacrifice of the New Covenant, \par
Living Bread that comes down from Heaven, \par
Bread changed into the Flesh, through the omnipresent Word, \par
Wheat of the elect, \par
Banquet of Angels, \par
Memorial of the passion of Jesus Christ, \par
Testament of God's love for us, \par
Mystery of faith, \par
Miracle beyond comprehension, \par
Covenant of love and peace, \par
Source of all graces, \par
Pledge of the glorious resurrection, \par
Encouragement on the way to eternity, \par
Refuge of sinners, \par
Support of the weak and fallen, \par
Physician of the sick, \par
Consolation of the afflicted, \par
Divine Prisoner of our altars, \par
\end{litany}
\begin{litany}Be merciful, \responce{Spare us, O Lord.}\par 
\end{litany}
\begin{litany}Be merciful, \responce{Graciously hear us, O Lord.}\par 
\end{litany}
\begin{litany}From an unworthy reception of your Body and Blood, \responce{Spare us, O Lord.}\par 
From all pride, \par
From all dangers of the soul, \par
From lack of love of neighbour, \par
From every occasion of sin, \par
\end{litany}
\begin{litany}Through your holy Incarnation, \responce{Deliver us, O Lord.}\par 
Through your bitter passion and death, \par
Through your desire to eat this Passover with your disciples, \par
Through your infinite love by which you have instituted this Divine Sacrament, \par
Through your precious Blood, which you have left us on our altars, \par
Through your profound humility with which you have washed the feet of your apostles, \par
Through the five wounds of your most holy Body, \par
\end{litany}
\begin{litany}We sinners beseech you \responce{Hear us, O Lord.}\par 
That you may preserve and increase our faith, reverence, and devotion towards this admirable Sacrament, we beseech you, \par
That through this Food of Angels, you may cleanse us of all evils from within, we beseech you, \par
That you may preserve us in your divine grace, we beseech you, \par
That you may deliver us form all the snares of the enemy, we beseech you, \par
That you may strengthen and defend us with this Heavenly Food at the Hour of our death, we beseech you, \par
Son of God, source of compassion and mercy, we beseech you, \par
\end{litany}
\medbreak
\noindent\begin{minipage}{\textwidth}
\begin{litany}Lamb of God, you take away the sins of the world, \responce{Spare us, O Lord.}\par 
\end{litany}
\begin{litany}Lamb of God, you take away the sins of the world, \responce{Graciously hear us, O Lord.}\par 
\end{litany}
\begin{litany}Lamb of God, you take away the sins of the world, \responce{Have mercy on us.}\par 
\end{litany}
\end{minipage}
\medbreak
\begin{samepage}
\prayer{L.}{\textbf{O Lord Jesus, truly present on our altars under the appearance of Bread and Wine.}}
\prayer{All}{Fortify us with your Body and Blood.}
\end{samepage}
\medbreak
\begin{samepage}
\prayer{L.}{\textbf{Let us pray:}\\ O God, who in this wonderful Sacrament has left us a memorial of your passion, grant us the grace, we beseech you, so to venerate the sacred mysteries of your Body and Blood, that we may ever continue to feel within ourselves the blessed fruit of your Redemption, who live and reign, one God for ever and ever.}
\prayer{All}{Amen.}
\end{samepage}
\medbreak
 \newpage 

\subsection{Prayer of Adoration of the Blessed Sacrament}

O Lord, Jesus Christ, in this adoration I wish to respond to your invitation: ``Watch and pray.'' Therefore, I desire to honour you, to confess my faith and love for you, truly present in this Most Holy Sacrament. You promised to remain with us to the end of time. I thank you for this promise and ask for forgiveness for my lack of love and weak faith. Strengthen and purify me by your Holy Spirit. Illumine the eyes of my mind, that I may honour you, not only with my lips, but with the holiness of my life.
\medbreak

Jesus Christ, Son of the living God, you who are true God and true Man, I believe in your unique presence in the Most Holy Sacrament of the Eucharist and I thank you for so many blessings which you have showered upon me. I adore you for your Mother, whom you have given me as my Mother and Protectress. Following her example, I desire to live in union with you, my supreme God.
\medbreak

I salute you, O most merciful heart of my Lord, Jesus! I desire to make reparation for all the outrages committed by your enemies against your divine presence in this Sacrament. Through this act of humility and adoration, I wish to adore you in all places on earth, where you are neglected and less reverenced.
\medbreak

O Jesus, I love you with all my heart. I am sorry that in the past I have saddened you, my infinite Goodness. I resolve with the help of your grace not to offend you anymore and now, as poor as I am, I consecrate myself totally to you. I renounce my will, my senses, my desires and everything I possess. From this moment on, direct me according to your liking. I long only for your divine love, and I pray for the grace of perseverance in this love to the end and to fulfil perfectly your Holy Will. I commend to you the poor souls in Purgatory, particularly those who had a special devotion to the Most Holy Sacrament and to the Most Blessed Virgin Mary. I commend to you, also, the unfortunate sinners.
\medbreak

Finally, my Beloved Saviour, I join all my intentions with those of your Sacred Heart and thus joined, I offer them to your Eternal Father and I ask Him in your Holy Name to accept and graciously answer them. Amen.
\medbreak
\newpage
\newpageheader{Life of Prayer and Sacrifice with Christ}{Divine Praises}

\section{Divine Praises}

\begin{litany}My Lord and my God! \par
O Jesus in the Most Holy Sacrament, have mercy on us! \par
May the Most Holy Sacrament be praised at every moment. \par
Hail, True Body, born of the Virgin Mary. \par
May the Most Holy Sacrament be honoured now and forever. \par
Amen. \par
\end{litany}
\medbreak
\begin{litany}Blessed be God! \par
Blessed be his Holy Name! \par
Blessed be Jesus Christ, true God and true Man. \par
Blessed be the name of Jesus. \par
Blessed be his Most Sacred Heart. \par
Blessed be his Most Precious Blood. \par
Blessed be Jesus in the Most Holy Sacrament of the Altar. \par
Blessed be the Holy Spirit, the Paraclete. \par
Blessed be the great Mother of God, Mary most Holy. \par
Blessed be her Holy and Immaculate Conception. \par
Blessed be her glorious Assumption. \par
Blessed be the Name of Mary, Virgin and Mother. \par
Blessed be St. Joseph, her most chaste spouse. \par
Blessed be God in his angels and in his Saints. \par
\end{litany}
\medbreak
\newpageheader{Life of Prayer and Sacrifice with Christ}{An Act of Spiritual Holy Communion}

\section{An Act of Spiritual Holy Communion}

O Lord, Jesus Christ, truly present in the Blessed Sacrament, I believe in you, I adore you and I love you with all my heart. I am sorry for all my sins, because I desire to receive you in Holy Communion and remain with you. Come, O Jesus to my heart, spiritually at least and dwell in it forever\ldots
\medbreak

May the Body of our Lord Jesus Christ guard me for life everlasting.
\medbreak

I welcome you, Jesus into my poor heart. I dedicate myself totally to you. I am yours in life and in death. Amen.
\medbreak
\newpage
\newpageheader{Life of Prayer and Sacrifice with Christ}{Preparation For Holy Mass}

\section{Preparation For Holy Mass}

\subsection{Prayer of St. Ambrose}

I draw near, loving Lord Jesus Christ, to the table of your most delightful banquet in fear and trembling, a sinner, presuming not upon my own merits, but trusting rather in your goodness and mercy.
\medbreak

I have a heart and body defiled by my many offenses, a mind and tongue over which I have kept no good watch.
\medbreak

Therefore, O loving God, O awesome Majesty, I turn in my misery, caught in snares, to you the fountain of mercy, hastening to you for healing, flying to you for protection; and while I do not look forward to having you as Judge, I long to have you as Savior.
\medbreak

To you, O Lord, I display my wounds, to you I uncover my shame. I am aware of my many and great sins, for which I fear, but I hope in your mercies, which are without number.
\medbreak

Look upon me, then, with eyes of mercy, Lord Jesus Christ, eternal King, God and Man, crucified for mankind.
\medbreak

Listen to me, as I place my hope in you, have pity on me, full of miseries and sins, you, who will never cease to let the fountain of compassion flow.
\medbreak

Hail, O Saving Victim, offered for me and for the whole human race on the wood of the Cross.
\medbreak

Hail, O noble and precious Blood, flowing from the wounds of Jesus Christ, my crucified Lord, and washing away the sins of all the world.
\medbreak

Remember, Lord, your creature, whom you redeemed by your Blood. I am repentant of my sins, I desire to put right what I have done.
\medbreak

Take from me, therefore, most merciful Father, all my iniquities and sins, so that, purified in mind and body, I may worthily taste the Holy of Holies.
\medbreak

And grant that this sacred foretaste of your Body and Blood which I, though unworthy, intend to receive, may be the remission of my sins, the perfect cleansing of my faults, the banishment of shameful thoughts, and the rebirth of right sentiments; and may it encourage a wholesome and effective performance of deeds pleasing to you and be a most firm defense of body and soul against the snares of my enemies. Amen.
\medbreak

\subsection{Prayer of St. Thomas Aquinas}

Almighty eternal God, behold, I come to the Sacrament of your Only Begotten Son, our Lord Jesus Christ, as one sick to the physician of life, as one unclean to the fountain of mercy, as one blind to the light of eternal brightness, as one poor and needy to the Lord of heaven and earth.
\medbreak

I ask, therefore, for the abundance of your immense generosity, that you may graciously cure my sickness, wash away my defilement, give light to my blindness, enrich my poverty, clothe my nakedness, so that I may receive the bread of Angels, the King of kings and Lord of lords, with such reverence and humility, such contrition and devotion, such purity and faith, such purpose and intention as are conducive to the salvation of my soul.
\medbreak

Grant, I pray, that I may receive not only the Sacrament of the Lord's Body and Blood, but also the reality and power of that Sacrament. O most gentle God, grant that I may so receive the Body of your Only Begotten Son our Lord Jesus Christ, which he took from the Virgin Mary, that I may be made worthy to be incorporated into his Mystical Body and to be counted among its members.
\medbreak

O most loving Father, grant that I may at last gaze for ever upon the unveiled face of your beloved Son, whom I, a wayfarer, propose to receive now veiled under these species: Who lives and reigns with you for ever and ever. Amen.
\medbreak
 \newpage 

\subsection{Prayer to Our Lady}

O most blessed Virgin Mary, Mother of tenderness and mercy, I, a miserable and unworthy sinner, fly to you with all the affection of my heart and I beseech your motherly love, that, as you stood by your most dear Son, while he hung on the Cross, so, in your kindness, you may be pleased to stand by me, a poor sinner, and all Priests who today are offering the Sacrifice here and throughout the entire holy Church, so that with your gracious help we may offer a worthy and acceptable oblation in the sight of the most high and undivided Trinity. Amen.
\medbreak

\subsection{The Priest's Declaration of Purpose before Mass}

My intention is to celebrate Mass and to consecrate the Body and Blood of our Lord Jesus Christ according to the Rite of Holy Roman Church, to the praise of almighty God and all the Church triumphant, for my good and that of all the Church militant, for all who have commended themselves to my prayers in general and in particular, and for the welfare of Holy Roman Church. Amen.
\medbreak

May the almighty and merciful Lord grant us joy with peace, amendment of life, room for true repentance, the grace and consolation of the Holy Spirit and perseverance in good works. Amen.
\medbreak

\subsection{Preface II of the Holy Eucharist}

It is truly right and just, our duty and our salvation, always and everywhere to give you thanks, Lord, holy Father, almighty and eternal God, through Christ our Lord.
\medbreak

For at the last Supper with his Apostles, establishing for the ages to come the saving memorial of the Cross, he offered himself to you as the unblemished Lamb, the acceptable gift of perfect praise.
\medbreak

Nourishing your faithful by this sacred mystery, you make them holy, so that the human race, bounded by one word, may be enlightened by one faith and united by one bond of charity.
\medbreak

\subsection{Prayer to St. Joseph}

God, who has conferred upon us a royal priesthood, we pray to you give us grace to minister at your holy altar with hearts as clean and lives as blameless as that of blessed Joseph, who was found worthy to hold in his arms and, with all reverence, to carry you the only-begotten Son, born of the Virgin Mary. Enable us this day to receive worthily the sacred Body and Blood of your Son, and enable us to win an everlasting reward in the world to come; through the same Christ our Lord. Amen.
\medbreak
\newpage
\newpageheader{Life of Prayer and Sacrifice with Christ}{Thanksgiving After Holy Mass}

\section{Thanksgiving After Holy Mass}

\subsection{Eucharistic Prayer drawn from the teaching of the Twelve Apostles (Didache)}

\begin{samepage}
\begin{litany}\prayer{}{We thank you, Father,\par for the holy vine of David, your servant,\par which you have made known to us\par through your servant, Jesus.\par \textbf{\textit{To you be glory for ever and ever!}}}\end{litany}
\end{samepage}
\medbreak
\begin{samepage}
\begin{litany}\prayer{}{We thank you, Father,\par for the life and knowledge\par you have made known to us\par through your servant, Jesus.\par \textbf{\textit{To you be glory for ever and ever!}}}\end{litany}
\end{samepage}
\medbreak
\begin{samepage}
\begin{litany}\prayer{}{In the same way that this bread which is now broken\par was scattered upon the mountains,\par was gathered up again, to become one,\par so may your Church be gathered together\par from the ends of the earth into your Kingdom,\par for yours is the glory and the power\par through Jesus Christ for ever.}\end{litany}
\end{samepage}
\medbreak
\begin{samepage}
\begin{litany}\prayer{}{We thank you, Holy Father, for your Holy Name,\par which you have made to dwell in our hearts,\par and for the knowledge, faith and immortality\par which you have made known to us\par through your servant, Jesus.\par \textbf{\textit{To you be glory for ever and ever!}}}\end{litany}
\end{samepage}
\medbreak
\begin{samepage}
\begin{litany}\prayer{}{May your grace come down and may this world pass away.\par Hosanna to God of David!\par Who is a saint, let him approach!\par Who is not, let him repent!\par Come, Lord Jesus! Amen.}\end{litany}
\end{samepage}
\medbreak

\subsection{Prayer of St. Thomas Aquinas}

I give you thanks, Lord, holy Father, almighty and eternal God, who have been pleased to nourish me, a sinner and your unworthy servant, with the precious Body and Blood of your Son, our Lord Jesus Christ: this through no merits of mine, but due solely to the graciousness of your mercy.
\medbreak

And I pray that this Holy Communion may not be for me an offense to be punished, but a saving plea for forgiveness. May it be for me the armor of faith, and the shield of good will. May it cancel my faults, destroy concupiscence and carnal passion, increase charity and patience, humility and obedience and all the virtues, may it be a firm defense against the snares of all my enemies, both visible and invisible, the complete calming of my impulses, both of the flesh and of the spirit, a firm adherence to you, the one true God, and the joyful completion of my life's course.
\medbreak

And I beseech you to lead me, a sinner, to that banquet beyond all telling, where with your Son and the Holy Spirit you are the true light of your Saints, fullness of satisfied desire, eternal gladness, consummate delight and perfect happiness. Through Christ our Lord. Amen.
\medbreak

\subsection{Petitions to the Most Holy Redeemer}

\begin{samepage}
\begin{litany}\prayer{}{Soul of Christ, sanctify me.\par Body of Christ, save me.\par Blood of Christ, embolden me.\par Water from the side of Christ, wash me.}\end{litany}
\end{samepage}
\medbreak
\begin{samepage}
\begin{litany}\prayer{}{Passion of Christ, strengthen me.\par O good Jesus, hear me.\par Within your wounds hide me.\par Never permit me to be parted from you.}\end{litany}
\end{samepage}
\medbreak
\begin{samepage}
\begin{litany}\prayer{}{From the evil Enemy defend me.\par At the hour of my death call me\par and bid me come to you,\par that with your Saints I may praise you\par for age upon age. Amen.}\end{litany}
\end{samepage}
\medbreak

\subsection{Prayer of self-offering to Jesus Christ}

Receive, Lord, my entire freedom. Accept the whole of my memory, my intellect and my will. Whatever I have or possess, it was you who gave it to me; I restore it to you in full, and I surrender it completely to the guidance of your will. Give me only love of you together with your grace, and I am rich enough and ask for nothing more. Amen.
\medbreak

\subsection[The Universal Prayer]{The Universal Prayer \\ {\small\normalfont (attributed to Pope Clement XI)}}

I believe, O Lord, but may I believe more firmly; I hope, but may I hope more securely; I love, but may I love more ardently; I sorrow, but may I sorrow more deeply. I adore you as my first beginning; I long for you as my last end; I praise you as my constant benefactor; I invoke you as my gracious protector.
\medbreak

By your wisdom direct me, by your righteousness restrain me, by your indulgence console me, by your power protect me.
\medbreak

I offer you, Lord, my thoughts to be directed to you, my words, to be about you, my deeds, to respect your will, my trials, to be endured for you.
\medbreak

I will whatever you will, I will it because you will it, I will it in the way you will it, I will it for as long as you will it.
\medbreak

Lord, enlighten my understanding, I pray: arouse my will, cleanse my heart, sanctify my soul.
\medbreak

May I weep for past sins, repel future temptations, correct evil inclinations, nurture appropriate virtues.
\medbreak

Give me, good God, love for you, hatred for myself, zeal for my neighbor, contempt for the world.
\medbreak

May I strive to obey superiors, to help those dependent on me, to have care for my friends, forgiveness for my enemies.
\medbreak

May I conquer sensuality by austerity, avarice by generosity, anger by gentleness, lukewarmness by fervor.
\medbreak

Render me prudent in planning, steadfast in dangers, patient in adversity, humble in prosperity.
\medbreak

Make me, O Lord, attentive at prayer, moderate at meals, diligent in work, steadfast in intent.
\medbreak

May I be careful to maintain interior innocence, outward modesty, exemplary behavior, a regular life.
\medbreak

May I be always watchful in subduing nature, in nourishing grace, in observing your law, in winning salvation.
\medbreak

May I learn from you how precarious are earthly things, how great divine things, how fleeting is time, how lasting things eternal.
\medbreak

Grant that I may prepare for death, fear judgment, flee hell, gain paradise. Through Christ our Lord. Amen.
\medbreak

\subsection{Prayer to our Lady}

O Mary, Virgin and Mother most holy, behold, I have received your most dear Son, whom you conceived in your immaculate womb, brought forth, nursed and embraced most tenderly.
\medbreak

Behold him at whose sight you used to rejoice and be filled with all delight; him whom, humbly and lovingly, once again I present and offer him to you to be clasped in your arms, to be loved by your heart, and to be offered up to the Most Holy Trinity as the supreme worship of adoration, for your own honor and glory and for my needs and for those of the whole world.
\medbreak

I ask you therefore, most loving Mother: entreat for me the forgiveness of all my sins and, in abundant measure, the grace of serving him in the future more faithfully, and at the last, final grace, so that with you I may praise him for all the ages of ages. Amen.
\medbreak
 \newpage 

\subsection{An Act of Worship}

I adore you Jesus, truly present in my heart, born of the Virgin Mary, the only-begotten Son and the glory of the Father. I thank you my Redeemer that you deigned to come to me, a miserable and poor sinner.
\medbreak

Together with Mary, I offer you to your Father, and with you and through you, may eternal glory, thanksgiving and prayer for world peace go up before his heavenly throne. Illumine my mind and heart that I may renew within myself the desire to belong to you alone. For you live for ever and ever. Amen.
\medbreak
\newpage
\newpageheader{Life of Prayer and Sacrifice with Christ}{Prayers for Special Occasions}

\section{Prayers for Special Occasions}

\subsection[Prayer before the homily]{Prayer before the homily \\ {\small\normalfont (St. Anselm)}}

My God, grant me eloquence, the fullness of gentleness and wisdom that would not arouse me with an arrogance and superiority over others. I beg you to implant in my mouth the words of consolation, strength and encouragement. May I, through the power of the Holy Spirit inspire the good to be better, convert the wayward and, through my words and example, direct them on the path of righteousness. May the words you give to your servant be as hot as missiles and like fiery arrows that would penetrate and enkindle the souls of the listeners to a holy fear and love of you. Amen.
\medbreak

\subsection{Prayer after the homily}

God the Holy Spirit, ignite our hearts and souls with the fire of your love that we may all serve you in submission and fidelity to your Word. May the efforts of my priestly service bring to me and to all my listeners and abundance of fruits. May their faith and hope be reflected in their daily living. For you live and reign as the Love of the Father and of the Son, one God for endless ages. Amen.
\medbreak

\subsection{Prayer of the Confessor before administering the Sacrament of Reconciliation}

God, Father of mercy, giver of all gifts and charisms, confer upon me the gift of your wisdom, responsibility and unique discretion to know the spiritual needs of the penitents that I may, as their confessor, help them in their process of conversion of life. Grant me the grace to discern accurately the true state of their consciences, that I may form them in truth and in the spirit of the Gospel.
\medbreak

God, Father of compassion, give me the grace to be always a helpful servant to penitents, in their difficult process of conversion; to lift them up from an even worse downfall. Give me the grace to sincerely sympathise with them over their offence to God and ignorance of His love. As a minister of your mysteries, may I communicate your merciful love that leads man, overburdened with guilt, to reconciliation with you and with the Church, which, through the sins of the faithful, sustains many wounds and injustices.
\medbreak

Lord, Jesus Christ you came into the world to seek out and save those who were lost, to call to repentance the sinner, to heal the broken-hearted. You desire not the death of the sinner, but await his conversion. Give me the grace to know for certain, that your love is stronger and greater than all our sins. May my intention to administer this sacrament be always pure and clear, that God may be honoured in everything and penitents be sanctified, spiritually liberated and healed. May my priestly zeal be prudent, my love pure and selfless, because I am administering this sacrament in Your Name, on your behalf, when I become somehow one with you, Jesus Christ, the Eternal Priest. Lord Jesus Christ, grant me the grace to be fully aware of your presence and your saving work. Make me merciful, kind and wise in asking questions, in carrying out pastoral dialogues, cautious and foreseeing in giving spiritual advice, in imposing the sacramental penance, in encouraging and strengthening the realisation of ones vocation in life.
\medbreak

O Lord, you have sent the Holy Spirit on your Apostles to forgive sins and through the breath of this same Holy Spirit you have bestowed the power to loose and to bind. May this Holy Spirit, through his power, make my ministry successful and fruitful. Be with me, O Holy Spirit, especially in difficult situations and cases, that in showing understanding and compassion to penitents, I may never depart from the sound teaching of the Church.
\medbreak

Strengthen me O Lord, one God in the unity of the Most Holy Trinity, with your spiritual strength, that in helping others, I myself may not be rejected.\\ Mother of Mercy, pray for us!
\medbreak

\subsection{Prayer of the Confessor after administering the Sacrament of Reconciliation}

O Lord Jesus Christ, our Redeemer and Conqueror of sin, I thank you that you have become the sacrifice of forgiveness for our sins and the sins of the whole world. In your wounds, lay our healing! You did not spare your Precious Blood to deliver us from the misery of our sins. Give me, and all my penitents, the grace to participate fully in your victory over sin that we may die to sin that thus it may never rule over us; and that we may live in your grace and grow in newness of life. I entrust to your infinite mercy, all my penitents that they may keep their good resolutions, on the road of new life. Surround with your merciful love especially those who for some reason, did not receive the sacrament of absolution, that they may not carry grudges against anybody, but come to you in due time, according to your Holy Will.
\medbreak

O Lord of clemency, through your all powerful grace, make up for all my deficiencies, limitations, human weaknesses, errors, unawareness and shortcomings, committed while administering the sacrament of mercy. Be merciful to me!
\medbreak

O Lord Jesus Christ, always rich in your mercy, you change the unjust into just, to the depressed you offer the hope and joy of a new life. Grant us the grace, we pray, that we may never waste the gifts received from you. As we, experience your merciful love, may we, too, be witnesses of this love toward others. \\ Mary, Refuge of sinners, pray for us!
\medbreak

\subsection{Prayer before study}

O Lord, our Almighty God, you are the source of light and wisdom, pour into my will, my memory and all my faculties the grace of your blessing. Give me sharpness of mind to understand, a ability to remember, an ability to learn, a skill to interpret and an abundant grace of eloquence. Shed your light at the beginning of my study, direct its progress and bring it to a fulfilling outcome. For you live and reign for ever and ever. Amen.
\medbreak

\subsection{Prayer after study}

I thank you, Lord God for light and help in the efforts of my memory and mind. Give me the grace to use this knowledge for your glory, for my service to the Church of Christ, for the benefit of my Order and for the instruction of souls. Let the flame of my knowledge contribute to knowing you better and loving you more. For you alone are one God in the unity of the Most Holy Trinity. Amen.
\medbreak

\subsection[Prayer for the Sanctification of Labour]{Prayer for the Sanctification of Labour \\ {\small\normalfont St. Pius X}}

Glorious St. Joseph, model of all who are devoted to labor, obtain for me the grace to work in the spirit of penance in expiation of my many sins; to work conscientiously by placing love of duty above my inclinations; to gratefully and joyously deem it an honor to employ and to develop by labor the gifts I have received from God, to work methodically, peacefully, and in moderation and patience, without ever shrinking from it through weariness or difficulty to work; above all, with purity of intention and unselfishness, having unceasingly before my eyes death and the account I have to render of time lost, talents unused, good not done, and vain complacency in success, so baneful to the work of God. All for Jesus, all for Mary, all to imitate thee, O patriarch St. Joseph! This shall be my motto for life and eternity. Amen.
\medbreak
\newpage
\includeimage{crucifix.pdf}\pagestyle{empty}
\hspace{0pt}
\vfill

\textit{``Our fathers showed us how to live the Mystery of the
Cross as expressed in a life of willingly embraced difficulties and
mortifications. Thus they became more like Christ in suffering for
the sins of the world.'' (art 34).}
\medbreak
Almighty God, your Son, obedient to your Holy Will suffered the agony of the Cross that all people may be saved. Grant
us, we pray that knowing the mystery of the Redemption of the
world, we may obtain its blessed fruits in heaven. Through Christ
our Lord. Amen.
\vfill
\textit{Crucifix found by Queen Elizabeth Rakuszanka, in the Krakow
school. The chapel of Our Lady of Jasna Gora, 2nd part of the
XVth century}
\hspace{0pt}
\newpage
\newpageheader{Life of Prayer and Sacrifice with Christ}{Prayers to Christ Crucified}

\section{Prayers to Christ Crucified}

\subsection{Litany of the Holy Cross}

\begin{litany}Lord, have mercy, \responce{Lord, have mercy.}\par 
\end{litany}
\begin{litany}Christ, have mercy, \responce{Christ, have mercy.}\par 
\end{litany}
\begin{litany}Lord, have mercy, \responce{Lord, have mercy.}\par 
\end{litany}
\begin{litany}Christ, hear us, \responce{Christ, graciously hear us.}\par 
\end{litany}
\begin{litany}God the Father of Heaven, \responce{Have mercy on us.}\par 
God the Son, Redeemer of the World, \par
God the Holy Spirit, \par
Holy Trinity, One God, \par
\end{litany}
\medbreak
\begin{litany}Holy Cross, tree of life, \responce{Save us.}\par 
Holy Cross, sign of all good, \par
Holy Cross, ornament of Christianity, \par
Holy Cross, instrument of Redemption, \par
Holy Cross, source of all blessings, \par
Holy Cross, consolation of the afflicted, \par
Holy Cross, terror of demons, \par
Holy Cross, cathedral of patience, \par
Holy Cross, defence of Christians, \par
Holy Cross, key to the heavenly Kingdom, \par
Holy Cross, refuge of the oppressed and tempted, \par
Holy Cross, hope of sinners, \par
Holy Cross, light of the just, \par
Holy Cross, rest of the perfect, \par
Holy Cross, consolation of the suffering Church, \par
Holy Cross, glory of the triumphant Church, \par
Holy Cross, sign predicted by Patriarchs and Prophets, \par
Holy Cross, joy of the Apostles, \par
Holy Cross, crown of Martyrs, \par
Holy Cross, joy of Confessors, \par
Holy Cross, shield of Virgins, \par
Holy Cross, pledge of eternal happiness, \par
Holy Cross, sprinkled with the Blood of Jesus Christ, \par
Holy Cross, which prepared eternal happiness for us, \par
Holy Cross, on which death was conquered, \par
\end{litany}
\medbreak
\begin{litany}Jesus Christ, crucified Saviour, \responce{Have mercy on us.}\par 
Jesus Christ, who suffered the shame of the Cross, \par
Jesus Christ, obedient to your heavenly Father to death on the Cross, \par
Jesus Christ, who with your Body on the Cross, has paid for our sins, \par
Jesus Christ, who have nailed to the Cross the sentence of death for us, \par
Jesus Christ, who through your Cross, have defeated Satan, \par
Jesus Christ, who through your Cross have restored to us the freedom of the sons of God, \par
Jesus Christ, who with the sign of the Cross will come to judge the living and the dead, \par
\end{litany}
\medbreak
\begin{litany}Be merciful, \responce{Spare us, O Jesus.}\par 
\end{litany}
\begin{litany}Be merciful, \responce{Hear us, O Jesus.}\par 
\end{litany}
\begin{litany}From all sin, \responce{Deliver us, O Jesus.}\par 
From the snares of our eternal enemy, \par
From eternal damnation, \par
From the fear of death, \par
Through your Holy Cross, \par
Through your five wounds, \par
Through your precious Blood, \par
Through your death, \par
\end{litany}
\medbreak
\begin{litany}We sinners implore you, \responce{Hear us, O Jesus.}\par 
That we may be true lovers of your Cross, \par
That we may glorify ourselves only in the Holy Cross, \par
That you may, through the power of your Cross deliver us from all dangers of body and soul, \par
\end{litany}
\medbreak
\noindent\begin{minipage}{\textwidth}
\begin{litany}Lamb of God, you take away the sins of the world, \responce{Spare us, O Lord.}\par 
\end{litany}
\begin{litany}Lamb of God, you take away the sins of the world, \responce{Graciously hear us, O Lord.}\par 
\end{litany}
\begin{litany}Lamb of God, you take away the sins of the world, \responce{Have mercy on us.}\par 
\end{litany}
\end{minipage}
\medbreak
\begin{samepage}
\prayer{L.}{\textbf{We adore you O Christ, and we praise you.}}
\prayer{All}{Because by your Holy Cross, you have redeemed the world.}
\end{samepage}
\medbreak
\begin{samepage}
\prayer{L.}{\textbf{Let us pray:}\\ Almighty God, your Son, obedient to your Holy Will, suffered the agony of the Cross, that all people may be brought to salvation. Grant that we, who know the mystery of Redemption on earth, may receive its fruits in heaven. Through Christ, our Lord.}
\prayer{All}{Amen.}
\end{samepage}
\medbreak

\subsection{Stations of the Cross}

\subsubsection*{The First Station \\ {\small\normalfont Jesus is condemned to death}}

\begin{samepage}
\prayer{L.}{\textbf{We adore you, O Christ, and we praise you!}}
\prayer{All}{Because by your Holy Cross, you have redeemed the world.}
\end{samepage}
\medbreak

\textit{``\ldots ~but the Jews shouted still louder to crucify him\ldots ~so in the end Pilate handed him over to them to be crucified \ldots ''}
\medbreak

Every road of religious life begins with the first station of the cross. The decision to abandon the world, and to petition for admission at the Monastery door has passed. Was I fully aware that, like Christ, I was entering the road of the cross; truly a pilgrimage of the cross? Was I mindful of the truth that the sweetness of mystic experience will dawn upon me, only after the night of affliction?
\medbreak

O most holy Will of God, may I seek only you and may the fulfilment of your Will be my greatest delight. Lord Jesus, may you be for my soul as a true vine to its branches.
\medbreak

\subsubsection*{The Second Station \\ {\small\normalfont Jesus accepts his cross}}

\begin{samepage}
\prayer{L.}{\textbf{We adore you, O Christ, and we praise you!}}
\prayer{All}{Because by your Holy Cross, you have redeemed the world.}
\end{samepage}
\medbreak

Jesus began the road of his suffering in the Garden of Gethsemane. Later, the night interrogations, the crowning with thorns, the scourging, exhausted him, even before his passion.
\medbreak

On the day of my admission I felt very energetic and enthusiastic, ready to begin a new adventure, a new state of life. Because of my weakness, God, in His wisdom, did not allow me to foresee the heavy weight of the religious and priestly life. He did not want me to be disheartened before I had learned to depend on His grace alone.
\medbreak

O Jesus, teach me to hold your image with the cross in my hands and in my heart, that I may have the wisdom to understand these words: ``Once having made the choice, I have to choose anew constantly.'' Teach me to remember the glorious message of Easter.
\medbreak

\subsubsection*{The Third Station \\ {\small\normalfont Jesus falls the first time}}

\begin{samepage}
\prayer{L.}{\textbf{We adore you, O Christ, and we praise you!}}
\prayer{All}{Because by your Holy Cross, you have redeemed the world.}
\end{samepage}
\medbreak

Jesus, because of the weight of the cross, and the noise of the pushing crowd, you begin to feel weak, unstable and shaky. You begin to stumble on the stony road and all of a sudden you fall to the ground.
\medbreak

This can happen to me too, especially at the beginning of my religious formation. I discover in myself more and more areas that are totally alienated from the ideals of Holy Scripture, of my religious rules and of the ascetical writings which I am reading. These ideals contrast with my real human weakness. Then, I ask myself: ``Is this the right place for me?''
\medbreak

O Lord, be my strength when I become discouraged. I cannot depart from you. I will arise again and go on. Once again, I will look at my Brothers and discover in them a hidden greatness.
\medbreak
 \newpage 

\subsubsection*{The Fourth Station \\ {\small\normalfont Jesus meets his sorrowful Mother}}

\begin{samepage}
\prayer{L.}{\textbf{We adore you, O Christ, and we praise you!}}
\prayer{All}{Because by your Holy Cross, you have redeemed the world.}
\end{samepage}
\medbreak

For a short time, the eyes of Jesus and Mary, his Mother, met in a sorrow beyond description. It was just a brief moment and yet it contained everything: their love, their pain, their determination and their obedience to the Will of the Eternal Father.
\medbreak

It is a common conviction, that the decision to choose the consecrated religious life, is born in the heart of the Mother of God. Every good inspiration comes to us from heaven through the mediation of Mary, as an answer to her request: ``they have no wine'' (Jn 2:4). This is very true in every vocation to the Pauline Order.
\medbreak

O Queen of Czestochowa! I am with you, I remember, and am vigilant to cherish the love of Jesus for me, the love that you have given me. May this love and devotion to your sweet name reign in my heart and engage all my emotions.
\medbreak

\subsubsection*{The Fifth Station \\ {\small\normalfont Simon of Cyrene helps Jesus to carry his cross}}

\begin{samepage}
\prayer{L.}{\textbf{We adore you, O Christ, and we praise you!}}
\prayer{All}{Because by your Holy Cross, you have redeemed the world.}
\end{samepage}
\medbreak

Simon, probably forced to carry the cross, did not realise whom he was helping, the important thing was he actually helped. A kind-hearted person would respond in the same way.
\medbreak

In the time of distress, God does not leave us alone, but sends us a friend with encouragement. In the case of religious community life, it may be a spiritual director, a confessor, a superior or a brother. Through them, our Heavenly Father watches over us and supports us in the time of our inner needs.
\medbreak

May we always have the courage and strength to ask for, or accept, the help offered to us by other people, but always arranged by Divine Providence!
\medbreak

\subsubsection*{The Sixth Station \\ {\small\normalfont Veronica wipes the Face of Jesus}}

\begin{samepage}
\prayer{L.}{\textbf{We adore you, O Christ, and we praise you!}}
\prayer{All}{Because by your Holy Cross, you have redeemed the world.}
\end{samepage}
\medbreak

Out of great compassion, a courageous woman, by the name of Veronica wipes the bloody, sweaty face of her beloved Master on his way to Calvary. O St. Veronica you have done what I would do when I look on the afflicted face of my beloved Redeemer.
\medbreak

Often, we see in our Monasteries persons, weighed down with keen, personal suffering of sickness or age. There is a need for great souls, filled with goodness, who would hasten with help and bring them solace and comfort. Sometimes, just a little kindness, a smile, a good word and above all a prayer, changes unexpectedly sadness into joy.
\medbreak

O Jesus Christ, teach me to dry the tears, not only of my brothers in religious life, but also those that stream down the faces of my brothers and sisters in the confessional and at the Monastery door.
\medbreak

\subsubsection*{The Seventh Station \\ {\small\normalfont Jesus falls the second time}}

\begin{samepage}
\prayer{L.}{\textbf{We adore you, O Christ, and we praise you!}}
\prayer{All}{Because by your Holy Cross, you have redeemed the world.}
\end{samepage}
\medbreak

What caused Jesus to fall again? Certainly not the cross, Simon of Cyrene carried it. Yes, Jesus noticed an evil desire in my soul beginning to take root. That is the real reason for the Saviour's fall.
\medbreak

Jesus cautions me with the words, once addressed to St. Peter: ``\ldots ~Simon, Simon, the devil has demanded your soul \ldots '' I ask and then answer with St. Paul the Apostle: ``I am a miserable man! Who will deliver me from death of the flesh? The grace of God, through Jesus Christ \ldots ''
\medbreak

O Lord, defend me from Satan's trap and from the frailty of my will.
\medbreak

\subsubsection*{The Eighth Station \\ {\small\normalfont Jesus consoles the women of Jerusalem}}

\begin{samepage}
\prayer{L.}{\textbf{We adore you, O Christ, and we praise you!}}
\prayer{All}{Because by your Holy Cross, you have redeemed the world.}
\end{samepage}
\medbreak

``Daughters of Jerusalem do not weep for me, but for your children.'' Often, the world feels sorry for those who have left it, in order to serve God and his people exclusively. In its eyes, they seem to be foolish and naïve for not developing and using wisely their talents and for embracing all types of inconveniences \ldots ~but it is a false, pagan sympathy. Behold our Lord and God who teaches us in the light of eternity by saying: Come, you blessed of my Father, inherit the kingdom, prepared for you from the very foundation of the world.''
\medbreak

O Jesus, teach me to guard and increase the grace of my monastic, hermitical and apostolic vocation to live for God in solitude: ``Alone with God, alone.''
\medbreak

\subsubsection*{The Ninth Station \\ {\small\normalfont Jesus falls the third time}}

\begin{samepage}
\prayer{L.}{\textbf{We adore you, O Christ, and we praise you!}}
\prayer{All}{Because by your Holy Cross, you have redeemed the world.}
\end{samepage}
\medbreak

A great loss of physical strength caused Jesus to collapse again, shortly before he arrived at the hill of Calvary. Before his public mission, Jesus was tempted by the devil upon another hill. O Lord, it was not the devil who struck you down. It was me. Through my arrogance and pride, I pushed and knocked you over, I, a man particularly consecrated to you. The pleasure of passion has begotten sin and sin brought forth death.
\medbreak

O Holy Spirit, give me the gift of prudence and courage that I may always resist the lustful thoughts and actions that tarnish purity of heart. O Immaculate Virgin Mary, my Mother, cover me with your mantle of purity and innocence and help me to keep my vow of chastity untainted.
\medbreak

\subsubsection*{The Tenth Station \\ {\small\normalfont Jesus is stripped of his garments}}

\begin{samepage}
\prayer{L.}{\textbf{We adore you, O Christ, and we praise you!}}
\prayer{All}{Because by your Holy Cross, you have redeemed the world.}
\end{samepage}
\medbreak

The world in its immoderate appetite for sensuality, goes so far as it dares to disrobe even the God-Man himself of his divine and human dignity. May the darkness of which the Gospels speak, envelop the hill of Golgotha and cover the humiliated Lord, exposed to the jeer of the crowd. Every person, even the consecrated, may strip himself of sanctifying grace, by breaking the virtues and the vows of chastity, poverty and obedience. Through a grievous infidelity to the religious consecration, a Father or a Brother exposes himself to a most humiliating darkness.
\medbreak

O Jesus, my Lord! May I never depart from your Holy Will, because of the lack of vigilance on my part. May the sign of my habit, constantly remind me of my consecration and may it strengthen the faith of the people I serve.
\medbreak

\subsubsection*{The Eleventh Station \\ {\small\normalfont Jesus is nailed to the Cross}}

\begin{samepage}
\prayer{L.}{\textbf{We adore you, O Christ, and we praise you!}}
\prayer{All}{Because by your Holy Cross, you have redeemed the world.}
\end{samepage}
\medbreak

``When they arrived at Golgotha, they crucified him there.'' Jesus himself stretched out his arms and feet for the nailing. This was his decision to obey his Father in everything.
\medbreak

O Jesus, my Lord! For you I have left my family and friends; for you, I have disowned my property; for you, I have given up the right to exercise my personal freedom; for you, I have decided to follow the evangelical counsels, the constitutions and directory of my Pauline Order. This is my cross! This is my inadequate imitation of you!
\medbreak

``O Lord Jesus! to you, I want to offer everything, and for you, I desire to live and die.'' Defend me, O Lord Jesus from the temptation to come down from the cross of my religious vows, pronounced in the Order of St. Paul the First Hermit.
\medbreak

\subsubsection*{The Twelfth Station \\ {\small\normalfont Jesus dies on the Cross}}

\begin{samepage}
\prayer{L.}{\textbf{We adore you, O Christ, and we praise you!}}
\prayer{All}{Because by your Holy Cross, you have redeemed the world.}
\end{samepage}
\medbreak

When we were still in the slavery of sin, Jesus delivered himself for our sake to death, a death on the cross. St John, the beloved Apostle, testified to this love when he wrote: ``Those who were his own, he loved them to the end.''
\medbreak

The Crucified Lord awaits my self-denial in the monastic community life. He expects my share in his bitter passion and crucifixion. The Lord allowed his Sacred Heart to be pierced to show me the height and depth, the length and breadth of his love for me. His Heart was opened wide to be my refuge and my hiding place, whenever I feel ashamed because of my shortcomings and even serious failures.
\medbreak

O Lord Jesus, I confess my weakness. I know that without you, I can do nothing. Inflame my soul with the fire of your love that I may love you above all things and my brothers for your sake, remembering the words of St. Paul, the Apostle: ``\ldots ~but without love, nothing can help me\ldots '' (2 Cor 13:1-13).
\medbreak

\subsubsection*{The Thirteenth Station \\ {\small\normalfont Jesus is taken down from the Cross}}

\begin{samepage}
\prayer{L.}{\textbf{We adore you, O Christ, and we praise you!}}
\prayer{All}{Because by your Holy Cross, you have redeemed the world.}
\end{samepage}
\medbreak

Mary, the Mother of all Sorrows, holds the dead body of her Son. With a broken heart, she cries out: ``All of you, who pass by, see if there is any sorrow like my sorrow?''
\medbreak

I earnestly desire, O Mother of all Sorrows, to place the moment of my death into your hands. At the foot of the cross, your dying Son has entrusted me to you, and you to me. Now, I have a right to ask you to be with me at the hour of my death.
\medbreak

O my sweet Jesus, I will be eternally grateful to you in heaven, that when on earth, I have followed your way, the simplest and the shortest, through Mary, our Mother. Give me the grace to obey always the request of your Mother: ``Do whatever he tells you'' (Jn 2:5).
\medbreak

\subsubsection*{The Fourteenth Station \\ {\small\normalfont Jesus is buried in the tomb}}

\begin{samepage}
\prayer{L.}{\textbf{We adore you, O Christ, and we praise you!}}
\prayer{All}{Because by your Holy Cross, you have redeemed the world.}
\end{samepage}
\medbreak

The body of our beloved Redeemer, wrapped in the shroud, lies in peace in the niche hewn in the rock.
\medbreak

There will the last station of my earthly pilgrimage be. Clothed in my white habit, with the Rosary at my side, and with the cross in my hands, my mortal body will be laid out in a humble, wooden coffin, but my Brothers will pray for my immortal soul: ``Eternal rest grant unto him, O Lord, and let perpetual light shine upon him. May he rest in peace!'' My mortal body will be laid to rest among my Brothers, awaiting the glorious resurrection.
\medbreak

O, my beloved Saviour, even now I thank you for the grace of religious life, for the fraternal charity, prayers and care of my Brothers. Be merciful to me on the day of my judgement!
\medbreak

\subsection{I will remain under the Cross}

I will remain under your Cross, my Lord and Redeemer, when people praise me. I will recall your humiliation, that my vanity may be humiliated.
\medbreak

I will remain under your Cross, my Lord and Redeemer, when people reproach me. I will recall your silence, that my lips may be sealed.
\medbreak

I will remain under your Cross, my Lord and Redeemer, when people will persecute me. I will remember the scourging of your sacred Body, that my self-love may be abased.
\medbreak

I will remain under your Cross, my Lord and Redeemer, when I will suffer shortage of necessities, and I will remember your poverty. Teach me to rejoice in my material and spiritual deficiencies.
\medbreak

I will remain under your Cross, my Lord and Redeemer, when evil thoughts disturb my mind. I will remember your crowning with thorns, that my thoughts may be calm and collected.
\medbreak

I will remain under your Cross, my Lord and Redeemer, when evil passions storm my heart. I will be mindful of your bruised body, that I may overcome the temptation of the flesh.
\medbreak

I will remain under your Cross, my Lord and Redeemer, when people judge me unjustly. I will pray for them: ``Our Father, forgive them.'' Teach me to forgive always.
\medbreak

I will remain under your Cross, my Lord and Redeemer, when my conscience accuses me and leads me to desperation. I will recall your words spoken to the thief: ``Today you will be with me in paradise.''
\medbreak

I will remain under your Cross, my Lord and Redeemer, when the world and people reject me. I will remember the words, addressed to your Mother: ``Woman behold your son.'' She will console me.
\medbreak

I will remain under your Cross, my Lord and Redeemer, when the darkness overshadows my soul. Your cross will give me light. When my strength begins to fail, your cross will fortify me. When the path of my life starts to deviate, your cross will show me the way.
\medbreak

I will remain under your Cross, my Lord and Redeemer, when people will not welcome me into their homes and hearts. Then, I will place my head and heart on the wood of your cross.
\medbreak

My Lord and Redeemer, always and everywhere I want to live and die under your cross.
\medbreak

Under your cross, I wish to shed my bitter tears for my sins. On your cross, you have wrought the salvation of the human race. Under your cross, I also desire to work and shed tears for my salvation.
\medbreak
 \newpage 

\subsection[Prayer to Jesus Christ]{Prayer to Jesus Christ \\ {\small\normalfont (St. Augustine)}}

\begin{litany}\prayer{}{Lord Jesus, let me know myself and know Thee,\par And desire nothing save only Thee.\par Let me hate myself and love Thee.\par Let me do everything for the sake of Thee.\par Let me humble myself and exalt Thee.\par Let me think of nothing except Thee.\par Let me die to myself and live in Thee.\par Let me accept whatever happens as from Thee.\par Let me banish self and follow Thee,\par And ever desire to follow Thee.\par Let me fly from myself and take refuge in Thee,\par That I may deserve to be defended by Thee.\par Let me fear for myself, let me fear Thee,\par And let me be among those who are chosen by Thee.\par Let me distrust myself and put my trust in Thee.\par Let me be willing to obey for the sake of Thee.\par Let me cling to nothing save only to Thee,\par And let me be poor because of Thee.\par Look upon me, that I may love Thee.\par Call me that I may see Thee,\par And for ever enjoy Thee. Amen.}\end{litany}

\subsection{Various Prayers}

My good and dearest Jesus, lo, I kneel before you beseeching and praying to you with all the ardour of my soul to engrave deep and vivid impressions of faith, hope, and charity upon my heart, with true repentance for my sins, and a very firm resolve to make amends. Meanwhile I ponder over your five wounds, dwelling on them with deep compassion and grief, and recalling the words that the prophet David wrote long ago, good Jesus, concerning yourself: They have pierced my hands and my feet; they have counted all my bones!
\medbreak
\hrline\medbreak
 \newpage \medbreak

O Lord Jesus Christ, may your Passion strengthen, protect and defend me. May your Wounds be my food and drink, which refresh me, inebriate me and fill me with joy. May the sprinkling of your Blood, wash me from all my sins. May your death be for me a source of inexhaustible life. May your Cross be my eternal glory. In all of this, may I find strength, joy, health and sweetness of heart. For you live and reign for ever and ever. Amen
\medbreak
\hrline\medbreak

O, Lord God, we pray, give the grace of perfect love to your servant. Through the merits of your passion, death and resurrection, through the intercession of the Mother of God, ever Virgin, preserve us, our loved ones, our families and benefactors and our home in the grace of sanctity. Dispel from us all evil.
\medbreak

Never allow us to be overcome by temptation and rescue us from all dangers. Watch over the Church, our Order, our nation, our parents, friends and all who are in need of your mercy.
\medbreak

Guard our Holy Father, all the bishops of the world and especially our Father General. To all the faithful departed, grant eternal peace in your Kingdom. May we all praise you in heaven with your Son and our Lord, Jesus Christ, who lives and reigns with you in the unity of the Holy Spirit, one God, for ever and ever. Amen.
\medbreak
\hrline\medbreak

O, Lord Jesus Christ, Master of my life! Remove from me the spirit of laziness, depression and empty talkativeness. Grant to your servant a spirit of joy, peace, patience and love. Lord, allow me to see my sins and never judge my brother. For you are blessed, for ever and ever. Amen.
\medbreak
 \newpage 

\subsection{Prayers for humility}

\subsubsection{Invocations for humility}

\begin{litany}From the desire for an easy life, \responce{Deliver me O Jesus.}\par 
From the desire of voicing my opinion, \par
From the desire of justifying myself, \par
From the desire of being honoured, \par
From the desire of being loved, \par
From the desire of being respected, \par
From the desire of being praised, \par
From the desire of being understood, \par
From the desire of being sought after, \par
From the desire of being raised above others, \par
From the desire of sparing myself hardships, \par
From the desire of seeing my will done, \par
From the desire of being asked for advice, \par
From the desire of being thought well of, \par
From the desire of having my opinion taken notice of, \par
From the desire of being needed, \par
From the fear of failure, \par
From the fear of difficulties, \par
From the fear of criticism, \par
From the fear of humiliation, \par
From the fear of insults, \par
From the fear of being forgotten, \par
From the fear of being removed, \par
From the fear of injustice, \par
From the fear of being judged, \par
From the fear of being falsely accused, \par
From the fear of being jeered at, \par
From the fear of being mocked, \par
From the fear of being calumniated, \par
From the fear of being insulted, \par
From the fear of being scorned, \par
\end{litany}
\medbreak
\begin{litany}That others may be loved more than me, \responce{Teach me to desire it, my Jesus.}\par 
That others may be more respected than me, \par
That others may gain, while I continually lose importance in the sight of all, \par
That others may be noticed, and that I be always unnoticed, \par
That others may be praised and that I be forgotten, \par
That others may be raised above me in everything, \par
That others may be holier than me, and that I be such as I can be by God's grace, \par
\end{litany}
\medbreak
\begin{samepage}
\prayer{L.}{\textbf{Jesus meek and humble of Heart.}}
\prayer{All}{Make our hearts like unto Thine.}
\end{samepage}
\medbreak
\begin{samepage}
\prayer{L.}{\textbf{Let us pray:}\\ Lord you resist the proud while you give grace to the humble, grant us the virtue of true humility of which your Only-Begotten Son gave an example. Never allow us to arouse your wrath through our pride, but grant that through our humility, we may obtain the gift of your grace. Through Christ our Lord.}
\prayer{All}{Amen.}
\end{samepage}
\medbreak

\subsubsection{Litany for humility}

\begin{litany}O Jesus! meek and humble of heart, \responce{Hear me.}\par 
\end{litany}
\begin{litany}From the desire of being esteemed, \responce{Deliver me Jesus.}\par 
From the desire of being loved, \par
From the desire of being extolled, \par
From the desire of being honoured, \par
From the desire of being praised, \par
From the desire of being preferred to others, \par
From the desire of being consulted, \par
From the desire of being approved, \par
From the fear of being humiliated, \par
From the fear of being despised, \par
From the fear of suffering rebukes, \par
From the fear of being calumniated, \par
From the fear of being forgotten, \par
From the fear of being ridiculed, \par
From the fear of being wronged, \par
From the fear of being suspected, \par
\end{litany}
\medbreak
\begin{litany}That other may be loved more than I, \responce{Jesus grant me the grace to desire it.}\par 
That others may be esteemed more than I, \par
That in the opinion of the world others may increase, and I may decrease, \par
That others may be chosen and I may be set aside, \par
That others may be praised and I may be unnoticed, \par
That others may be preferred to me in everything, \par
That others become holier than I, provided that I may become as holy as I should, \par
\end{litany}
\medbreak

\subsubsection[Prayer for humility]{Prayer for humility \\ {\small\normalfont (composed by Mother Teresa of Calcutta)}}

O, Jesus, your heart was filled with silence and humility. I humbly bow before you with a true sorrow as I recognise the devastation caused by the pride of my heart.
\medbreak

Lord, in your infinite mercy, give me the grace that I would not speak much about myself, but about your Gospel message, that I may not give in to curiosity and not overburden others with my problems or try to control their lives.
\medbreak

Lord, Jesus, I am sorry if I ever have yielded to the temptation of sadness and self-pity, doubt or discouragement. I desire to accept with joy and even with humour all the contradictions of life, which I have to face day by day. I cannot pay much attention to the mistakes of others, but have to remember the fact that we are all weak.
\medbreak

O, Jesus Christ, my Lord! You did not respond or defend yourself when you were falsely accused. Give me the ability to welcome bitter remarks, unjust and unfair accusations, offences and even insults, as if they were veritable pearls. May my will be free from personal, selfish stubbornness, and be willing to fulfil the wishes of others.
\medbreak

In the past, I have been accepting of the praises of others and enjoyed being seen and acknowledged. O, my Master, help me from now on to see the light of holiness, often hidden in the shadow of humility and to cherish those moments of my life, marked with insolence, neglect, ignorance or ingratitude.
\medbreak

O, my Lord! Up to now I have been kind towards people who love and respect me. Free me of this weakness and make me well-disposed towards all with kindness, meekness and generosity, especially towards those who are less lovable and ungrateful.
\medbreak

O, Holy Spirit of truth and love! You imprint the image of Jesus in our souls. Bestow upon me the grace to fall in love with the modesty of my Saviour, Jesus Christ. He did not try to please the people with his glory or to impress them with his majesty and power. He took upon himself our human nature, which we have wounded by sin.
\medbreak

I humbly ask you, inspire me with the longing to avoid everything that would inspire people to value me, to love me, to admire me or to follow me, because of my personality or achievements. May everything I do and talk about, inspire others to love more and more the Sacred Heart of Jesus and the Immaculate Heart of Mary.
\medbreak

O Divine Master! You are the true friend of sinners, of tax-collectors, of abandoned and forgotten children, of lonely strangers, of broken hearts, and of all whose lives may be compared to a dimming lamp or a swaying reed.
\medbreak

O, Divine Master! You hide your divine majesty under the veil of the suffering face of the Man of Sorrows! May I never abuse my human dignity or personal authority, to justify my decisions or actions. Instead, may I offer everything for your glory!
\medbreak

For you live and reign for ever and ever. Amen.
\medbreak
\newpage
\newpageheader{Life of Prayer and Sacrifice with Christ}{Litanies}

\section{Litanies}

\label{litanieshtmlholynamejesus}

\subsection{Litany of the Most Holy Name of Jesus}

\begin{litany}Lord, have mercy, \responce{Lord, have mercy.}\par 
\end{litany}
\begin{litany}Christ, have mercy, \responce{Christ, have mercy.}\par 
\end{litany}
\begin{litany}Lord, have mercy, \responce{Lord, have mercy.}\par 
\end{litany}
\begin{litany}Christ, hear us, \responce{Christ, graciously hear us.}\par 
\end{litany}
\begin{litany}God the Father of Heaven, \responce{Have mercy on us.}\par 
God the Son, the Saviour of the World, \par
God the Holy Spirit, \par
Holy Trinity, One God, \par
\end{litany}
\medbreak
\begin{litany}Jesus, Son of the living God, \par
Jesus, splendour of the Father, \par
Jesus, brightness of eternal light, \par
Jesus, King of glory, \par
Jesus, Son of the Virgin Mary, \par
Jesus, most amiable, \par
Jesus, most admirable, \par
Jesus, mighty God, \par
Jesus, Father of the world to come, \par
Jesus, Angel of the great counsel, \par
Jesus, most powerful, \par
Jesus, most patient, \par
Jesus, most obedient, \par
Jesus, meek and humble of heart, \par
Jesus, lover of chastity, \par
Jesus, lover of us, \par
Jesus, God of peace, \par
Jesus, author of life, \par
Jesus, model of virtues \par
Jesus, zealous for souls, \par
Jesus, our God, \par
Jesus, our refuge, \par
Jesus, Father of the poor, \par
Jesus, treasure of the faithful, \par
Jesus, Good Shepherd, \par
Jesus, true light, \par
Jesus, eternal wisdom, \par
Jesus, infinite Goodness, \par
Jesus, our Way and our Life, \par
Jesus, joy of angels, \par
Jesus, King of Patriarchs, \par
Jesus, Master of the Apostles, \par
Jesus, Teacher of Evangelists, \par
Jesus, Strength of the Martyrs, \par
Jesus, light of Confessors, \par
Jesus, purity of virgins, \par
Jesus, Crown of all Saints, \par
\end{litany}
\begin{litany}Be merciful, \responce{Spare us, O Jesus.}\par 
\end{litany}
\begin{litany}Be merciful, \responce{Graciously hear us, O Jesus,}\par 
\end{litany}
\begin{litany}From all evil, \responce{Jesus, deliver us.}\par 
From all sin, \par
From your wrath, \par
From the snares of the devil, \par
From the spirit of fornication, \par
From everlasting death, \par
From neglect of your inspirations, \par
Through the mystery of you holy Incarnation, \par
Through your Nativity, \par
Through your Infancy, \par
Through your most holy Life, \par
Through your Labours, \par
Through your Agony and Passion, \par
Through your Cross and dereliction, \par
Through your faintness and weariness, \par
Through your Death and Burial, \par
Through your Resurrection, \par
Through your Ascension, \par
Through your institution of the Most Holy Eucharist, \par
Through your joys, \par
Through your glory, \par
\end{litany}
\medbreak
\noindent\begin{minipage}{\textwidth}
\begin{litany}Lamb of God, you take away the sins of the world, \responce{Spare us, O Lord.}\par 
\end{litany}
\begin{litany}Lamb of God, you take away the sins of the world, \responce{Graciously hear us, O Lord.}\par 
\end{litany}
\begin{litany}Lamb of God, you take away the sins of the world, \responce{Have mercy on us.}\par 
\end{litany}
\end{minipage}
\medbreak
\begin{samepage}
\prayer{L.}{\textbf{Jesus, hear us.}}
\prayer{All}{Jesus, graciously hear us,}
\end{samepage}
\medbreak
\begin{samepage}
\prayer{L.}{\textbf{Let us pray:}\\ O Lord Jesus Christ, who has said: \textit{``Ask and you shall receive; seek, and you shall find; knock, and it shall be opened unto you''}; grant to us who ask, we beseech you, the gift of your most divine love, that we may ever love you with all our heart, and in all our words and actions, and never cease from praising you, who live and reign, world without end.}
\prayer{All}{Amen.}
\end{samepage}
\medbreak

\instruction{or:}
\medbreak
\begin{samepage}
\prayer{L.}{\textbf{Let us pray:}\\Make us, O Lord, to have both a perpetual fear and love of your holy name, for you never fail to govern those, whom you founded upon the strength of your love, who live and reign, world without end.}
\prayer{All}{Amen.}
\end{samepage}
\medbreak

\label{litanieshtmlsacredheart}

\subsection{Litany of the Sacred Heart of Jesus}

\begin{litany}Lord, have mercy, \responce{Lord, have mercy.}\par 
\end{litany}
\begin{litany}Christ, have mercy, \responce{Christ, have mercy.}\par 
\end{litany}
\begin{litany}Lord, have mercy, \responce{Lord, have mercy.}\par 
\end{litany}
\begin{litany}Christ, hear us, \responce{Christ, graciously hear us.}\par 
\end{litany}
\begin{litany}God the Father of Heaven, \responce{Have mercy on us.}\par 
God the Son, Redeemer of the World, \par
God the Holy Spirit, \par
Holy Trinity, One God, \par
\end{litany}
\medbreak
\begin{litany}Heart of Jesus, Son of the Eternal Father, \par
Heart of Jesus, formed by the Holy Spirit in the womb of the Virgin Mother, \par
Heart of Jesus, substantially united to the Word of God, \par
Heart of Jesus, of infinite majesty, \par
Heart of Jesus, holy Temple of God, \par
Heart of Jesus, tabernacle of the Most High, \par
Heart of Jesus, house of God and gate of heaven, \par
Heart of Jesus, burning furnace of charity, \par
Heart of Jesus, vessel of justice and love, \par
Heart of Jesus, full of goodness and love, \par
Heart of Jesus, abyss of all virtues, \par
Heart of Jesus, most worthy of all praise, \par
Heart of Jesus, king and centre of all hearts, \par
Heart of Jesus, in whom are all the treasures of wisdom and knowledge, \par
Heart of Jesus, in whom dwells the fullness of the divinity, \par
Heart of Jesus, in whom the Father was well pleased, \par
Heart of Jesus, of whose fullness we have all received, \par
Heart of Jesus, desire of the everlasting hills, \par
Heart of Jesus, patient and abounding in mercy, \par
Heart of Jesus, enriching all who invoke Thee, \par
Heart of Jesus, fountain of life and holiness, \par
Heart of Jesus, propitiation for our sins, \par
Heart of Jesus, filled with reproaches, \par
Heart of Jesus, bruised for our offences, \par
Heart of Jesus, obedient unto death, \par
Heart of Jesus, pierced with a lance, \par
Heart of Jesus, source of all consolation, \par
Heart of Jesus, our life and resurrection, \par
Heart of Jesus, our peace and reconciliation, \par
Heart of Jesus, victim for our sins, \par
Heart of Jesus, salvation of those who hope in Thee, \par
Heart of Jesus, hope of those who die in Thee, \par
Heart of Jesus, delight of all the Saints, \par
\end{litany}
\medbreak
\noindent\begin{minipage}{\textwidth}
\begin{litany}Lamb of God, you take away the sins of the world, \responce{Spare us, O Lord.}\par 
\end{litany}
\begin{litany}Lamb of God, you take away the sins of the world, \responce{Graciously hear us, O Lord.}\par 
\end{litany}
\begin{litany}Lamb of God, you take away the sins of the world, \responce{Have mercy on us.}\par 
\end{litany}
\end{minipage}
\medbreak
\begin{samepage}
\prayer{L.}{\textbf{Jesus meek and humble of heart}}
\prayer{All}{Make our hearts like unto Thine.}
\end{samepage}
\medbreak
\begin{samepage}
\prayer{L.}{\textbf{Let us pray:}\\ Almighty and eternal God, look upon the Heart of Thy most beloved Son and upon the praises and satisfaction which He offers Thee in the name of sinners; and to those who implore Thy mercy, in Thy great goodness, grant forgiveness in the name of the same Jesus Christ, Thy Son, who lives and reigns with Thee and the Holy Spirit, one God, for ever and ever.}
\prayer{All}{Amen.}
\end{samepage}
\medbreak

\label{litanieshtmlbloodjesus}

\subsection{Litany of the Most Precious Blood of Jesus}

\begin{litany}Lord, have mercy, \responce{Lord, have mercy.}\par 
\end{litany}
\begin{litany}Christ, have mercy, \responce{Christ, have mercy.}\par 
\end{litany}
\begin{litany}Lord, have mercy, \responce{Lord, have mercy.}\par 
\end{litany}
\begin{litany}Christ, hear us, \responce{Christ, graciously hear us.}\par 
\end{litany}
\begin{litany}God the Father of Heaven, \responce{Have mercy on us.}\par 
God the Son, Redeemer of the World, \par
God the Holy Spirit, \par
Holy Trinity, One God, \par
\end{litany}
\medbreak
\begin{litany}Blood of Christ, only-begotten Son of the Eternal Father, \responce{Save us.}\par 
Blood of Christ, incarnate Word of God, \par
Blood of Christ, of the New and Eternal Covenant, \par
Blood of Christ, falling upon the earth in Agony, \par
Blood of Christ, shed profusely at the Scourging, \par
Blood of Christ, flowing forth in the Crowning with thorns, \par
Blood of Christ, poured out on the Cross, \par
Blood of Christ, Price of our salvation, \par
Blood of Jesus, without which there is no forgiveness, \par
Blood of Jesus, Eucharistic drink and refreshment of souls, \par
Blood of Jesus, river of mercy, \par
Blood of Jesus, Victor over demons, \par
Blood of Jesus, Courage of Martyrs, \par
Blood of Jesus, Strength of Confessors, \par
Blood of Jesus, bringing forth virgins, \par
Blood of Jesus, Help of those in peril, \par
Blood of Jesus, Relief of the burdened, \par
Blood of Jesus, Solace in sorrow, \par
Blood of Jesus, Hope of the penitent, \par
Blood of Jesus, Consolation of the dying, \par
Blood of Jesus, Peace and Tenderness of our hearts, \par
Blood of Jesus, Pledge of Eternal Life, \par
Blood of Jesus, freeing souls from Purgatory, \par
Blood of Jesus, most worthy of all glory and honour, \par
\end{litany}
\medbreak
\noindent\begin{minipage}{\textwidth}
\begin{litany}Lamb of God, you take away the sins of the world, \responce{Spare us, O Lord.}\par 
\end{litany}
\begin{litany}Lamb of God, you take away the sins of the world, \responce{Graciously hear us, O Lord.}\par 
\end{litany}
\begin{litany}Lamb of God, you take away the sins of the world, \responce{Have mercy on us.}\par 
\end{litany}
\end{minipage}
\medbreak
\begin{samepage}
\prayer{L.}{\textbf{You have redeemed us, O Lord, in your Blood.}}
\prayer{All}{And made of us a kingdom for our God.}
\end{samepage}
\medbreak
\begin{samepage}
\prayer{L.}{\textbf{Let us pray:}\\ Almighty and Eternal God, you have appointed your only begotten Son the Redeemer of the world, and willed to be appeased by his Blood. Grant, we beseech you, that we may worthily adore this price of our salvation, and through its power, be safeguarded from the evils of this present life, so that we may rejoice in its fruits forever in Heaven. Through the same Christ Our Lord.}
\prayer{All}{Amen.}
\end{samepage}
\medbreak

\label{litanieshtmlwoundsjesus}

\subsection{Litany of the Five Wounds of the Lord Jesus Christ}

\begin{litany}Lord, have mercy, \responce{Lord, have mercy.}\par 
\end{litany}
\begin{litany}Christ, have mercy, \responce{Christ, have mercy.}\par 
\end{litany}
\begin{litany}Lord, have mercy, \responce{Lord, have mercy.}\par 
\end{litany}
\begin{litany}Christ, hear us, \responce{Christ, graciously hear us.}\par 
\end{litany}
\begin{litany}God the Father of Heaven, \responce{Have mercy on us.}\par 
God the Son, Redeemer of the World, \par
God the Holy Spirit, \par
Holy Trinity, One God, \par
\end{litany}
\medbreak
\begin{litany}Wounds of Christ, which have been divinised, \par
Wounds of Christ, of infinite mercy, \par
Wounds of Christ, giving us peace, \par
Wounds of Christ, propitiation for our sins, \par
Wounds of Christ, sweetness of our hearts, \par
Wounds of Christ, bleeding rubies, decorating the body of Christ, \par
Wounds of Christ, the greatest lights, \par
Wounds of Christ, New Covenant ratified in the Blood of the Lamb, \par
Wounds of Christ, Sign of Divine Wisdom, \par
Wounds of Christ, Gate of Heaven, \par
Wounds of Christ, Spring of the living waters, \par
Wounds of Christ, our Life and Resurrection, \par
Wounds of Christ, full of the Redeemer's merits, \par
Wounds of Christ, inexhaustible treasure of the Redeemer's grace, \par
Wounds of Christ, our peace and reconciliation, \par
Wounds of Christ, defence of all who put their trust in you, \par
Wounds of Christ, our Redemption, \par
Wounds of Christ, Rock of faith, \par
Wounds of Christ, Sign of victory of the Lamb, \par
Wounds of Christ, Everlasting glory, \par
Wounds of Christ, Mirror of God's love, \par
Wounds of Christ, Mistresses of perfection, \par
Wounds of Christ, Refuge of sinners, \par
Wounds of Christ, Five Stones of David, \par
Wounds of Christ, Five precincts of Jerusalem's pool, \par
Wounds of Christ, Health of the sick, \par
Wounds of Christ, joy of the afflicted, \par
Wounds of Christ, Commemoration of the faithful, \par
Wounds of Christ, Joy of the Angels, \par
Wounds of Christ, Happiness of Patriarchs, \par
Wounds of Christ, Promise of the Prophets, \par
Wounds of Christ, Fortress of Apostles, \par
Wounds of Christ, Victory of Martyrs, \par
Wounds of Christ, Teacher of the Lord's disciples, \par
Wounds of Christ, Hope of Confessors, \par
Wounds of Christ, Embellishments of Virgins, \par
Wounds of Christ, Crown of all Saints, \par
\end{litany}
\medbreak
\begin{litany}From every sin, \responce{Wounds of Christ, deliver us.}\par 
From a sudden and unexpected death, \par
From hunger, plague, fire and war, \par
From the wrath of God, \par
From the spirit of impurity, \par
From Satan's snares and temptations, \par
From eternal damnation, \par
We sinners, cry to you, \par
\end{litany}
\medbreak
\noindent\begin{minipage}{\textwidth}
\begin{litany}Lamb of God, you take away the sins of the world, \responce{Spare us, O Lord.}\par 
\end{litany}
\begin{litany}Lamb of God, you take away the sins of the world, \responce{Graciously hear us, O Lord.}\par 
\end{litany}
\begin{litany}Lamb of God, you take away the sins of the world, \responce{Have mercy on us.}\par 
\end{litany}
\end{minipage}
\medbreak
\begin{samepage}
\prayer{L.}{\textbf{They have pierced my hands and my feet.}}
\prayer{All}{They have counted all my bones.}
\end{samepage}
\medbreak
\begin{samepage}
\prayer{L.}{\textbf{Let us pray:}\\ Almighty and Eternal God, out of obedience, your Son, Jesus Christ, hung upon the Cross for our salvation. Look down upon his Five Wounds, which are the signs of his passion, death and glorious resurrection. Grant, that we may obtain the grace of his Redemption. For you live and reign for ever and ever.}
\prayer{All}{Amen.}
\end{samepage}
\medbreak
\newpage
\includeimage{our-lady-queen.pdf}\pagestyle{empty}
\hspace{0pt}
\vfill

\textit{``It is our conviction that in the course of some very turbu-
lent history associated with our Order (of which there is no lack),
the protection of Mary has been a support in striving for holiness
and making effective our mission to those peoples to whom Divine
Providence has sent us'' (art. 60).}
\medbreak
O, Lord, our God, may the intercession of the Most Holy
Virgin Mary protect us always, that freed from all dangers, we may
rejoice in your peace. Through Christ our Lord. Amen.
\vfill
\textit{A painting,
produced in the workshop of Jasna Gora, XVII century. From the
Art collections of Jasna Gora.}
\hspace{0pt}
\newpage
\newpageheader{Mary, Mother and Patroness of the Order}{Litany of the Blessed Virgin Mary}

\chapter{Mary, Mother and Patroness of the Order}

\label{loretohtmllitany}

\section[Litany of the Blessed Virgin Mary]{Litany of the Blessed Virgin Mary \\ {\small\normalfont (The Litany of Loreto)}}

\begin{litany}Lord, have mercy, \responce{Lord, have mercy.}\par 
\end{litany}
\begin{litany}Christ, have mercy, \responce{Christ, have mercy.}\par 
\end{litany}
\begin{litany}Lord, have mercy, \responce{Lord, have mercy.}\par 
\end{litany}
\begin{litany}Christ, hear us, \responce{Christ, graciously hear us.}\par 
\end{litany}
\begin{litany}God the Father of Heaven, \responce{Have mercy on us.}\par 
God the Son, Redeemer of the World, \par
God the Holy Spirit, \par
Holy Trinity, One God, \par
\end{litany}
\medbreak
\begin{litany}Holy Mary, \responce{Pray for us.}\par 
Holy Mother of God, \par
Holy Virgin of virgins, \par
Mother of Christ, \par
Mother of the Church, \par
Mother of our Order \par
Mother of Divine Grace, \par
Mother Most Pure, \par
Mother Most Chaste, \par
Mother Inviolate, \par
Mother Undefiled, \par
Mother most Amiable, \par
Mother most Admirable, \par
Mother of Good Counsel, \par
Mother of our Creator, \par
Mother of our Saviour, \par
Virgin most Prudent, \par
Virgin most Venerable, \par
Virgin most Renowned, \par
Virgin most Powerful, \par
Virgin most Merciful, \par
Virgin most Faithful, \par
Mirror of Justice, \par
Seat of Wisdom, \par
Cause of our Joy, \par
Spiritual Vessel, \par
Vessel of Honour, \par
Singular Vessel of Devotion, \par
Mystical Rose, \par
Tower of David, \par
Tower of Ivory, \par
House of Gold, \par
Ark of the Covenant, \par
Gate of Heaven, \par
Morning Star, \par
Health of the Sick, \par
Refuge of Sinners, \par
Comforter of the Afflicted, \par
Help of Christians, \par
Queen of Angels, \par
Queen of Patriarchs, \par
Queen of Prophets, \par
Queen of Apostles, \par
Queen of Martyrs, \par
Queen of Confessors, \par
Queen of Hermits, \par
Queen of Virgins, \par
Queen of All Saints, \par
Queen conceived without Original Sin, \par
Queen assumed into Heaven, \par
Queen of the Most Holy Rosary, \par
Queen of Families, \par
Queen of Peace, \par
Queen of Poland, \par
Queen of the World, \par
\end{litany}
\medbreak
\noindent\begin{minipage}{\textwidth}
\begin{litany}Lamb of God, you take away the sins of the world, \responce{Spare us, O Lord.}\par 
\end{litany}
\begin{litany}Lamb of God, you take away the sins of the world, \responce{Graciously hear us, O Lord.}\par 
\end{litany}
\begin{litany}Lamb of God, you take away the sins of the world, \responce{Have mercy on us.}\par 
\end{litany}
\end{minipage}
\medbreak
\begin{samepage}
\prayer{L.}{\textbf{Pray for us, O Holy Mother of God.}}
\prayer{All}{That we may be made worthy of the promises of Christ.}
\end{samepage}
\medbreak
\begin{samepage}
\prayer{L.}{\textbf{Let us pray:}\\ Grant we beseech Thee, O Lord God, that we your servants may enjoy perpetual health of mind and body, and by the glorious intercession of the Blessed Mary, ever Virgin, be delivered from present sorrow and enjoy everlasting happiness. Through Christ our Lord.}
\prayer{All}{Amen.}
\end{samepage}
\medbreak
\newpageheader{Mary, Mother and Patroness of the Order}{Prayer of St. Bernard (The Memorare)}

\section[Prayer of St. Bernard]{Prayer of St. Bernard \\ {\small\normalfont (The Memorare)}}

Remember, O most gracious Virgin Mary, that never was it known that anyone who fled to thy protection, implored thy help, or sought thy intercession, was left unaided. Inspired by this confidence, I fly unto thee, O Virgin of virgins, my Mother. To thee do I come, before thee I stand, sinful and sorrowful. O Mother of the Word Incarnate, despise not my petitions, but in thy mercy, hear and answer me. Amen.
\medbreak
\newpage
\includeimage{our-lady-jasna-gora.pdf}\pagestyle{empty}
\hspace{0pt}
\vfill

\textit{``Jasna Gora and other Marian sanctuaries, which play an
important role in the spiritual formation of people...constitute
something of great value. They endow our Order with a special
charism, as well as a serious mission within our ministry'' (Art.
62).}
\medbreak
Grant, we beseech Thee, O Lord God, that we thy servants may enjoy perpetual health of mind and body, and by the glorious intercession of blessed Mary, ever Virgin, may we be freed from present sorrow, and rejoice in eternal happiness. Through Christ our Lord. Amen.
\vfill
\textit{Our Lady of Jasna Góra in the company of St Paul and St Thomas. The painting studio of Jasna Góra, from the collections of Jasna Góra, from the XVII century.}
\hspace{0pt}
\newpage
\newpageheader{Mary, Mother and Patroness of the Order}{Litany of Our Blessed Virgin Mary ``Kleczki Jasnagorskie''}

\section{Litany of the Blessed Virgin Mary \\``Kleczki Jasnagorskie''}

\gresetlastline{justified}
\gresetinitiallines{0}
\gresethyphen{force}
\gregorioscore{music/kyrie-i}
\begin{nstabbing}
\>Fi-\>li, Redemptor mundi \>De-\>us,\\
\>Spi-\>ritus Sancte \>De-\>us,\\
\>San-\>cta Trinitas, unus \>De-\>us,\\
\end{nstabbing}
\gregorioscore{music/kyrie-ii}\medbreak
\begin{litany}\underline{Sa}ncta Dei \underline{Ge}netrix, \responce{ora pro nobis.}\par 
\underline{Sa}ncta Virgo \underline{vi}rginum, \par
\underline{Ma}ter \underline{Chri}sti, \par
\underline{Ma}ter Ecc\underline{les}iae, \par
\underline{Ma}ter Ordinis \underline{no}stri \par
\underline{Ma}ter Divinae \underline{gra}tiae, \par
\underline{Ma}ter pu\underline{ris}sima, \par
\underline{Ma}ter cas\underline{tis}sima, \par
\underline{Ma}ter invio\underline{la}ta, \par
\underline{Ma}ter inteme\underline{ra}ta, \par
\underline{Ma}ter a\underline{ma}bilis, \par
\underline{Ma}ter admi\underline{ra}bilis, \par
\underline{Ma}ter boni Con\underline{si}lii, \par
\underline{Ma}ter Crea\underline{to}ris, \par
\underline{Ma}ter Salva\underline{to}ris, \par
\underline{Vir}go pruden\underline{tis}sima, \par
\underline{Vir}go vene\underline{ran}da, \par
\underline{Vir}go praedi\underline{can}da, \par
\underline{Vir}go \underline{po}tens, \par
\underline{Vir}go \underline{cle}mens, \par
\underline{Vir}go fi\underline{de}lis, \par
\underline{Spe}culum iu\underline{sti}tiae, \par
\underline{Cau}sa nostrae lae\underline{ti}tiae, \par
\underline{Vas} spiritu\underline{a}le, \par
\underline{Vas} hono\underline{ra}bile, \par
\underline{Vas} insigne devoti\underline{o}nis, \par
\underline{Ro}sa \underline{my}stica, \par
\underline{Tur}ris Da\underline{vi}dica, \par
\underline{Tur}ris e\underline{bur}nea, \par
\underline{Do}mus \underline{au}rea, \par
\underline{Foe}deris \underline{ar}ca, \par
\underline{Ia}nua \underline{cae}li, \par
\underline{Ste}lla matu\underline{ti}na, \par
\underline{Sa}lus infir\underline{mo}rum, \par
\underline{Re}fugium pecca\underline{to}rum, \par
\underline{Con}solatrix afflic\underline{to}rum, \par
\underline{Au}xilium Christia\underline{no}rum, \par
\underline{Re}gina Ange\underline{lo}rum, \par
\underline{Re}gina Patriar\underline{cha}rum, \par
\underline{Re}gina Prophe\underline{ta}rum, \par
\underline{Re}gina Aposto\underline{lo}rum, \par
\underline{Re}gina \underline{Mar}tyrum, \par
\underline{Re}gina Confes\underline{so}rum, \par
\underline{Re}gina Eremi\underline{ta}rum, \par
\underline{Re}gina \underline{Vir}ginum, \par
\underline{Re}gina Sanctorum \underline{om}nium, \par
\underline{Re}gina sine labe originali con\underline{cep}ta, \par
\underline{Re}gina in caelum as\underline{sum}pta, \par
\underline{Re}gina Sacratissimi Ro\underline{sa}rii, \par
\underline{Re}gina famili\underline{a}rum, \par
\underline{Re}gina \underline{pa}cis, \par
\underline{Re}gina \underline{Po}loniae \par
\underline{Re}gina \underline{Mu}ndi \par
\end{litany}
\medbreak
\gresetlastline{justified}
\gresetinitiallines{0}
\gresethyphen{force}
\gregorioscore{music/agnus-die}

\subsection{Sub tuum praesidium}

\gresetlastline{justified}
\gresetinitiallines{1}
\ifdefined\musicillumination
\musicillumination{S}
\fi
\gregorioscore{music/sub-tuum-praesidium}
\label{kleczkihtmlmonstrate}

\subsection[Monstra Te]{Monstra Te\\ {\small\normalfont (Cantatur: expansis manibus)}}

\gresetlastline{justified}
\gresetinitiallines{1}
\gresethyphen{force}
\ifdefined\musicillumination
\musicillumination{M}
\fi
\gregorioscore{music/monstra-te}

\subsection{Versus et responsiones}

\gresetlastline{justified}
\gresetinitiallines{0}
\gresethyphen{force}
\gregorioscore{music/versus-et-responsiones-i}\gresetlastline{justified}
\gresetinitiallines{1}
\gresethyphen{force}
\greillumination{ }
\gregorioscore{music/versus-et-responsiones-ii}\gresetlastline{justified}
\gresetinitiallines{1}
\gresethyphen{force}
\greillumination{ }
\gregorioscore{music/versus-et-responsiones-iii}\gresetlastline{justified}
\gresetinitiallines{1}
\gresethyphen{force}
\greillumination{ }
\gregorioscore{music/versus-et-responsiones-iv}

\subsection*{Tempore aestivo \\ {\small\normalfont (Pro serenitate vel pluvia)}}

\addcontentsline{toc}{subsection}{Tempore aestivo 
(Pro serenitate vel pluvia)}

\gresetlastline{justified}
\gresetinitiallines{1}
\gresethyphen{force}
\greillumination{ }
\gregorioscore{music/versus-et-responsiones-v-i}%
\begin{nstabbing}
\>II. pluviam \>con-\>gru-\>en-\>tem.
\end{nstabbing}
\gresetinitiallines{0}
\gregorioscore{music/versus-et-responsiones-v-ii}

\subsection{Orationes}

\gresetinitiallines{1}
\gresethyphen{auto}
\gresetlastline{ragged}
\ifdefined\musicillumination
\musicillumination{O}
\fi
\gregorioscore{music/orationes-i}\gresetinitiallines{1}
\gresethyphen{auto}
\gresetlastline{ragged}
\ifdefined\musicillumination
\musicillumination[36pt]{D}
\fi
\gregorioscore{music/orationes-ii}\gresetinitiallines{1}
\gresethyphen{auto}
\gresetlastline{ragged}
\ifdefined\musicillumination
\musicillumination[36pt]{D}
\fi
\gregorioscore{music/orationes-iii}\gresetinitiallines{1}
\gresethyphen{auto}
\gresetlastline{ragged}
\ifdefined\musicillumination
\musicillumination[36pt]{D}
\fi
\gregorioscore{music/orationes-iv}\gresetinitiallines{1}
\gresethyphen{auto}
\gresetlastline{ragged}
\ifdefined\musicillumination
\musicillumination[36pt]{A}
\fi
\gregorioscore{music/orationes-v-i}\gresetinitiallines{1}
\gresethyphen{auto}
\gresetlastline{ragged}
\ifdefined\musicillumination
\musicillumination[36pt]{D}
\fi
\gregorioscore{music/orationes-v-ii}\gresetinitiallines{1}
\gresethyphen{auto}
\gresetlastline{ragged}
\ifdefined\musicillumination
\musicillumination[36pt]{H}
\fi
\gregorioscore{music/orationes-vi}\gresetinitiallines{0}
\gresethyphen{auto}
\gresetlastline{ragged}
\gregorioscore{music/per-christum}

\subsection{Chorale}

\gresetlastline{justified}
\gresetinitiallines{0}
\gregorioscore{music/o-maria-virgo-honesta}

\subsection[Psalmus 129]{Psalmus 129 \\ {\small\normalfont De profundis}}

\gresetlastline{justified}
\gresetinitiallines{0}
\gresethyphen{force}
\gregorioscore{music/de-profundis}\begin{psalm}Fiant aures tuae intendéntes \red{*}\par
\hspace{1em}in vocem deprecationis meae.\par
\medbreak
Si iniquitátes observáveris, Domine, \red{*}\par
\hspace{1em}Domine, quis sustinébit.\par
\medbreak
Quia apud te propitiátio est, \red{*}\par
\hspace{1em}et timébimus te.\par
\medbreak
Sustínui te, Domine: \red{\GreDagger}\par
\hspace{1em}sustínuit ánima mea in verbo eius, \red{*}\par
\hspace{1em}sperávit ánima mea in Domino.\par
\medbreak
Magis quam custodes auroram, \red{*}\par
\hspace{1em}speret Israel in Domino.\par
\medbreak
Quia apud Dominum misericordia, \red{*}\par
\hspace{1em}et copiosa apud eum redémptio.\par
\medbreak
Et ipse rédimet Israel \red{*}\par
\hspace{1em}ex omnibus iniquitátibus eius.\par
\medbreak
Requiem aeternam \red{*} dona eis Domine.\par
\hspace{1em}Et lux perpetua \red{*} luceat eis.\par
\medbreak
\end{psalm}
\medbreak

\subsection*{Salve Regina}
\addcontentsline{toc}{subsection}{Salve Regina}

\gresetlastline{justified}
\gresetinitiallines{1}
\ifdefined\musicillumination
\musicillumination{S}
\fi
\gregorioscore{music/salve-regina}\medbreak
\begin{samepage}
\prayer{\Vbar.}{\textbf{Ora pro nobis Sancta Dei Genitrix.}}
\prayer{\Rbar.}{Ut digni efficiamur promissionibus Christi.}
\end{samepage}
\medbreak
\begin{samepage}
\prayer{\Vbar.}{\textbf{Oremus:}\\ Omnipotens sempiterne Deus, qui gloriosae Virginis Matris Mariae corpus et animam ut dignum Filii Tui habitaculum effici mereretur, Spiritu Sancto cooperante praeparasti, da, ut cuius commemoratione laetamur eius pia intercessione ab instantibus malis et a morte perpetua libereamur. Per eundem Christum Dominum nostrum. Amen.}
\prayer{\Rbar.}{Amen.}
\end{samepage}
\medbreak
\newpage
\newpageheader{Mary, Mother and Patroness of the Order}{Marian Antiphons}

\section*{Antiphonae Mariales \\ {\small\normalfont (Cantuale Paulinorum, ca. 1526, \\ Jasna Gora Archive, R583 (I-215), p. 15-19)}}

\subsection*{Ave Regina caelorum}
\addcontentsline{toc}{subsection}{Ave Regina caelorum}

\gresetlastline{justified}
\gresetinitiallines{1}
\ifdefined\musicillumination
\musicillumination{A}
\fi
\gregorioscore{music/ave-regina-caelorum}

\subsection*{Regina caeli}
\addcontentsline{toc}{subsection}{Regina caeli}

\gresetlastline{justified}
\gresetinitiallines{1}
\ifdefined\musicillumination
\musicillumination{R}
\fi
\gregorioscore{music/regina-caeli}

\subsection*{Ave spes nostra}
\addcontentsline{toc}{subsection}{Ave spes nostra}

\gresetlastline{justified}
\gresetinitiallines{1}
\ifdefined\musicillumination
\musicillumination{A}
\fi
\gregorioscore{music/ave-spes-nostra}\newpage
\includeimage{czestochowa.pdf}\pagestyle{empty}
\hspace{0pt}
\vfill

\textit{``In the sixth month, the Angel Gabriel was sent by God to
a town in Galilee, called Nazareth, to a Virgin betrothed to a man
named Joseph, of the House of David; and the Virgin’s name was
Mary'' (Lk 1:26-27)}
\medbreak
Lord, through your will the Eternal Word became incarnate
in the womb of the Virgin Mary, grant that we who have confessed
the Saviour to be true God and True Man, may merit a participation
in His Divine Nature, who lives and reigns with you and the Holy
Spirit, one God for ever and ever. Amen.
\vfill
\textit{St. Luke paints Image of Our Lady of Czestochowa. Workshop of
Jasna Gora, Art collections, XVII century}
\hspace{0pt}
\newpage
\newpageheader{Mary, Mother and Patroness of the Order}{The Holy Rosary}

\section{The Holy Rosary}

\instruction{Our religious Constitutions remind us that ``we honour Mary each day by reciting at least one portion of the Holy Rosary and praying the Litany of Loreto'' (norm 66).}
\medbreak

\instruction{From among the fifteen mysteries of the Holy Rosary, only two are not mentioned in the Holy Scriptures. This is a reason why Pope Pius XII called the Rosary ``a summary of the Gospel.'' This is confirmed also by our Constitution: ``Our devotion to the Most Holy Virgin is based on the Gospels and teaching of the Church'' (art. 61)}
\medbreak

\subsection*{Part One --- The Joyful Mysteries}

\subsubsection*{The first joyful mystery \\ {\small\normalfont The Annunciation:}}

``In the sixth month, the Angel Gabriel was sent by God to a town in Galilee called Nazareth, to a Virgin, whose name was Mary. The Angel said to her: `Rejoice, so highly favoured\ldots You are to conceive and bear a Son, and you will name him Jesus' '' (Lk 1:26-31)
\medbreak

\subsubsection*{The second joyful mystery \\ {\small\normalfont The Visitation:}}

``Mary set out at that time and went as quickly as she could to a town in the hill country of Judah\ldots Now as soon as Elizabeth heard Mary's greeting\ldots she gave a loud cry and said, `Of all women you are the most blessed, and blessed is the fruit of your womb' '' (Lk 1:39-42).
\medbreak

\subsubsection*{The third joyful mystery \\ {\small\normalfont The Nativity:}}

``While they were there the time came for her to have her Child, and she gave birth to a Son, her first-born. She wrapped him in swaddling clothes, and laid him in a manger because there was no room for them at the inn'' (Lk 2: 2-7)
\medbreak

\subsubsection*{The fourth joyful mystery \\ {\small\normalfont The Presentation:}}

``And when the day came for them to be purified as laid down by the Law of Moses, they took him up to Jerusalem to present him to the Lord\ldots and also offered in sacrifice\ldots a pair of turtledoves or two young pigeons'' (Lk 2: 22-24)
\medbreak

\subsubsection*{The fifth joyful mystery \\ {\small\normalfont The finding of the Child Jesus in the Temple:}}

``Every year his parents used to go to Jerusalem for the feast of the Passover. When they were on their way home after the feast, the boy Jesus stayed behind in Jerusalem\ldots Three days later, they found him in the Temple, sitting among the doctors, listening to them and asking them questions'' (Lk 2:41-46).
\medbreak

\subsection*{Part Two --- The Mysteries of Light}

\subsubsection*{The first mystery of light \\ {\small\normalfont Christ's Baptism in the Jordan:}}

``The Spirit descended like a dove on him. And a voice came from heaven: You are my Son, the Beloved, with whom I am well pleased.'' (Mk 1:10)
\medbreak

\subsubsection*{The second mystery of light \\ {\small\normalfont Christ's self-revelation at the marriage at Cana:}}

``Jesus changed water into wine, the first of his signs, and revealed his glory, and his disciples believed in him.'' (Jn 2:11).
\medbreak

\subsubsection*{The third mystery of light \\ {\small\normalfont Christ's proclamation of the Kingdom of God with a call to conversion:}}

``The time is fulfilled, and the Kingdom of God has come near, repent, and believe in the Good News.'' (Mk 1:15)
\medbreak

\subsubsection*{The fourth mystery of light \\ {\small\normalfont The Transfiguration:}}

``And he was transfigured before them, and his face shone like the sun, and his clothes became dazzling white.'' (Mt 17: 2)
\medbreak

\subsubsection*{The fifth mystery of light \\ {\small\normalfont The Institution of the Eucharist:}}

``He took bread, gave thanks, broke it and gave it to them, saying, This is my Body, which is given for you. Do this in remembrance of me.'' (Lk 22:19).
\medbreak

\subsection*{Part Three --- The Sorrowful Mysteries}

\subsubsection*{The first sorrowful mystery \\ {\small\normalfont The Agony in the Garden:}}

``He then left to make his way as usual to the Mount of Olives, with the disciples following\ldots Then he withdrew from them\ldots , and knelt down and prayed `Father', he said, `if you are willing, take this cup away from me; yet let your will be done, not mine' '' (Lk 22:41-46)
\medbreak

\subsubsection*{The second sorrowful mystery \\ {\small\normalfont The Scourging at the Pillar:}}

``As you can see, this man has done nothing that deserves death, so I will have him flogged and then let him go\ldots but they kept on shouting at the top of their voices, demanding that he should be crucified'' (Lk 23:22-23).
\medbreak

\subsubsection*{The third sorrowful mystery \\ {\small\normalfont The Crowning with Thorns:}}

``The governor's soldiers took Jesus with them into the praetorium and collected the whole cohort round him. They then stripped him and made him wear a scarlet cloak, and having twisted some thorns into a crown, they put this on his head and placed a reed in his right hand. To make fun of him, they knelt to him saying, `Hail, king of the Jews' '' (Mt 27:27-28).
\medbreak

\subsubsection*{The fourth sorrowful mystery \\ {\small\normalfont The Carrying of the Cross:}}

``And when they had finished making fun of him, they took off the cloak and dressed him in his own clothes and led him away to crucify him'' (Mt 27:31)
\medbreak

\subsubsection*{The fifth sorrowful mystery \\ {\small\normalfont The Crucifixion and Death of our Lord:}}

``When they reached the place called ``The Skull'', they crucified him there and the two criminals also, one on the right, the other on the left. Jesus said, `Father forgive them; they do not know what they are doing.' Then, they cast lots to share his clothing'' (Lk 23:33-34)
\medbreak

\subsection*{Part Four --- The Glorious Mysteries}

\subsubsection*{The first glorious mystery \\ {\small\normalfont The Resurrection:}}

``And all at once there was a violent earthquake, for the Angel of the Lord, descending from heaven, came and rolled away the stone and sat on it \ldots ~But the Angel spoke and said to the women: `There is no need for you to be afraid. I know you are looking for Jesus, who was crucified. He is not here, for he has risen, as he said he would' '' (Mt 28:5-6).
\medbreak

\subsubsection*{The second glorious mystery \\ {\small\normalfont The Ascension:}}

``And so the Lord Jesus, after he had spoken to them, was taken up into heaven: there at the right hand of God, he took his place, while they going out preached everywhere, the Lord working with them and confirming the word by the signs that accompanied it'' (Mk 16:19-20).
\medbreak

\subsubsection*{The third glorious mystery \\ {\small\normalfont The Descent of the Holy Spirit:}}

``When Pentecost day came round, they had all met in one room, when suddenly they heard what sounded like a powerful wind from heaven, the noise of which filled the entire house in which they were sitting; and something appeared to them that seemed like tongues of fire; these separated and came to rest on the head of each of them. They were filled with the Holy Spirit, and began to speak foreign languages as the Spirit gave them the gift of speech'' (Acts 2:1-4).
\medbreak

\subsubsection*{The fourth glorious mystery \\ {\small\normalfont The Assumption of Mary into Heaven:}}

``Finally, the Immaculate Virgin, preserved free from all guilt of original sin, on the completion of her earthly sojourn, was taken up body and soul into heavenly glory, and exalted by the Lord as Queen of the universe, that she might be the more fully conformed to her Son, the Lord of lords and the conqueror of sin and death.'' (LG 59)
\medbreak

\subsubsection*{The fifth glorious mystery \\ {\small\normalfont The Coronation of the Blessed Virgin Mary as the Queen of Heaven:}}

``Now, a great sign appeared in heaven: a Woman, adorned with the sun, standing on the moon, and with the twelve stars on her head for a crown'' (Rev 12:1).
\medbreak
\newpage
\newpageheader{Mary, Mother and Patroness of the Order}{An act of Consecration to the Mother of God}

\section{An act of Consecration to the \\ Mother of God}

Mother of God, Immaculate Mary, to you I consecrate my body and soul, all my prayers, works, joys and sufferings, all that I am and possess. With an eager heart, I dedicate myself to you in the service of love. I give you complete freedom in making use of me for the salvation of men and the assistance of the Holy Church of which you are the Mother.
\medbreak

From now on, I wish to do everything with you, through you and for you. I know that by my own strength, I shall accomplish nothing. You, however, can do everything that is the will of your Son and you are always victorious. Grant then, O Help of Christians, that my family, parish and my country, may truly be yours and your Son's Kingdom. Amen.
\medbreak
\newpageheader{Mary, Mother and Patroness of the Order}{Invocation to Implore the Intercession of the Immaculate Conception}

\section{Invocation to Implore the Intercession of the Immaculate Conception}

O Most Holy Virgin, I believe in your Holy and Immaculate Conception, always pure and undefiled, and I confess it before the whole world.
\medbreak

O Virgin most pure, through your Immaculate Conception and through your great dignity as Mother of God, implore for me, from your Son, the virtue of deep humility, of perfect love, of intact chastity of heart, body and soul, the grace of final perseverance in my beloved vocation, the gift of worthy prayer and the grace of a happy death. Amen.
\medbreak
\newpage
\newpageheader{Mary, Mother and Patroness of the Order}{An act of Reparation to the Immaculate Heart of Mary}

\section{An act of Reparation to the Immaculate Heart of Mary}

O Immaculate Heart of Mary, masterpiece of Almighty God, most honoured vessel of the Most Holy Trinity; it is our greatest joy to venerate and honour you. You have loved us so much! After God, we owe everything to you! You gave us your beloved Son, Jesus. For our sake, you shared all his sufferings. Your suffering was ``as deep as the sea''. But, there are some unhappy souls, who instead of loving you without end, inspired by the devil himself, dare to throw at your Name and your most exalted privileges, their hideous blasphemies.
\medbreak

Every day, millions of Christians through their sins, renew the passion of Christ and pierce your Immaculate Heart. Unfortunately, we, ourselves, through our faults, do the same. Forgive us, O Queen and Mother of Mercy all our infidelities toward you. Today, more than ever before, we promise with the help of God's grace to honour, to love and imitate you. In the spirit of reparation, we renew our total consecration to you.
\medbreak

We offer you all what we are and possess. We want to be true to this consecration by the way we live, by working for an increase of your glory and the kingship of your Son.
\medbreak

For all the blasphemies and insults directed against you by words and deeds, --- we make reparation to your heart, O Mary.
\medbreak

For the profanations of your holy images, --- we make reparation to your heart, O Mary.
\medbreak

For outrages and ingratitude with which we have wounded your Maternal heart --- we make reparation to your heart, O Mary.
\medbreak

For the ignorance and indifference of many Christians, --- we make reparation to your heart, O Mary.
\medbreak

United with all Angels, Saints and chosen souls we honour you in heaven and on earth, --- we make reparation to your heart, O Mary.
\medbreak

Accept our filial homage. Show yourself to be our Mother and make us worthy to be numbered among your servants and children. Amen.
\medbreak
\newpageheader{Mary, Mother and Patroness of the Order}{Filial and Apostolic Consecration to the Heart of Mary}

\section{Filial and Apostolic Consecration to the Heart of Mary}

O Mother of God ever Virgin! I consecrate myself to your maternal love, as your son and servant, that you may form within me your Son Jesus, the messenger of the Father, anointed by the Holy Spirit, to proclaim the Good News to the poor.
\medbreak

I entrust myself into your hands that you may send me to all people; your children and my brothers and sisters. Help me to show them the face of the Father.
\medbreak

Teach me to preserve in my heart the Word of God, as you did, that I may become a living Gospel. Make of me an obedient instrument of your maternal love, that I may bring forth many sons and daughters born through the preaching of the Gospel.
\medbreak
\begin{litany}O Mother, behold your son, send me! \par
O Mother, behold your son, speak through me! \par
O Mother, behold your son, love through me! \par
O Mother, behold your son, watch over me! \par
O Mother, help me that announcing the Gospel to others, I myself may not be excluded from the Kingdom. \par
O Mother, in you I place all my hope and I trust that I may never will be disappointed. \par
Amen. \par
\end{litany}
\medbreak
\newpageheader{Mary, Mother and Patroness of the Order}{Prayer to the Queen of Hermits for our Order}

\section{Prayer to the Queen of Hermits for our Order}

Our Lady of Jasna Gora, Mother of the Church, singular Patroness of the Pauline Order and Queen of Hermits! We wish to express our eager thanksgiving for your blessed presence in our Religious Family. Be ever honoured for being our unfailing Hope and Refuge in the history of the Order, especially in its most difficult moments. Along with you, we praise the Divine Providence, which has entrusted to us many new places for our apostolic works and the establishment of shrines, particularly that of Jasna Gora, as the primary centre of your glory and our Pauline service.
\medbreak

We want to serve you in the Church and our Nation worthily, especially at the shrine of Jasna Gora, where Almighty God has deigned to glorify your miraculous Image with a rare devotion of the faithful. We thank you for your maternal goodness demonstrated in the past, and we ask you today to continue to shower upon us all the graces we need.
\medbreak

Through your maternal presence, make our consciences sensitive so that we may repent of the sins of our past, and help us to make good resolutions for the future.
\medbreak

O Mother, forgive us our negligence in imitating your virtues and our shortcomings in fostering your glory among people.
\medbreak

O Mother, forgive us our unwillingness to entrust ourselves completely to you, without any hesitation.
\medbreak

O Mother, forgive us our weaknesses and defects in straying from the ways of those who are ``totally consecrated to you''.
\medbreak

O Humble Servant of the Lord! Teach us to know Jesus more deeply! Help us to become a ``new creation'', by being faithful to our religious vows and to the spiritual heritage of our Fathers. Help us to build up our fraternal Community in our Religious Houses, by our mutual love, respect and trust. United with our Superiors, may we radiate the light of the Gospel!
\medbreak

O Mother of the Church! Animate us in apostolic zeal! Inspire us with new pastoral initiatives and teach us to join contemplation with evangelization, that we may serve the faithful well in their struggles to enrich their life of faith with greater efficacy.
\medbreak

Finally, O Teacher of many Pauline generations! Ask your Son, the Lord of the Harvest, to send new hosts of workers for service in our Order. May our Novitiates and Seminaries be your schools of Religious and Priestly life. May there be a new Pentecost within our Monastic Family to the glory of the Most Holy Trinity and the joy of the People of God. Amen.
\medbreak
\newpageheader{Mary, Mother and Patroness of the Order}{An Act of Consecration to Mary the Queen and Protectress of the Order}

\section{An Act of Consecration to Mary the Queen and Protectress of the Order}

O Holy Mary, Mother of God, ever Virgin! I Brother \rubric{N. N.} choose you this day as my Mother, Queen and Protectress. I resolve never to abandon you, and I promise to work with every thought, word and action to further your glory. I beg you, to take me as your eternal servant. Be at my side always! Do not forsake me at the hour of my death. Amen.
\medbreak
\newpageheader{Mary, Mother and Patroness of the Order}{An Act of self-offering}

\section{An Act of self-offering}

O my Lady, the only Mother and Queen on earth and in heaven! Behold, I, an unworthy servant, kneel before you and entrust myself entirely to your mighty protection. You were always a mother to me, but especially now in this holy Order. Instruct and defend me as your son. Teach me the virtue of humility and angelic purity. Guide me by the hand amongst many hardships and temptations and never allow me to forget you.
\medbreak

O my Mother! After God you are my only hope. Into your hands I entrust my life and labour. To you, I give my heart, that you may reign over it as my Lady and Queen. You taught me how to love Jesus.
\medbreak

O my Queen! To serve you is my only happiness. I do not have any other desire, but to think, say and do everything that serves your glory. Amen.
\medbreak
\newpageheader{Mary, Mother and Patroness of the Order}{An act of Consecration of Seminarians of the Pauline Order to the Queen of Jasna Gora}

\section[An act of Consecration of Seminarians of the Pauline Order to the Queen of Jasna Gora.]{An act of Consecration of Seminarians of the Pauline Order to the Queen of Jasna Gora. \\ {\small\normalfont (Jasna Gora, 9 September 1961)}}

O powerful Queen of Poland and of the world!\\ An hour of great decision for the whole of humanity today has arrived. It challenges and impels us to fall to our knees before your Miraculous Image and to turn our eyes and hearts toward you.
\medbreak

We, the young members of the Pauline Order at this hallowed place, chosen by you as the Shrine of the Nation, wish to offer you our humble homage of an unwavering faith and deep gratitude, and to express our desire to dedicate our lives to your exclusive service. As Jesus, your Son and our Saviour and Leader, in his time chose his apostles and disciples, likewise today, you invite each one of us personally, to join the ranks of this Order, that carries on the joyous mission, entrusted by God, to proclaim the Kingdom of the Sacred Heart of Jesus on Earth.
\medbreak

We are called to reveal you to the nation and to the world. Our nation and all of humanity, have an obligation, to know you better, and love you more ardently. You are the only true Mother and a great Teacher given to man. Therefore, he has a duty to accept and trust in you, knowing that the salvation of the world came through you.\\ O Our Lady and Queen!
\medbreak

From the bottom of our young hearts, we thank you above all for the gift of our good and believing parents, who enkindled in our hearts a great love for you and for the true Catholic Church.
\medbreak

We thank you for the wonderful and glorious gift of our Monastic and Priestly vocation.
\medbreak

We thank you with all our heart for your constant protection of our Order, especially of our houses of studies.
\medbreak

We thank you for our professors and superiors and for the graces showered upon us with such abundance that help us in achieving our goal: the holy priesthood and holiness. O how much we ought to love and honour you! Forgive us all our former and present transgressions, sins and negligence.
\medbreak

Today, with your help, we want to begin writing a new book of life for our Seminary, for our nation and for the whole world.
\medbreak

O Beloved Lady of Jasna Gora! Embrace all of us, our Religious Family, and particularly our Seminary with your royal heart! Accept us as your perpetual possession, all we are and have: our past, our present and our future. From now on, we are yours and yours alone! Use us well for your own royal designs. Enclose our houses of studies and religious properties within your arms and protect them from all dangers. Remove from us all mediocrity, compromise and discouragement. Excite us with great aspirations and with a steadfast faith in Jesus Christ and his final victory.
\medbreak

Finally, we ask you, our Beloved Mother, that you may teach us how to love Our Lord Jesus Christ with all our hearts, with the example of St. Paul the Apostle, who still cries out: ``The love of Christ urges us\ldots ~Nothing therefore can come between us and the love of Christ, even if we are troubled or worried, or being persecuted, or lacking food or clothes, or being threatened or even attacked\ldots ~For I am certain of this: neither death nor life, no angel, no prince, nothing that exists, nothing still to come, not any power, or height or depth, nor any created thing, can ever come between us and the love of God made visible in Christ Jesus our Lord'' (Rom. 8:35-39). Amen.
\medbreak
\newpageheader{Mary, Mother and Patroness of the Order}{An act of self-offering to the Most Blessed Virgin Mary}

\section[An act of self-offering to the Most Blessed Virgin Mary]{An act of self-offering to the Most Blessed Virgin Mary \\ {\small\normalfont (A prayer of Novices and Seminarians)}}

O My Dearest Mother, My Lady, My Queen and My Hope! I consecrate myself totally to you in the presence of Jesus, My Lord, abiding now in my heart.
\medbreak

I offer to you my body and my soul. I commend to you, My Most Loving Mother, all the activities and sufferings of my life, even the smallest. O Mother, for whom nothing is too small or too great, be my strongest and never failing support and comfort.
\medbreak

I offer my whole life: minute by minute, hour by hour, day by day. In a special way, I commend to you the moment of my death, knowing that you will not fail to come. I offer you my Purgatory and I rejoice at the thought that through my suffering there, I will come to you purified.
\medbreak

I offer you my Heaven, my Eternal Heaven, and I rejoice at the thought of seeing you there, My Dearest Mother, whom I love so dearly and whom I have not yet seen. There, you will make me happy, showing me your face!
\medbreak

O My Beloved Mother, I commend to you my religious vows, that you may preserve them solely for Jesus. In a special way I entrust to you my vow of chastity. As you called yourself a ``Servant of the Lord'', teach me a great sense of humility of mind and heart.
\medbreak

Finally, O Dearest Mother, I ask you with all my heart, even today, to say to My Lord, with your blessed lips these four words: ``Make him a saint.'' O My Mother, I beg you for these four words! It is so easy for you, but for me, it means an infinite benefit.
\medbreak

You know how ardently I desire to become a saint, no matter how much I should suffer. I do not ask for anything else that is in Heaven or on earth. Tell this to Jesus, whom I carry now in my heart.
\medbreak

O My Mother, I offer you also Jesus, who is your most precious Treasure. O My Mother, I desire nothing in this life; neither health, nor happiness nor anything else in the world. Grant me only the grace to love Jesus and you and to be with you in Heaven forever. Amen.
\medbreak
\newpageheader{Mary, Mother and Patroness of the Order}{Renewal of the Covenant with Mary the Mother and Protectress of the Order}

\section[Renewal of the Covenant with Mary the Mother and Protectress of the Order]{Renewal of the Covenant with Mary the Mother and Protectress of the Order \\ {\small\normalfont \rubric{(16\textsuperscript{th}th January)}}}

\instruction{``On the feast of the Queen of Hermits, Mother and Protectress of the Order, observed on January 16th, each Community will renew its consecration to Mary. As part of the consecration, we offer first thanks, we then place ourselves totally at her disposal in our future action for the good of the Church and we beg her protection for ourselves.'' (norm 66a).}
\medbreak

\instruction{During the Conventual Mass presided over by the Prior, after the homily the principal celebrant, joined by all the members of the Monastery on their knees, recite the following prayer:}
\medbreak

O Mother and singular Protectress of our Order! Today on your feast day of the ``Queen of Hermits'', we desire to renew our covenant of love with you. Above all, accept our grateful thanksgiving!
\medbreak

\rubric{We will all repeat:} \textbf{We thank you, O Mary!}
\medbreak
\begin{litany}For the grace that Blessed Eusebius gathered together into one great flame, all your dispersed servants, \responce{We thank you, O Mary!}\par 
For the grace that in your goodness, you have chosen the shrine of Jasna Gora, as a place of your special presence, \responce{We thank you, O Mary!}\par 
For the grace that has called us the sons of St. Paul the First Hermit, to serve the People of God \responce{We thank you, O Mary!}\par 
For the grace that in the most difficult moments, as the Mother of Mercy, you have sustained the very existence of the Order, by showing your unfailing help, \responce{We thank you, O Mary!}\par 
For the grace, that there are always among us your saintly servants, \responce{We thank you, O Mary!}\par 
\end{litany}
\medbreak
\begin{litany}In a spirit of gratitude for the grace of our holy vocation, we renew our dedication to you as your perpetual possession, \responce{We dedicate ourselves to you, O Mary.}\par 
Take into your sovereignty our Order, all its members, in particular our youth, \responce{We dedicate ourselves to you, O Mary.}\par 
In the days of struggles between good and evil, in the hour of great trials, we want to assist you by our faithful love, therefore, \responce{We dedicate ourselves to you, O Mary.}\par 
\end{litany}
\medbreak

Lead each one of us and our whole Pauline Family to the full growth in holiness and apostolic action, in the service of the Church and its saving mission.
\medbreak

Watch over us at every hour of life, but especially in the last hour of our passage into eternity, that with you, we may praise the One God, in the Most Holy Trinity, for ever and ever. Amen.
\medbreak

\instruction{In closing, we may sing the hymn: ``Monstra te esse matrem'' p.~\pageref{kleczkihtmlmonstrate} or another, appropriate hymn.}
\medbreak
\newpageheader{Mary, Mother and Patroness of the Order}{An act of consecration of the Order to the Most Holy Mother of God}

\section[An act of consecration of the Order to the Most Holy Mother of God]{An act of consecration of the Order to the Most Holy Mother of God \\ {\small\normalfont (3\textsuperscript{rd}rd May 1962)}}

\begin{samepage}
\prayer{L.}{\textbf{O Mother of the Church! Remember most gracious Virgin Mary, our Queen, the Mother of our vocations and the Lady of our hearts, that we are a part of the chosen heritage of your Son, Jesus Christ. You were present at the birth of our Order and supported the founding Holy Fathers and patriarchs with your maternal inspirations.\\ O Mediatrix of all graces! Throughout the centuries, with your countless graces, you have generously enriched our brothers.\\ O Good Protectress! With your protective mantle you have shielded us against the coldness and malice of the world with its falsehood and perversity.\\ O Good Mother! Behold, we stand before you once again, and admit, that you are not loved enough in the Religious Families of the Church today. However, following the great and noble example of our national traditions and being aware of the seriousness of the present moment and encouraged by our Church leaders, we wish to foster in our Polish religious houses an intimate, childlike love and attachment for you, our Lady and our Queen. With a spirit of responsibility for the solemn vows of the Nation, pronounced on the threshold of the Great Jubilee of Poland's Baptism, we renew our own religious vows, to live in holiness and to lead the Nation to Christ, the Light. Therefore, today, we come to you again, O Our Lady, with faith and full of confidence, to pay due homage of thanksgiving, honour and dedication.}}
\prayer{All}{We thank you, Our Lady, for all graces that we have received from your maternal hands. We thank you for the grace of our vocation to the consecrated life, for that of our brothers and for those whom you will send to us.\\ We thank you for a thousand years of the consecrated life in our Homeland.\\ We thank you for all our joys and sorrow, for rest and leisure, for the fruits of our labour and the disappointments we have encountered.\\}
\end{samepage}
\medbreak
\begin{samepage}
\prayer{L.}{\textbf{Today, we bring you our humblest homage along with all our brothers.}}
\prayer{All}{Be forever blessed, the Fullness of Grace! Be forever blessed, O Mother of God! Let us praise you, O Holy Virgin! Give us the grace to honour you always, to spread devotion toward you, and to win all hearts to love you.}
\end{samepage}
\medbreak
\begin{samepage}
\prayer{L.}{\textbf{We consecrate and dedicate ourselves to you, Our Queen. To you we wish to belong.}}
\prayer{All}{Therefore, we pray that you may take under your protection our religious families, our brothers, our homes and our labour. In our daily toil, when the devil tries to destroy the Holy Church, the work of your Son, through his falsehood, malice, persecution, sin and all kinds of deceit, with great confidence, we take refuge under your maternal mantle.}
\end{samepage}
\medbreak
\begin{samepage}
\prayer{L.}{\textbf{O Most Holy Servant of the Lord! We place and entrust to you our religious vows and our humble petitions.}}
\prayer{All}{O Mother Most Pure! Grant us and our brothers the strength to preserve our chastity.}
\end{samepage}
\medbreak
\begin{samepage}
\prayer{L.}{\textbf{O Virgin Most Poor!}}
\prayer{All}{Give us and our brothers the grace to love the vow of Evangelical poverty and detachment from the world.}
\end{samepage}
\medbreak
\begin{samepage}
\prayer{L.}{\textbf{O Virgin Most Humble!}}
\prayer{All}{Teach us and our brothers a humble obedience to our superiors.}
\end{samepage}
\medbreak
\begin{samepage}
\prayer{L.}{\textbf{O Virgin Most Powerful!}}
\prayer{All}{Give us strength against the powers of darkness.}
\end{samepage}
\medbreak
\begin{samepage}
\prayer{L.}{\textbf{O Virgin Most faithful!}}
\prayer{All}{Help us to be faithful to our vocation.}
\end{samepage}
\medbreak
\begin{samepage}
\prayer{L.}{\textbf{Queen of Apostles!}}
\prayer{All}{We ask you to give us the spirit of the Apostles. May our lives be a continued new leaven for an increase of the Kingdom of your Son on earth, and especially in our country.}
\end{samepage}
\medbreak
\begin{samepage}
\prayer{L.}{\textbf{Queen of Peace!}}
\prayer{All}{We pray for peace in our homes.}
\end{samepage}
\medbreak
\begin{samepage}
\prayer{L.}{\textbf{O Mother of Beautiful Love!}}
\prayer{All}{Enkindle your love in our hearts.}
\end{samepage}
\medbreak
\begin{samepage}
\prayer{L.}{\textbf{Queen of Heaven!}}
\prayer{All}{Lead us and all our brothers to the Kingdom of Heaven.}
\end{samepage}
\medbreak
\begin{samepage}
\prayer{L.}{\textbf{We fly to your patronage,}}
\prayer{All}{O Holy Mother of God, despise not our petitions in our necessities, but deliver us always from all dangers. O Glorious and Blessed Virgin. We humbly ask you to be our Lady, our Intercessor, and our Advocate. Reconcile us to your Son. Recommend us to your Son. Bring us to your Son.}
\end{samepage}
\medbreak
\newpage
\includeimage{guardian-angel.pdf}\pagestyle{empty}
\hspace{0pt}
\vfill

\textit{``See, I am sending an angel before you, to guard you on
your way and bring you to the place I have prepared. Be attentive
to him and heed his voice'' (Ex 23: 20-21).}
\medbreak
God our Father, in your loving providence, you send your
holy angels to watch over us. Hear our prayers, defend us always
by their protection and let us share your life with them for ever.
Through Christ our Lord. Amen.
\vfill
\textit{A Guardian Angel. Workshop of Jasna Gora. The Jasna Gora Art
Collection, XVII century.}
\hspace{0pt}
\newpage
\newpageheader{Joined to the Guardian Angels}{Invocations to the Holy Angels}

\chapter{Joined to the Guardian Angels}

\instruction{``Consequently, over and above our Marian devotion, we should take care especially to foster devotion to the Holy Angels'' (art. 32). The Pauline tradition and the faith teach us, that liturgy and prayer is a work of the whole Church; both of the pilgrim Church on earth and of the glorious one in heaven. In the Constitution on the Sacred Liturgy of the Second Vatican Council, we read: ``In the earthly liturgy we take part in a foretaste of that heavenly liturgy\ldots with the Heavenly Hosts, we sing a hymn of glory to the Lord'' (Sacrosanctum Concilium 8).}
\medbreak

\instruction{By honouring our Holy Angels, we want to bring our daily life ever closer to the Kingdom of God. Let us keep in mind that the primary obligation of Angels is to worship God. The liturgical texts constantly remind us of this. We are also encouraged to follow the Holy Angels in our worship and praise of God. We do not add anything to God's glory. Our praise and worship of God, contribute to our salvation.}
\medbreak

\label{invocationsangelshtmlinvocations}

\section{Invocations to the Holy Angels}

\begin{litany}Holy Mary, Queen of Angels, \responce{Pray for us.}\par 
Holy Angels, our Guardians, \par
Holy Angels, who sing unceasingly ``Holy'' to God, \par
Holy Angels, who illuminate our thoughts with divine light, \par
Holy Angels, who out of love, watch over people, \par
Holy Angels, who are sent to help those on the way to salvation, \par
Holy Angels, who constantly contemplate the face of our heavenly Father, \par
Holy Angels, who rejoice over every penitent sinner, \par
Holy Angels, who rescued Lot from among those who were doomed, \par
Holy Angels, who proclaimed joy and peace at the Nativity of our Lord, \par
Holy Angels, who attended our Lord in the wilderness, \par
Holy Angels, who appeared at the tomb, \par
Holy Angels, who comforted the disciples at the Ascension of our Lord, \par
Holy Angels, who will separate the good from the evil, \par
Holy Angels, who present our prayers before the throne of God, \par
Holy Angels, who lead the souls of the just to heaven, \par
Holy Angels, who keep guard over nations and countries, \par
Holy Angels, who deliver people from danger, \par
\end{litany}
\medbreak
\begin{litany}We sinners, cry unto you, \responce{Hear us, O Lord.}\par 
Through the prayers of our Holy Guardian Angels, \par
That you may preserve children and youth from sin, \par
That you may strengthen the fidelity of our Order to its vocation, \par
That you may protect the Church of Christ from all the attacks of our enemies, \par
That you may inspire us to a greater devotion to the Holy Angels, \par
That you may give us the grace of final perseverance, \par
\end{litany}
\medbreak
\begin{samepage}
\prayer{L.}{\textbf{In the presence of the Angels I will sing to you.}}
\prayer{All}{And adore before your holy temple.}
\end{samepage}
\medbreak
\begin{samepage}
\prayer{L.}{\textbf{Let us pray:}\\ O God, who in your inexpressible Providence, deigned to send your Holy Angels as our guardians, we humbly ask you, that we may find in their constant protection our defence, and enjoy their company throughout eternity.\\ Through Christ our Lord.}
\prayer{All}{Amen.}
\end{samepage}
\medbreak
 \newpage 

\subsection*{Prayer to the Angels}

Holy choirs of Angels! You were the first from among all creatures to recognize the supreme authority of God over you, pay him homage and adore him; teach us to humble ourselves before God's majesty.
\medbreak

Holy servants of the Lord! You always stand before the throne of God, ready to obey his commands; teach us to remain faithful in fulfilling God's holy will.
\medbreak

St. Michael the Archangel, prince of the heavenly hosts, who with the cry ``Who is like God'' defeated Lucifer and his angels; inspire us with a zeal to work for the greater glory of God.
\medbreak

St. Gabriel the Archangel, who announced to the Holy Virgin Mary that she would become the Mother of the Son of God by the power of the Holy Spirit; illuminate our minds, that like Mary, we may accept all revealed truths of God with a deep humility.
\medbreak

St. Raphael the Archangel, who happily guided Tobias and protected him against an evil spirit; be our companion on our earthly pilgrimage to the heavenly homeland.
\medbreak

Holy Angels, who unceasingly contemplate the face of the heavenly Father; teach us to see God and his perfection everywhere.
\medbreak

Holy Angels, who sang the glory of God and announced his peace to the people of good will at the birth of Jesus; bring the peace of Christ to our hearts.
\medbreak

Holy Angels, in whose company, Christ will come to judge the living and the dead; lead us on the way to salvation.
\medbreak

Holy Angels, who stand before the throne of God; present to him our prayers and deeds, that he may kindly accept them, forgive all our sins and grant us the grace to join you in heaven forever and ever. Amen.
\medbreak
\newpage
\newpageheader{Joined to the Guardian Angels}{Litany to the Guardian Angel}
\label{litanyguardianangelshtmllitany}

\section{Litany to the Guardian Angel}

\begin{litany}Lord, have mercy, \responce{Lord, have mercy.}\par 
\end{litany}
\begin{litany}Christ, have mercy, \responce{Christ, have mercy.}\par 
\end{litany}
\begin{litany}Lord, have mercy, \responce{Lord, have mercy.}\par 
\end{litany}
\begin{litany}Christ, hear us, \responce{Christ, graciously hear us.}\par 
\end{litany}
\begin{litany}God the Father of Heaven, \responce{Have mercy on us.}\par 
God the Son, Redeemer of the World, \par
God the Holy Spirit, \par
Holy Trinity, One God, \par
\end{litany}
\medbreak
\begin{litany}Holy Mary, Queen of Heaven, \responce{Pray for us.}\par 
\end{litany}
\begin{litany}Holy Angel, my Guardian, \responce{Pray for me.}\par 
Holy Angel, my Protector in all dangers, \par
Holy Angel, my Defence in all afflictions, \par
Holy Angel, my most faithful Lover, \par
Holy Angel, my Preceptor, \par
Holy Angel, my Guide, \par
Holy Angel, Witness of all my actions, \par
Holy Angel, my Helper in all my difficulties, \par
Holy Angel, my Negotiator with God, \par
Holy Angel, my Advocate, \par
Holy Angel, lover of Chastity, \par
Holy Angel, lover of Innocence, \par
Holy Angel, most obedient to God, \par
Holy Angel, Director of my soul, \par
Holy Angel, model of Purity, \par
Holy Angel, model of Docility, \par
Holy Angel, my Counsellor in doubt, \par
Holy Angel, my Guardian through life, \par
Holy Angel, my shield at the hour of Death, \par
\end{litany}
\medbreak
\noindent\begin{minipage}{\textwidth}
\begin{litany}Lamb of God, you take away the sins of the world, \responce{Spare us, O Lord.}\par 
\end{litany}
\begin{litany}Lamb of God, you take away the sins of the world, \responce{Graciously hear us, O Lord.}\par 
\end{litany}
\begin{litany}Lamb of God, you take away the sins of the world, \responce{Have mercy on us.}\par 
\end{litany}
\end{minipage}
\medbreak
\begin{samepage}
\prayer{L.}{\textbf{The Lord has put his Angels in charge of you,}}
\prayer{All}{To keep you in all your ways.}
\end{samepage}
\medbreak
\begin{samepage}
\prayer{L.}{\textbf{Let us pray:}\\ O God, Who with unspeakable providence vouchsafe to send Thy Angels to be our Guardians, mercifully grant, that we, Thy supplicants, may be always defended by their protection and enjoy their eternal society, through Jesus Christ, Thy Son, Our Lord, Who lives and reigns with you, in the unity of the Holy Spirit, One God, world without end.}
\prayer{All}{Amen.}
\end{samepage}
\medbreak
\newpage
\includeimage{saint-michael.pdf}\pagestyle{empty}
\hspace{0pt}
\vfill

\textit{``At that time Michael will stand up, the great prince, who
mounts guard over your people'' (Dn 12: 1)}
\medbreak
O God, who arranges the services of Angels and men in a
wonderful order; mercifully grant that our life may be protected on
earth by those who always do your service in heaven. Through our
Lord, Jesus Christ, your Son, who lives and reigns with you, in the
unity of the Holy Spirit, One God, for ever and ever. Amen.
\vfill
\textit{Icon of St. Michael the Archangel}
\hspace{0pt}
\newpage
\newpageheader{Joined to the Guardian Angels}{Saint Michael the Archangel}

\section{Litany to St. Michael the Archangel}

\begin{litany}Lord, have mercy, \responce{Lord, have mercy.}\par 
\end{litany}
\begin{litany}Christ, have mercy, \responce{Christ, have mercy.}\par 
\end{litany}
\begin{litany}Lord, have mercy, \responce{Lord, have mercy.}\par 
\end{litany}
\begin{litany}Christ, hear us, \responce{Christ, graciously hear us.}\par 
\end{litany}
\begin{litany}God the Father of Heaven, \responce{Have mercy on us.}\par 
God the Son, Redeemer of the World, \par
God the Holy Spirit, \par
Holy Trinity, One God, \par
\end{litany}
\medbreak
\begin{litany}Holy Mary, Queen of Angels, \responce{Pray for us.}\par 
St. Michael the Archangel, \par
Most glorious attendant of the Triune Divinity, \par
Standing at the right of the Altar of Incense, \par
Ambassador of Paradise, \par
Glorious Prince of the heavenly armies, \par
Leader of the angelic hosts, \par
Warrior who thrust Satan into Hell, \par
Defender against the wickedness and snares of the devil, \par
Standard-bearer of God's armies, \par
Defender of divine glory, \par
First defender of the Kingship of Christ, \par
Strength of God, \par
Invincible prince and warrior, \par
Angel of peace, \par
Guardian of the Christian Faith, \par
Champion of God's people, \par
Guardian angel of the Eucharist, \par
Defender of the Church, \par
Protector of the Sovereign Pontiff, \par
Angel of Catholic Action, \par
Powerful intercessor of Christians, \par
Bravest defender of those who hope in God, \par
Guardian of our souls and bodies, \par
Healer of the sick, \par
Help of those in their agony, \par
Consoler of the souls in Purgatory, \par
God's messenger for the souls of the just, \par
Terror of the evil spirits, \par
Victorious in battle against evil, \par
Guardian and Patron of the Universal Church \par
\end{litany}
\medbreak
\noindent\begin{minipage}{\textwidth}
\begin{litany}Lamb of God, you take away the sins of the world, \responce{Spare us, O Lord.}\par 
\end{litany}
\begin{litany}Lamb of God, you take away the sins of the world, \responce{Graciously hear us, O Lord.}\par 
\end{litany}
\begin{litany}Lamb of God, you take away the sins of the world, \responce{Have mercy on us.}\par 
\end{litany}
\end{minipage}
\medbreak
\begin{samepage}
\prayer{L.}{\textbf{Pray for us, O glorious St. Michael,}}
\prayer{All}{That we may be made worthy of the promises of Christ.}
\end{samepage}
\medbreak
\begin{samepage}
\prayer{L.}{\textbf{Let us pray:}\\ Relying, O Lord, upon the intercession of your blessed Archangel Michael, we humbly beg of you, that the Holy Eucharist in whose presence we kneel, may make our souls holy and pleasing to you. \\ Through Christ Our Lord.}
\prayer{All}{Amen.}
\end{samepage}
\medbreak

\subsection{Act of dedication to St. Michael the Archangel}

St. Michael the Archangel, the greatest Prince of heaven, the most faithful Guardian of the Church, behold, I am most unworthy to see your face, but full of confidence in your goodness, encouraged by the power of your prayers and numerous benefits, I stand before you in the company of my Guardian angel and in the presence of all Angels of heaven whom I take as witnesses of my devotion towards you.
\medbreak

I choose you this day for my special defender and protector. I firmly resolve to honour you always and to promote devotion to you. Be at my side always that I may never offend God, either by thought, word or deed. Defend me against the temptations of the devil, particularly directed against faith and purity and at the hour of my death, pray for me and lead me to the Eternal Home. Amen.
\medbreak

\subsection{Prayer to St. Michael the Archangel for the Church and the Holy Father}

St. Michael the Archangel, Patron of the Universal Church and a special Protector of the Holy See! We pray to you with lively faith and trust that you may defend and watch over the Holy Spouse of Christ, the Church. Show yourself now a great Guardian of the Church in this time of godlessness. Protect the Vicar of Christ on earth. Inspire the faithful with great honour due to the Vicar of Christ. May they respect and obey him with love as a Father of all the children of God. O Prince and Guardian of God's Holy People!
\medbreak

Pray that the Church may grow in numbers and in holiness. Inspire young people to become apostles and defenders of God's Holy Laws. Never allow hatred to triumph. With your victorious sword, defend the See of St. Peter. Come in haste to his defence in every need. Join together all the hearts of all the faithful with the Centre of the Catholic world. Amen.
\medbreak

\subsection{Prayer to St. Michael the Archangel}

Saint Michael Archangel, defend us in battle, be our safeguard against the wickedness and snares of the devil; may God rebuke him, we humbly pray; and do thou, O Prince of the heavenly host, by the power of God, cast into hell Satan and all the other evil spirits who wander through the world for the ruin of souls. Amen
\medbreak
\newpage
\includeimage{saint-paul.pdf}\pagestyle{empty}
\hspace{0pt}
\vfill

\textit{``By celebrating their anniversaries the Church proclaims the
achievement of the paschal mystery in the saints who have suffered
and have been glorified with Christ'' (S.L. 104).}
\medbreak
Almighty, eternal God, you permit us to venerate the merits of your
saints; we ask you to grant us your rich forgiveness through their
intercession. Through Christ our Lord. Amen.
\vfill
\textit{St. Paul the First Hermit by Karol Dank wart. Refectory of Jasna
Gora, second half of the XVII century}
\hspace{0pt}
\newpage
\newpageheader{Contemplation of God in His Saints}{Contemplation of God in His Saints}

\chapter{Contemplation of God in His Saints}

\instruction{The Second Vatican Council reminds us that in honouring the saints, the Church proclaims the paschal mystery. The Church presents their exemplary lives that steer us on our way to the Eternal Father and allow us to participate in their abundant merits (cf. Sacrosanctum Concilium 104). In our Constitutions we read the following: ``We are led to do this through the example of St. Paul the First Hermit in his communing one-on-one with God \ldots ~We are further encouraged by the tradition of our Fathers, but above all, through the example of Christ and his Mother praying. It is in prayer that we are able to recognize the plan of God for us and for the world'' (art. 26).}
\medbreak

\instruction{In our Order, we honour particularly St. Joseph, St. Paul the First Hermit, St. Augustine, St. Casimir and St. Hedwig. We follow the virtues of our Founder, Blessed Eusebius. Their examples of holiness inspire us in our yearning for Christian perfection.}
\medbreak
\newpage
\newpageheader{Contemplation of God in His Saints}{Saint Joseph}

\section{Saint Joseph}

\instruction{St. Joseph was a righteous man to whom divine Providence has entrusted the privilege to participate decisively in the greatest mysteries of salvation. It was he who obeyed the command of the Angel and took Mary into his house as his wife. He believed that she had truly conceived the Divine Child through the power of the Holy Spirit (cf. Mt 1:19-25). When mortal danger for the Child came from King Herod, ``Joseph got up and taking the Child and his Mother with him, left that night for Egypt'' (Mt 2:14). When the danger passed, he ``returned to the land of Israel, and settled in the town of Nazareth in Galilee. In this way the words spoken through the prophets were to be fulfilled `He will be called a Nazarene' ''(Mt 2:23). When Jesus began his public mission, St. Joseph appeared at the Jordan River, where he was baptized by St. John. Jesus was thought to be the son of Joseph (cf. Lk 3:23). In the Church's liturgy and in the piety of the faithful, St. Joseph stands as a man of justice, a most faithful servant of the living God, Spouse of the Most Holy Virgin Mary and a silent witness of divine mysteries. ``He stands on the threshold of the Old and New Testaments. In his person, the dignity of Patriarchs and the oracles of the Prophets, find their fulfilment'' (St. Bernardino of Siena). The cult to St. Joseph in the Pauline tradition does not exceed specifically the general veneration due to St. Joseph and practiced in the Catholic Church.}
\medbreak

\label{saintjosephhtmllitany}

\subsection{Litany of St. Joseph}

\begin{litany}Lord, have mercy, \responce{Lord, have mercy.}\par 
\end{litany}
\begin{litany}Christ, have mercy, \responce{Christ, have mercy.}\par 
\end{litany}
\begin{litany}Lord, have mercy, \responce{Lord, have mercy.}\par 
\end{litany}
\begin{litany}Christ, hear us, \responce{Christ, graciously hear us.}\par 
\end{litany}
\begin{litany}God the Father of Heaven, \responce{Have mercy on us.}\par 
God the Son, Redeemer of the World, \par
God the Holy Spirit, \par
Holy Trinity, One God, \par
\end{litany}
\medbreak
\begin{litany}Holy Mary, \responce{Pray for us.}\par 
St. Joseph, \par
Renowned offspring of David, \par
Light of Patriarchs, \par
Spouse of the Mother of God, \par
Chaste guardian of the Virgin, \par
Foster-father of the Son of God, \par
Diligent protector of Christ, \par
Head of the Holy Family, \par
Joseph most just, \par
Joseph most chaste, \par
Joseph most prudent, \par
Joseph most strong, \par
Joseph most obedient, \par
Joseph most faithful, \par
Mirror of patience, \par
Lover of poverty, \par
Model of artisans, \par
Glory of home life, \par
Guardian of virgins, \par
Pillar of families, \par
Solace of the wretched, \par
Hope of the sick, \par
Patron of the dying, \par
Terror of demons, \par
Protector of Holy Church, \par
\end{litany}
\medbreak
\noindent\begin{minipage}{\textwidth}
\begin{litany}Lamb of God, you take away the sins of the world, \responce{Spare us, O Lord.}\par 
\end{litany}
\begin{litany}Lamb of God, you take away the sins of the world, \responce{Graciously hear us, O Lord.}\par 
\end{litany}
\begin{litany}Lamb of God, you take away the sins of the world, \responce{Have mercy on us.}\par 
\end{litany}
\end{minipage}
\medbreak
\begin{samepage}
\prayer{L.}{\textbf{He made him the Lord of his house,}}
\prayer{All}{And prince over all his possessions.}
\end{samepage}
\medbreak
\begin{samepage}
\prayer{L.}{\textbf{Let us pray:}\\ O God, in your ineffable providence you were pleased to choose Blessed Joseph to be the spouse of your most holy Mother; grant, we beg you, that we may be worthy to have him as our intercessor in heaven whom on earth we venerate as our Protector. You who live and reign forever and ever.}
\prayer{All}{Amen.}
\end{samepage}
\medbreak

\subsection[Prayer to St. Joseph]{Prayer to St. Joseph \\ {\small\normalfont (recited after the Rosary)}}

O St. Joseph, we flee to your protection in all our necessities. Having invoked the help of your Holy Spouse Mary, we also ask for your aid. By the love that joined you to the Mother of God ever Virgin and by the fatherly care, with which you caressed the Child Jesus, we beg you humbly to look down upon the heritage obtained by the precious Blood of Jesus Christ and through your powerful intercession assist us in our needs.
\medbreak

O providential Guardian of the Holy Family, watch constantly over the posterity of Jesus Christ. Remove from us, our beloved Father, all error and corruption. Come to our rescue in our struggles with the powers of darkness. As once you saved the Child Jesus from the danger that imperilled his life, so today defend the Church of God from all adversities and snares of the devil. Surround us with your constant protection, that following your example, and strengthened with your help, we may live piously, die happily and obtain everlasting happiness in heaven. Amen.
\medbreak
\newpage
\includeimage{saint-paul-crown.pdf}\pagestyle{empty}
\hspace{0pt}
\vfill

\textit{``Along with our Marian devotion, we should particularly cultivate the devotion [\ldots] to St. Paul the First Hermit.'' (art. 32). }
\medbreak
O God, you have called St. Paul the First Hermit our 
Father into the desert to contemplate your name in solitude. Through
his intercession, may we reject all earthly desires and to seek your
love alone. Through Christ Our Lord. Amen.
\vfill
%\textit{St. Paul the First Hermit. Painting workshop of Jasna Gora, art collections, XVII century.}
\hspace{0pt}
\newpage
\newpageheader{Contemplation of God in His Saints}{Saint Paul the First Hermit}

\section{Saint Paul the First Hermit}

\instruction{In the light of the recent studies, St. Paul the First Hermit lived in the years between 228 -341. His personality and life style harmonize perfectly with the events and situation of the Church of the first centuries. The exact dates are not of great importance but the example, as an ideal model of the hermit and his personal union with God, remains always a great and worthy attraction for Christians to follow.}
\medbreak

\instruction{St. Paul the First Hermit is well known among the early Fathers of the Church and in the lives of the first anchorites. Down the centuries in Pauline Communities the title ``St. Paul the First Hermit'' was generally accepted and approved by the Holy See as an official name of the Order.}
\medbreak

\instruction{It was St. Jerome who wrote the life of St. Paul the First Hermit. According to his relations St. Paul was quite an educated man in Greek and Egyptian culture. He was endowed with a noble and gentle nature towards the people. During the persecution of the Church, inspired by the Holy Spirit, he left the family home and settled himself in the desert, where ``alone with God alone'', he spent all his earthy life in prayer and contemplation of God's mysteries.}
\medbreak

\instruction{There is no doubt that St. Paul the First Hermit was a historical figure. However, many extraordinary stories, related by St. Jerome (345-420), are not necessarily true or essential. For all Pauline generations St. Paul the First Hermit is not only the most venerable and beloved Patriarch of the Order, but above all a living symbol of imitation of Christ in the desert that originated in the past, but still flourishes today and will continue to inspire his spiritual sons for generations to come.}
\medbreak

\label{saintpaulhtmllitany}

\subsection{Litany of St. Paul the First Hermit}

\begin{litany}Lord, have mercy, \responce{Lord, have mercy.}\par 
\end{litany}
\begin{litany}Christ, have mercy, \responce{Christ, have mercy.}\par 
\end{litany}
\begin{litany}Lord, have mercy, \responce{Lord, have mercy.}\par 
\end{litany}
\begin{litany}Christ, hear us, \responce{Christ, graciously hear us.}\par 
\end{litany}
\begin{litany}God the Father of Heaven, \responce{Have mercy on us.}\par 
God the Son, Redeemer of the World, \par
God the Holy Spirit, \par
Holy Trinity, One God, \par
\end{litany}
\medbreak
\begin{litany}Holy Mary, Mother of God, \responce{Pray for us.}\par 
Holy Mary, Mother of our Order, \par
Holy Mary, hope for a better world, \par
Holy Mary, unfailing help in our striving for holiness, \par
Holy Guardian Angels, who watches over our consciences, \par
St. Paul, father of our Pauline Order, \par
St. Paul, lover of solitude, \par
St. Paul, an example of Christian perfection, \par
St. Paul, mirror of poverty, \par
St. Paul, true model of humility, \par
St. Paul, ornament of purity, \par
St. Paul, faithful servant of the Eternal King, \par
St. Paul, refuge of the tempted, \par
St. Paul, inflamed with divine love, \par
St. Paul, totally dedicated to God, \par
St. Paul, example of sobriety and abstinence, \par
St. Paul, wholly trusting in God's Providence, \par
St. Paul, summit of patience, \par
St. Paul, man without complaint, \par
St. Paul, admired by St. Anthony, \par
St. Paul, commending your spirit with hands raised to heaven, \par
St. Paul, honoured after death with the service of lions, \par
St. Paul, buried by St. Anthony, \par
St. Paul, taken to heaven by the Angels, \par
\end{litany}
\medbreak
\begin{litany}That we may live in a spirit of faith, hope and love, \responce{Intercede to God for us.}\par 
That we may abound with good deeds in our daily life, \par
That we may faithfully fulfil the duties of our state of life, \par
That we may use the goods of this world with prudence, \par
That we may unceasingly remember our death and God's judgment, \par
That we may follow the Crucified Saviour and be always faithful to his Sorrowful Mother, \par
\end{litany}
\medbreak
\noindent\begin{minipage}{\textwidth}
\begin{litany}Lamb of God, you take away the sins of the world, \responce{Spare us, O Lord.}\par 
\end{litany}
\begin{litany}Lamb of God, you take away the sins of the world, \responce{Graciously hear us, O Lord.}\par 
\end{litany}
\begin{litany}Lamb of God, you take away the sins of the world, \responce{Have mercy on us.}\par 
\end{litany}
\end{minipage}
\medbreak
\begin{samepage}
\prayer{L.}{\textbf{Pray for us Saint Paul the First Hermit.}}
\prayer{All}{That we may be made worthy of the promises of Christ.}
\end{samepage}
\medbreak
\begin{samepage}
\prayer{L.}{\textbf{Let us pray:}\\ God, you granted that Saint Paul, our Father, achieve sanctity in the desert, grant through his intercession, that we who are striving to develop in ourselves a spirit of prayer and service may come ever closer to you in love. Through Christ our Lord.}
\prayer{All}{Amen.}
\end{samepage}
\medbreak

\instruction{or:}
\medbreak
\begin{samepage}
\prayer{L.}{\textbf{Let us pray:}\\ O God, who led St. Paul our Father into the desert to contemplate the mystery of your Holy Name and arrive safely at his eternal dwellings; teach us through his intercession that we may despise worldly pleasures and seek everlasting happiness. Through Christ our Lord.}
\prayer{All}{Amen.}
\end{samepage}
\medbreak
\newpage
\includeimage{saint-paul-chapel.pdf}\pagestyle{empty}
\hspace{0pt}
\vfill

\textit{``But because of Christ, I have come to consider all these
advantages that I had as disadvantages. Not only that, but I believe
nothing can happen that will outweigh the supreme advantage of
knowing Christ Jesus my Lord'' (Phil 3:7-8).}
\medbreak
\noindent St. Paul, who now enjoys the bliss of Heaven,\\
kindly look down on your sons here below.\\
Draw them to yourself,\\
by the power of your prayer.
\vfill
%\textit{The chapel of St. Paul the Hermit, a work of Franciszek Zaor; Basilica of Jasna Gora, XVII century.}
\hspace{0pt}
\newpage
\newpageheader{Contemplation of God in His Saints}{Novena to St. Paul the First Hermit}

\subsection{Novena to St. Paul the First Hermit}

\instruction{Traditionally, we begin the Novena 9 days before the external solemnity of St. Paul the First Hermit. It is before the next Sunday after the 15th of January. The novena to St. Paul the First Hermit enjoys a special splendour. It should take place in the chapel or at another place dedicated to St. Paul the First Hermit. All members of the community take part in the Novena. The solemnly professed Monks wear the processional mantle. After the solemn entrance to the place of devotion and the placing of the relics of our Patriarch, one of the hymns honouring St. Paul the First Hermit is sung. The faithful are encouraged to join in the singing, alternately. It is favourable to sing the hymns with musical accompaniment, either organ or trumpet orchestra. A short homily should be delivered at the end of the devotion. The veneration of the relics takes place, followed by the blessing with the reliquary given to the participants according to the Order's ritual.}
\medbreak

\label{saintpaulnovenahtmlhymn}

\subsubsection*{Hymn I}
\addcontentsline{toc}{subsubsection}{Hymn I}

\rubric{Fathers and Brothers} 

\begin{center}
\noindent\includegraphics[width=\textwidth,height=\textheight,keepaspectratio]{modern/lilypond/d8/lily-0eeeee21-1}
\noindent\includegraphics[width=\textwidth,height=\textheight,keepaspectratio]{modern/lilypond/d8/lily-0eeeee21-2}
\noindent\includegraphics[width=\textwidth,height=\textheight,keepaspectratio]{modern/lilypond/d8/lily-0eeeee21-3}
\noindent\includegraphics[width=\textwidth,height=\textheight,keepaspectratio]{modern/lilypond/d8/lily-0eeeee21-4}
\noindent\includegraphics[width=\textwidth,height=\textheight,keepaspectratio]{modern/lilypond/d8/lily-0eeeee21-5}
\noindent\includegraphics[width=\textwidth,height=\textheight,keepaspectratio]{modern/lilypond/d8/lily-0eeeee21-6}
\end{center}
\medbreak
\begin{minipage}{\textwidth}\begin{addmargin}[1.5in]{0em}

\rubric{Congregation} 
\begin{psalm}It was then that Paul the Hermit\par
had deprived himself of all,\par
of his family, pompous honour,\par
and prosperity mortal.\par
\medbreak
Not delaying rather fleeing,\par
to the den of the desert,\par
while his sister had been calling,\par
to him that he should revert.\par
\medbreak
\end{psalm}
\medbreak
\end{addmargin}\end{minipage}
 \par\medbreak
\begin{minipage}{\textwidth}

\rubric{Fathers and Brothers} 
\begin{psalm}``Wherever you go or thus move,\par
Oh! my brother, I'm moaning,\par
Ah! your sister so sincerely,\par
as with tears be convincing.\par
\medbreak
Don't leave your paternal home, Paul,\par
to be thought inhumanely,\par
but I beg you to remain here,\par
and wish you to protect me''.\par
\medbreak
\end{psalm}
\medbreak
\end{minipage}
 \par\medbreak
\begin{minipage}{\textwidth}\begin{addmargin}[1.5in]{0em}

\rubric{Congregation} 
\begin{psalm}How your body, Ah, so fragile,\par
doesn't belong in the caves,\par
nor your talent of great measure,\par
hidden in those mountain graves.\par
\medbreak
For the people of the Thebaid,\par
desire to acknowledge you,\par
so horrendous would this loss be,\par
incredible, if it's true.\par
\medbreak
\end{psalm}
\medbreak
\end{addmargin}\end{minipage}
 \par\medbreak
\begin{minipage}{\textwidth}

\rubric{Fathers and Brothers} 
\begin{psalm}So behold our lover of caves,\par
who rejected tears and pleas,\par
he was the great self triumpher,\par
over sec'lar pageantries.\par
\medbreak
As he went on with celerity,\par
to the depth of the forest,\par
and no knowledge of returning,\par
did he wish in this world's rest.\par
\medbreak
\end{psalm}
\medbreak
\end{minipage}
 \par\medbreak
\begin{minipage}{\textwidth}\begin{addmargin}[1.5in]{0em}

\rubric{Congregation} 
\begin{psalm}Look! What's so admirable,\par
his eremitical face,\par
solely walking on this new path,\par
never trodden with such grace.\par
\medbreak
Not alone in this nude body,\par
covered up with some palm leaves,\par
for his pleasure fruit was given,\par
from Him in Whom he believed.\par
\medbreak
\end{psalm}
\medbreak
\end{addmargin}\end{minipage}
 \par\medbreak
\begin{minipage}{\textwidth}

\rubric{Fathers and Brothers} 
\begin{psalm}Now this hermit who was thirsty,\par
thus drank from the flowing font,\par
and so save his life immortal,\par
in world's cares he gave no want.\par
\medbreak
All his body how exhausted,\par
on the ground he formed his bed,\par
oh, this practice sternly followed,\par
he endured and with stance led.\par
\medbreak
\end{psalm}
\medbreak
\end{minipage}
 \par\medbreak
\begin{minipage}{\textwidth}\begin{addmargin}[1.5in]{0em}

\rubric{Congregation} 
\begin{psalm}For some sixty years in counting,\par
God's black raven brought repast,\par
one half loaf as daily measure,\par
bread for keening man's own fast.\par
\medbreak
At the coming of Anthony,\par
God so doubled the portion,\par
God thus merciful and clement,\par
that he cared for their ration.\par
\medbreak
\end{psalm}
\medbreak
\end{addmargin}\end{minipage}
 \par\medbreak
\begin{minipage}{\textwidth}

\rubric{Fathers and Brothers} 
\begin{psalm}Holy Paul said to Anthony,\par
``behold him, whom you have sought,\par
you look on a man of great age,\par
whose repose will soon be brought.\par
\medbreak
Now I pray you go and bring me\par
the cloak of Athanasius,\par
so to wrap around my body\par
and place in the earth precious.\par
\medbreak
\end{psalm}
\medbreak
\end{minipage}
 \par\medbreak
\begin{minipage}{\textwidth}\begin{addmargin}[1.5in]{0em}

\rubric{Congregation} 
\begin{psalm}With the thought of Paul's transition,\par
the old man rushed off distraught,\par
then returning all elated,\par
he saw Paul's soul upward brought.\par
\medbreak
High among the choirs of heaven,\par
this he pondered and conceived,\par
Anthony then asked the question,\par
``should I rejoice or thus grieve?''\par
\medbreak
\end{psalm}
\medbreak
\end{addmargin}\end{minipage}
 \par\medbreak
\begin{minipage}{\textwidth}

\rubric{Fathers and Brothers} 
\begin{psalm}What grief Anthony encountered,\par
from the death of our hermit,\par
yet the glory of Paul honoured\par
rather placed joy into it.\par
\medbreak
All the clothes he took with reverence,\par
just before the interment,\par
as he chanted with such fervour,\par
for the soul which to God went.\par
\medbreak
\end{psalm}
\medbreak
\end{minipage}
 \par\medbreak
\begin{minipage}{\textwidth}\begin{addmargin}[1.5in]{0em}

\rubric{Congregation} 
\begin{psalm}Oh! how famed was that entombment,\par
that Saint Paul had merited,\par
never had the world admired,\par
such an act, experience.\par
\medbreak
How two lions did dig the grave,\par
for the corpse of the hermit,\par
and by which they did reveal him,\par
on holiness before men.\par
\medbreak
\end{psalm}
\medbreak
\end{addmargin}\end{minipage}
 \par\medbreak
\begin{minipage}{\textwidth}

\rubric{Fathers and Brothers} 
\begin{psalm}Saint Paul, now we pray you look here,\par
on us from your place on high,\par
never let us be divided,\par
from your own paternal eye.\par
\medbreak
Do remember your dear brothers,\par
and desert them not on earth,\par
but so always interceding,\par
bring us to our King, our worth.\par
\medbreak
\end{psalm}
\medbreak
\end{minipage}
 \par\medbreak
 \newpage 
\label{saintpaulnovenahtmlhymn}

\subsubsection*{Hymn II}
\addcontentsline{toc}{subsubsection}{Hymn II}

\rubric{Fathers and Brothers} 

\begin{center}
\noindent\includegraphics[width=\textwidth,height=\textheight,keepaspectratio]{modern/lilypond/7d/lily-b6a5e964-1}
\noindent\includegraphics[width=\textwidth,height=\textheight,keepaspectratio]{modern/lilypond/7d/lily-b6a5e964-2}
\noindent\includegraphics[width=\textwidth,height=\textheight,keepaspectratio]{modern/lilypond/7d/lily-b6a5e964-3}
\end{center}
\medbreak
\begin{minipage}{\textwidth}\begin{addmargin}[1.5in]{0em}

\rubric{Congregation} 
\begin{psalm}This great solemnity,\par
for our Father's glory,\par
such a grand day allows,\par
to show it by our vows,\par
so joyously.\par
\medbreak
\end{psalm}
\medbreak
\end{addmargin}\end{minipage}
 \par\medbreak
\begin{minipage}{\textwidth}

\rubric{Fathers and Brothers} 
\begin{psalm}Let us now reconfirm,\par
how brave was his stature,\par
from the world retreating,\par
to heaven entreating,\par
much happier.\par
\medbreak
\end{psalm}
\medbreak
\end{minipage}
 \par\medbreak
\begin{minipage}{\textwidth}\begin{addmargin}[1.5in]{0em}

\rubric{Congregation} 
\begin{psalm}To leave this world so vain,\par
to seek the Saviour's reign,\par
he taught by his own call,\par
and admonished us all,\par
imitate me.\par
\medbreak
\end{psalm}
\medbreak
\end{addmargin}\end{minipage}
 \par\medbreak
\begin{minipage}{\textwidth}

\rubric{Fathers and Brothers} 
\begin{psalm}Alone with God alone,\par
all proud honours unknown,\par
he fled his heritage,\par
hid in his hermitage,\par
desiring more.\par
\medbreak
\end{psalm}
\medbreak
\end{minipage}
 \par\medbreak
\begin{minipage}{\textwidth}\begin{addmargin}[1.5in]{0em}

\rubric{Congregation} 
\begin{psalm}This is his hope for you,\par
withdrawn from sec'lar view,\par
as he sought out a den,\par
which for some time had been,\par
a place most true.\par
\medbreak
\end{psalm}
\medbreak
\end{addmargin}\end{minipage}
 \par\medbreak
\begin{minipage}{\textwidth}

\rubric{Fathers and Brothers} 
\begin{psalm}From palm branches so fine,\par
he wove clothes as a sign to,\par
keep his custody,\par
and tame his frail body,\par
through suffering.\par
\medbreak
\end{psalm}
\medbreak
\end{minipage}
 \par\medbreak
\begin{minipage}{\textwidth}\begin{addmargin}[1.5in]{0em}

\rubric{Congregation} 
\begin{psalm}Father most amiable,\par
what food! How terrible,\par
by which you kept your zeal,\par
yet no great banquet meal,\par
only palm dates.\par
\medbreak
\end{psalm}
\medbreak
\end{addmargin}\end{minipage}
 \par\medbreak
\begin{minipage}{\textwidth}

\rubric{Fathers and Brothers} 
\begin{psalm}Then once he took his meal,\par
the stream near did appeal,\par
the water was ice-cold,\par
refreshing him now old,\par
giving new life.\par
\medbreak
\end{psalm}
\medbreak
\end{minipage}
 \par\medbreak
\begin{minipage}{\textwidth}\begin{addmargin}[1.5in]{0em}

\rubric{Congregation} 
\begin{psalm}In order to give man,\par
the strength for a life's span,\par
the raven brought repast,\par
daily bread to holdfast,\par
at noontime came.\par
\medbreak
\end{psalm}
\medbreak
\end{addmargin}\end{minipage}
 \par\medbreak
\begin{minipage}{\textwidth}

\rubric{Fathers and Brothers} 
\begin{psalm}When Anthony had come,\par
once more portion in sum,\par
God doubled the measure,\par
granting his saints pleasure,\par
by the blest grain.\par
\medbreak
\end{psalm}
\medbreak
\end{minipage}
 \par\medbreak
\begin{minipage}{\textwidth}\begin{addmargin}[1.5in]{0em}

\rubric{Congregation} 
\begin{psalm}Soon Paul's transition came,\par
he reached heavenly fame,\par
wrapped in the mantle,\par
thus from Athanasius,\par
as was witnessed.\par
\medbreak
\end{psalm}
\medbreak
\end{addmargin}\end{minipage}
 \par\medbreak
\begin{minipage}{\textwidth}

\rubric{Fathers and Brothers} 
\begin{psalm}No tools there could be found,\par
so lions dug the ground,\par
their work was diligent,\par
nonetheless elegant,\par
they formed his grave.\par
\medbreak
\end{psalm}
\medbreak
\end{minipage}
 \par\medbreak
\begin{minipage}{\textwidth}\begin{addmargin}[1.5in]{0em}

\rubric{Congregation} 
\begin{psalm}What a rare memory,\par
of our Father's glory,\par
Anthony went away,\par
with the tunic that day,\par
the robe of palm.\par
\medbreak
\end{psalm}
\medbreak
\end{addmargin}\end{minipage}
 \par\medbreak
\begin{minipage}{\textwidth}

\rubric{Fathers and Brothers} 
\begin{psalm}This triumphant symbol,\par
in a world deceitful,\par
he overcame the flesh,\par
and struck down Satan's mesh,\par
through fervent prayer.\par
\medbreak
\end{psalm}
\medbreak
\end{minipage}
 \par\medbreak
\begin{minipage}{\textwidth}\begin{addmargin}[1.5in]{0em}

\rubric{Congregation} 
\begin{psalm}Here's our inheritance,\par
palm clothes for assurance,\par
which he gave all of us,\par
to be victorious,\par
over our foes.\par
\medbreak
\end{psalm}
\medbreak
\end{addmargin}\end{minipage}
 \par\medbreak
\begin{minipage}{\textwidth}

\rubric{Fathers and Brothers} 
\begin{psalm}So Brothers fight bravely,\par
with great virility,\par
your sins you can conquer,\par
and see God forever,\par
in his Kingdom.\par
\medbreak
\end{psalm}
\medbreak
\end{minipage}
 \par\medbreak
\begin{samepage}
\prayer{L.}{\textbf{Pray for us Saint Paul the First Hermit.}}
\prayer{All}{That we may be made worthy of the promises of Christ.}
\end{samepage}
\medbreak
\begin{samepage}
\prayer{L.}{\textbf{Let us pray:}\\ God, you granted that Saint Paul, our Father, achieve sanctity in the desert, grant through his intercession, that we who are striving to develop in ourselves a spirit of prayer and service may come ever closer to you in love. Through Christ our Lord.}
\prayer{All}{Amen.}
\end{samepage}
\medbreak

\instruction{or:}
\medbreak
\begin{samepage}
\prayer{L.}{\textbf{Let us pray:}\\ O God, who led St. Paul our Father into the desert to contemplate the mystery of your Holy Name and arrive safely at his eternal dwellings; teach us through his intercession that we may despise worldly pleasures and seek everlasting happiness. Through Christ our Lord.}
\prayer{All}{Amen.}
\end{samepage}
\medbreak
\newpage
\newpageheader{Contemplation of God in His Saints}{Hymns in honour of St. Paul the First Hermit}

\subsection{Hymns in honour of St. Paul the First Hermit}

\subsubsection*{I. Let us sing to our Patriarch}
\addcontentsline{toc}{subsubsection}{I. Let us sing to our Patriarch}

\begin{psalm}Let us sing to our Patriarch,\par
let us pay homage to our Father,\par
with pure heart and humbled mind,\par
united in unison, together.\par
\medbreak
Let us direct our steps toward him,\par
let us draw our strength from him,\par
may the example of his virtue,\par
shine forth in our soul's darkness.\par
\medbreak
May we honour Mary,\par
In humility, charity and silence,\par
may we stand at Jesus' side,\par
and share his triumphant crown.\par
\medbreak
O St. Paul who are already in Heaven,\par
today, we humbly pray,\par
accept our homage, a sign of our love,\par
may we live with you eternally.\par
\medbreak
\end{psalm}
\medbreak

\subsubsection*{II. O come my thirsting soul for peace}
\addcontentsline{toc}{subsubsection}{II. O come my thirsting soul for peace}

\begin{psalm}O come my thirsting soul for peace,\par
come under the foliage palm.\par
Rest for a while in its shade,\par
after the heated toil, put your head,\par
on the stone that closes the cave.\par
\medbreak
A gentle breeze of the desert shall cool you,\par
and the night's silence will close your eyes,\par
and God's hand, a mistress of miracles,\par
will never depart from you.\par
\medbreak
When you get new strength to do your work,\par
St. Paul will hold you by your hand,\par
and will show you new desert paths,\par
and will remove from you all your fears.\par
\medbreak
Life will lead to the only homeland above,\par
where there will be sadness or heart's longing no more,\par
now he leads them, his white-dressed sons,\par
on the thorny, virtuous way.\par
\medbreak
\end{psalm}
\medbreak

\instruction{For the Novices}
\medbreak
\begin{psalm}May you never loose sight of your goal,\par
watch and pray through your Novitiate time,\par
many have lost the grace of vocation,\par
you may share their fate as well, Dear Brother.\par
\medbreak
Therefore, watch for the Lord in solitude,\par
follow the example of St. Paul in a cave,\par
you will find your happiness from the early morning,\par
like your Patron in the desert.\par
\medbreak
\end{psalm}
\medbreak

\subsubsection*{III. Our Patriarch}
\addcontentsline{toc}{subsubsection}{III. Our Patriarch}

\begin{psalm}Our Patriarch, you are an example for all hermits,\par
you are the Father of the Pauline Order,\par
we honour you in humility,\par
and ask you for your help before God's Throne.\par
\medbreak
You have fulfilled all God's commandments,\par
you have despised all earthly pleasures quite early,\par
you have overcome all Satan's snares,\par
and have spent many years in the desert.\par
\medbreak
God himself made you strong,\par
a raven brought you daily bread,\par
wild beasts willingly served you,\par
for you lived not for this earth, but for heaven.\par
\medbreak
May your example encourage us,\par
to serve God in your Order, our Father,\par
may our vocation grow stronger,\par
that we may join you, after we depart\par
\medbreak
\end{psalm}
\medbreak

\subsubsection*{Novitiate Doors}
\addcontentsline{toc}{subsubsection}{Novitiate Doors}

\begin{psalm}How will we find God, in these blessed days?\par
Signs are posted on the Novitiate doors.\par
\medbreak
The Crucifix, a sign of Passion,\par
cries out: Love your Cross,\par
for the Cross comes from God's hand,\par
and leads you ever high.\par
\medbreak
The Image of Czestochowa's Lady,\par
gently whispers to you:\par
Trust your Heavenly Mother,\par
She is your Patroness forever.\par
\medbreak
The little sanctuary light,\par
shines brightly in the night\par
and reminds you:\par
these are the days filled with God's grace.\par
\medbreak
The black bird with bread in its beak,\par
this symbolic raven,\par
speaks to you:\par
Your real Bread is your Saviour, God himself.\par
\medbreak
Therefore, dear brother,\par
whenever you fix your eyes on these doors,\par
may your heart be inflamed,\par
and may you make a new step in love.\par
\medbreak
\end{psalm}
\medbreak
 \newpage 

\subsubsection*{Antiphone to St. Paul the First Hermit \\ {\small\normalfont Words and music: Fr Irenaeus Wybraniak OSPPE}}

\addcontentsline{toc}{subsubsection}{Antiphone to St. Paul the First Hermit Words and music: Fr Irenaeus Wybraniak OSPPE}

\begin{center}
\noindent\includegraphics[width=\textwidth,height=\textheight,keepaspectratio]{modern/lilypond/3b/lily-977e9a36-1}
\noindent\includegraphics[width=\textwidth,height=\textheight,keepaspectratio]{modern/lilypond/3b/lily-977e9a36-2}
\noindent\includegraphics[width=\textwidth,height=\textheight,keepaspectratio]{modern/lilypond/3b/lily-977e9a36-3}
\noindent\includegraphics[width=\textwidth,height=\textheight,keepaspectratio]{modern/lilypond/3b/lily-977e9a36-4}
\end{center}

\subsubsection*{Hymnus de sancto Paulo primo eremita \\ {\small\normalfont (Cantuale Paulinorum, ca 1526, Jasna Gora Archive ms. R 583, p. 230-234)}}

\addcontentsline{toc}{subsubsection}{Hymnus de sancto Paulo primo eremita
    (Cantuale Paulinorum, ca 1526, Jasna Gora Archive
    ms. R 583, p. 230-234)}

\gresetlastline{justified}
\gresetinitiallines{0}
\gresethyphen{auto}
\gregorioscore{music/ordo-sanctorum}\medbreak
\begin{psalm}2. Pauper mitis et docilis,\par
Celebs, pius, oboediens,\par
Modestus, sanctus, humilis\par
Iustus, verus et sapiens.\par
\medbreak
3. Quaesivit hunc Anthonius,\par
Quem Christus ei nuncius\par
Panem caelestem nuncius\par
Ieiuno corvus attulit.\par
\medbreak
4. Rosa fulgens et inclita\par
Baptista Christi praemio,\par
Paulus Eliae consita\par
Consimilis in invio.\par
\medbreak
5. Suum predixit exitum\par
Propheta pandens abdita,\par
Leo deplorat obitum\par
Flectens genua subdita.\par
\medbreak
6. Te transitu mirabilem\par
Sanctorum in splendoribus\par
Te acelis honorabilem\par
Translatum caeli civibus.\par
\medbreak
7. Veste  miranda praedicat\par
Anthonius, Hieronimus,\par
Laude multa magnificat\par
Te pater pie petimus.\par
\medbreak
8. Christi nos auge munere\par
Auctos fove precantibus\par
Pacis astringe foedere\par
Redde polorum sedibus.\par
\medbreak
9. Honor virtutis inclite,\par
Patri prolique genite,\par
Sacro spirato flamini\par
Sit sine meta termini.\par
\medbreak
\end{psalm}
\medbreak
\newpage
\includeimage{blessed-eusebius.pdf}\pagestyle{empty}
\hspace{0pt}
\vfill

\textit{``Anyone who does the will of my Father in heaven, he is
my brother and sister and mother'' (Mt 12:50).}
\vspace{1em}
O God, in Blessed Eusebius, you have shown us an exam-
ple of the Evangelical perfection; grant that passing through the
uncertain world events, we may love the things that last forever.
Through Christ our Lord. Amen.
\vfill
\textit{Blessed Eusebius, Paint shop of Jasna Gora.
Art collection of Jasna Gora, 1966}
\hspace{0pt}
\newpage
\newpageheader{Contemplation of God in His Saints}{Blessed Eusebius}

\section{Blessed Eusebius}

\instruction{Blessed Eusebius is considered a co-founder of the Order of the Brothers of St. Paul the First Hermit. From his initiatives the Pauline Order began to exist in Hungary in 1250. It was created from the two joined monasteries of hermits living in the Danube River forests: the Monastery of St. Jacob, situated upon the hill of Patacs and the Monastery of the Holy Cross. Blessed Eusebius, a Canon of Esztergom (lat. Strigonium) during the invasion of Turks in 1241, running for his life, found safety in the forests of Pecs, where he encountered many saintly men living individually in seclusion. In 1246 he gathered them together in the Monastery of the Holy Cross. By doing this, he laid the foundation for the religious community that lived under the rule of Bishop Bartholomew. Practically, he was the first provincial of the new Order which since 1250 has accepted a title: ``Order of the Brothers of St. Paul the First Hermit.'' In 1263 at the first General Chapter this title was approved, and in 1308 sanctioned by Pope Clement V and conveyed to them by Cardinal Gentilis. Because of his efforts to establish and gain approval from the Holy See of the new religious community and because of his saintly life, Eusebius, has enjoyed great popularity among the people since his death in 1270. They venerated him like a canonized Saint. His labours produced many blessed fruits of holiness in the many members of the Order he founded. In popular devotion he is invoked as ``Blessed.''}
\medbreak

\subsection{Litany of Blessed Eusebius}

\begin{litany}Lord, have mercy, \responce{Lord, have mercy.}\par 
\end{litany}
\begin{litany}Christ, have mercy, \responce{Christ, have mercy.}\par 
\end{litany}
\begin{litany}Lord, have mercy, \responce{Lord, have mercy.}\par 
\end{litany}
\begin{litany}Christ, hear us, \responce{Christ, graciously hear us.}\par 
\end{litany}
\begin{litany}God the Father of Heaven, \responce{Have mercy on us.}\par 
God the Son, Redeemer of the World, \par
God the Holy Spirit, \par
Holy Trinity, One God, \par
\end{litany}
\medbreak
\begin{litany}Holy Mary, Mother of God, \responce{Pray for us.}\par 
Holy Mary, Queen of Hermits, \par
Father Eusebius, glory of the Pauline Order, \par
Father Eusebius, example of piety from the days of your youth, \par
Father Eusebius, lover of silence and solitude, \par
Father Eusebius, adoring the Most Blessed Sacrament in isolation, \par
Father Eusebius, called to the eremitical life, \par
Father Eusebius, having miraculous visions, \par
Father Eusebius, faithful priest of Mother Church, \par
Father Eusebius, Canon of Estergom, \par
Father Eusebius, honour of the Hungarian people, \par
Father Eusebius, great treasure of the poor, \par
Father Eusebius, angel of purity, \par
Father Eusebius, inflamed with divine love, \par
Father Eusebius, living spring of humility, \par
Father Eusebius, example of perfect poverty, \par
Father Eusebius, untiring in the practice of penance, \par
Father Eusebius, a true devotee of the Blessed Virgin Mary, \par
Father Eusebius, faithful disciple of St. Paul and St. Anthony, \par
Father Eusebius, blooming lily of the desert, \par
Father Eusebius, enamoured of the Word of God, \par
Father Eusebius, sign of unity for hermits and monks, \par
Father Eusebius, renovator of the eremitical life, \par
Father Eusebius, great lesson of perseverance, \par
Father Eusebius, spurning everything for God, \par
Father Eusebius, loving all people, \par
Father Eusebius, perfect follower of Christ, \par
Father Eusebius, sanctified by a life of prayer, \par
Father Eusebius, contrite for sins at the hour of death, \par
Father Eusebius, ideal of union with God, \par
Father Eusebius, joy of Angels, \par
Father Eusebius, adoring the Face of God, \par
\end{litany}
\medbreak
\begin{litany}That we may persevere in faith, hope and love, \responce{Intercede to God for us.}\par 
That we may imitate the humility, obedience and justice of Jesus Christ, \par
That we may learn to discern perfectly good from evil, \par
That we may daily examine our consciences, \par
That we may sympathize with the poor and suffering, \par
That we may often renew ourselves in the Sacrament of Penance, \par
That we may use well our gifts and talents, \par
That we may discover the value of prayer and sacrifice, \par
That we may contemplate the Face of God someday, \par
\end{litany}
\medbreak
\begin{samepage}
\prayer{L.}{\textbf{Let us pray:}\\ O God, Almighty Father, you have shown your boundless mercy and abundance of graces to our Father Blessed Eusebius, by allowing him to serve you in solitude and contemplation. We ask you humbly, through his intercession, to inflame us with the flame of your divine love, that following his example, we may honour you all the days of our life. Through Christ our Lord.}
\prayer{All}{Amen.}
\end{samepage}
\medbreak
\newpage
\includeimage{saint-augustine.pdf}\pagestyle{empty}
\hspace{0pt}
\vfill

\textit{``You called us, Lord that we may search and praise you
with joy. You have created us for yourself; our hearts will remain
restless until they rest in you.'' (St. Augustine, Confessions)}
\medbreak
O unmatchable shepherd of the Lord's flock. You are an
example for all monks and priests. May your prayers implore for us
God's clemency.
\vfill
\textit{St. Augustine by Fr. Augustyn Jedrzejczyk Art Collection of Jasna
Góra, 1903}
\hspace{0pt}
\newpage
\newpageheader{Contemplation of God in His Saints}{Saint Augustine}

\section{Saint Augustine}

\instruction{St. Augustine (354-430), bishop of Hippo, Doctor of the Church is the Father and Legislator of our Order. As the result of his deep conversion in 385 after a dissolute life style, he was baptized by St. Ambrose at the age of 33 and admitted to the Catholic Church on December 8th. Shortly after, he resigned from the post as the teacher of rhetoric, sold his inheritance, and started a monastic life, which he organized in Tagaste (North Africa) in 391. The same year, he was ordained a priest and dedicated himself to the studies of Scriptures. In 396 he was consecrated Bishop of Hippo. In his pastoral administration, St. Augustine cared for the clergy and monastic life. He combated Manichees, the Donatists and Pelagius. His works: The Confessions, City of God and over 500 homilies are well known and appreciated in the West. For his sublime achievements in theology and spirituality, the Latin Church rewarded him with the title: ``Doctor of Grace.'' He composed a Rule that is observed by some 150 religious communities today. It was his Rule that was given to our Founding Fathers to follow: the Brothers of St. Paul the First Hermit. On December 13, 1308 with the decree ``Ubi Saecularia'', signed by Pope Clement V, through the mediation of his Legate, Cardinal Gentilis de Monte Florido, our Order was officially approved by the Church and the Rule of St. Augustine was assigned to it to live by.}
\medbreak

\subsection{Litany of St. Augustine}

\begin{litany}Lord, have mercy, \responce{Lord, have mercy.}\par 
\end{litany}
\begin{litany}Christ, have mercy, \responce{Christ, have mercy.}\par 
\end{litany}
\begin{litany}Lord, have mercy, \responce{Lord, have mercy.}\par 
\end{litany}
\begin{litany}Christ, hear us, \responce{Christ, graciously hear us.}\par 
\end{litany}
\begin{litany}God the Father of Heaven, \responce{Have mercy on us.}\par 
God the Son, Redeemer of the World, \par
God the Holy Spirit, \par
Holy Trinity, One God, \par
\end{litany}
\medbreak
\begin{litany}Holy Mary, \responce{Pray for us.}\par 
\end{litany}
\begin{litany}Holy Father Augustine, \responce{Implore the Heart of God for us.}\par 
Son for whom St. Monica shed bitter tears, \par
Father and Founder of the Order of Hermits, \par
Chosen light of Doctors, \par
Devotedly studying the Holy Scriptures, \par
Teacher of theologians, \par
Example of apostolic life, \par
Outstanding preacher of the Word of God, \par
Defender of faith, \par
Wounded with the love of Christ, \par
Bishop of great humility and mercy, \par
Light scattering the darkness of errors, \par
Implore the heart of God for us. \par
Source of divine eloquence, \par
Father of many orders living according to your rule, \par
Charitable defender of the troubled who come to you, \par
Consoler of the suffering and distressed, \par
Through your wondrous conversion, \par
Through the zeal with which you defended the faith, \par
Through your labours with which you supported the Church, \par
Lover of contemplation and community life, \par
Adorer of the Most Holy Trinity, \par
Ardent devotee of the Mother of God, \par
Our model and protector, \par
\end{litany}
\medbreak
\begin{litany}That we may constantly make progress on our way to perfection, \responce{Support us.}\par 
That we may faithfully observe the directions of your rule, \par
That we may develop the fraternal love of mutual care, \par
That we may always long for God, \par
That we may care for the salvation of others, \par
That we may be of one heart and of one soul united with God, \par
That our restless hearts may find their rest with you in God, \par
\end{litany}
\medbreak
\noindent\begin{minipage}{\textwidth}
\begin{litany}Lamb of God, you take away the sins of the world, \responce{Spare us, O Lord.}\par 
\end{litany}
\begin{litany}Lamb of God, you take away the sins of the world, \responce{Graciously hear us, O Lord.}\par 
\end{litany}
\begin{litany}Lamb of God, you take away the sins of the world, \responce{Have mercy on us.}\par 
\end{litany}
\end{minipage}
\medbreak
\begin{samepage}
\prayer{L.}{\textbf{Pray for us Saint Augustine, our Father,}}
\prayer{All}{That we may be made worthy of the promises of Christ.}
\end{samepage}
\medbreak
\begin{samepage}
\prayer{L.}{\textbf{Let us pray:}\\ Almighty God and Father, we pray that you renew your Church with the spirit which you gave to your Bishop, St. Augustine. May we all be filled with the same spirit and seek only you; for you are the source of all wisdom, and the giver of perfect love. Through Christ our Lord.}
\prayer{All}{Amen.}
\end{samepage}
\medbreak

\subsection{Prayer to St Augustine}

O St. Augustine, wondrous Doctor of Grace! You wrote and sang of God's merciful love that worked wonders in your soul. Help us to trust in God's help alone. You knew the tragedy of sin. You experienced the darkness of error and the chains of pride and sensuality and endured the pain of alienation from God. But the Lord saw the tears and heard the supplications of your saintly mother, Monica. Because of her unceasing prayer, God in his great mercy, rescued you from the darkness of Satan and immersed you in his divine light.
\medbreak

O great St. Augustine! Help us to discover God as our eternal Truth, as our true Love and as our beloved Eternity. Teach us to believe in Him and to remain in the state of grace always. Teach us to pray constantly to overcome the weakness of our flesh. Accompany us on our pilgrimage to the endless Life. Pray for us to know the Lord better, to love Him more and to be worthy to reign with Him for ever and ever. Amen.
\medbreak
\newpage
\includeimage{saint-casimir.pdf}\pagestyle{empty}
\hspace{0pt}
\vfill

\textit{``Happy those servants whom the master finds awake
when he comes. I tell you solemnly, he will put on an apron, sit
them down at table and wait on them'' (Lk 12:37).}
\medbreak
O Merciful God! Give us the grace to imitate St. Casimir,
our spiritual Brother with a joyful and loving heart and fulfil perfectly your commandment of love. Through Christ our Lord.
Amen.
\vfill
\textit{The Holy Communion of the Royal Jagiellonian Family in the miraculous Chapel.
Workshop of Jasna Gora. Art Collection of Jasna Gora,
2nd part of the XVIIth century.}
\hspace{0pt}
\newpage
\newpageheader{Contemplation of God in His Saints}{Saint Casimir}

\section{Saint Casimir}

\instruction{St. Casimir, a son of King Casimir Jaggielonczyk and Elizabeth Rakuska was born October 3, 1458 in Kraków at the royal castle of Wawel. He died March 4, 1484 and was buried at the Cathedral of Vilnius (Lithuania) where his relics are venerated to this day. His cult was approved by the Holy See in 1521 and ratified again by Pope Clement VIII in 1602. The same Pope also affirmed the liturgical texts of the Holy Mass and Divine Office and extended the feast of St. Casimir to be celebrated by the Universal Church.}
\medbreak

\instruction{St. Casimir was noted for his special grace of prayer and a filial devotion to the Mother of God, His favourite hymn to Mary was, ``Omni die, dic Mariae'' (sing to Mary every day). The text of this hymn was enclosed in his coffin. Along with his father he used to come to the Shrine of Our Lady of Jasna Gora. In 1472 the whole royal family prayed before the Miraculous Image and received Holy Communion. On that day, the royal family was admitted to the Pauline Confraternity, including the future saint, St. Casimir.}
\medbreak

\instruction{St. Casimir is a Saint of two nations: Poland and Lithuania. In 1950, Pope Pius XII proclaimed him patron of the youth of Lithuania.}
\medbreak

\subsection{Litany of St. Casimir}

\begin{litany}Lord, have mercy, \responce{Lord, have mercy.}\par 
\end{litany}
\begin{litany}Christ, have mercy, \responce{Christ, have mercy.}\par 
\end{litany}
\begin{litany}Lord, have mercy, \responce{Lord, have mercy.}\par 
\end{litany}
\begin{litany}Christ, hear us, \responce{Christ, graciously hear us.}\par 
\end{litany}
\begin{litany}God the Father of Heaven, \responce{Have mercy on us.}\par 
God the Son, Redeemer of the World, \par
God the Holy Spirit, \par
Holy Trinity, One God, \par
\end{litany}
\medbreak
\begin{litany}Holy Mary, \responce{Pray for us.}\par 
St. Casimir, \par
St. Casimir, lily of purity, \par
St. Casimir, torch burning brightly in the House of God, \par
St. Casimir, Patron of Lithuania and Poland, \par
St. Casimir, example for the powerful of the world, \par
St. Casimir, mirror of justice, \par
St. Casimir, strength of the weak, \par
St. Casimir, model of patience, \par
St. Casimir, detester of royal honours, \par
St. Casimir, imitator of God's generosity, \par
St. Casimir, faithful disciple of Christ, \par
St. Casimir, loving devotee of Mary, \par
St. Casimir, consoler of the distressed, \par
St. Casimir, visiting the sick, \par
St. Casimir, helping the dying, \par
St. Casimir, pure source of goodness, \par
St. Casimir, defender of the holy faith, \par
St. Casimir, father of the poor and orphans, \par
St. Casimir, youth of innocence, \par
St. Casimir, example of humility and meekness, \par
St. Casimir, joy and glory of Vilnius, \par
\end{litany}
\medbreak
\noindent\begin{minipage}{\textwidth}
\begin{litany}Lamb of God, you take away the sins of the world, \responce{Spare us, O Lord.}\par 
\end{litany}
\begin{litany}Lamb of God, you take away the sins of the world, \responce{Graciously hear us, O Lord.}\par 
\end{litany}
\begin{litany}Lamb of God, you take away the sins of the world, \responce{Have mercy on us.}\par 
\end{litany}
\end{minipage}
\medbreak
\begin{samepage}
\prayer{L.}{\textbf{Pray for us, Saint Casimir,}}
\prayer{All}{That we may be made worthy of the promises of Christ.}
\end{samepage}
\medbreak
\begin{samepage}
\prayer{L.}{\textbf{Let us pray:}\\ O God our Father in Heaven, you have given St. Casimir the grace of perseverance in the virtue of sanctity at the royal court, exposed to the temptations of the world. Grant us, through his intercession to withdraw our hearts from changeable allurements and seek the eternal joys that last forever. Through Christ our Lord.}
\prayer{All}{Amen.}
\end{samepage}
\medbreak
\newpage
\includeimage{saint-hedwig.pdf}\pagestyle{empty}
\hspace{0pt}
\vfill

\textit{``The Lord God blessed her richly and crowned her with
the diadem of pure gold.''}
\medbreak
Almighty, Eternal God! We pray that following the noble
example of the Queen St. Hedwig, we may serve generously all our
brothers and sisters in their temporal and spiritual needs. Through
Christ our Lord. Amen.
\vfill
\textit{The Queen St. Hedwig by Augustyn Jodrzejczyk; Art Collection of
Jasna Góra, 1933}
\hspace{0pt}
\newpage
\newpageheader{Contemplation of God in His Saints}{Saint Hedwig}

\section{Saint Hedwig}

\instruction{St. Hedwig, Queen of Poland (1374 - 1399). During the Canonization Mass, June 8, 1997 on the Great Meadows of Cracow, Pope John Paul II in his homily, said the following, ``It was from the Black Crucifix, venerated at the royal Cathedral of Wawel that St. Hedwig learned her deep wisdom and love for the poor and suffering. It was from the contemplation of that famed Cross and her personal union with the Crucified Lord that she never lost the ``best part'' (cf. Lk 10: 42), the presence of Christ. Today, we wish to kneel with you, St. Hedwig, at the foot of that Cross of Wawel, to be able to hear an echo of that divine lesson you once heard. From you, St. Hedwig, we want to learn how to realize this divine lesson in our days.''}
\medbreak

\instruction{St. Hedwig is very close to our Pauline Order for many reasons because of the historical ties between Poland and Hungary. Her father, Louis the Angevin was a great benefactor and supporter of the Pauline Order in Hungry and in Poland through his deputy, the Prince Ladislaus of Opole, who according to tradition, brought the miraculous image of Our Lady of Czestochowa to Poland in 1383 and entrusted it to the Pauline Monks, invited from the Hungarian Monastery Maria Nostra. We are in possession of a document issued on February 24, 1393 and signed by the King Ladisalus Jagiello, the King of Poland and husband of St. Hedwig, who confirmed the foundation of the Monastery of Jasna Gora, initiated by Prince Laislaus of Opole. It is most probable that the King was greatly influenced by his saintly wife, St. Hedwig, to make such a decision. Her mission was one of personal sacrifice to bring the nations to the Catholic Church. As a young lady, she was chosen as Queen of Poland by the Polish nobles. She fulfilled her mission of personal sacrifice when she broke her engagement with Wilhelm Habsburg, and married a much older and wild Duke of Lithuania, Ladislaus Jagiello, the future King of Poland, in order to bring him and all the Lithuanians to the Catholic Faith. Thus, she became the ``godmother'' of Catholic Lithuania and began a mission of evangelization. With her great heart and wisdom she embraced all her subjects with love, especially the poor and underprivileged. She gave away much of her personal treasures to found the Jagiellonian University in Kraków, one of the best universities in Europe to this day.}
\medbreak

\instruction{The life of St. Hedwig is deeply associated with the mystery of the Cross. It serves the Pauline Order as an eminent reminder of its mission of the Cross. The first Monastery of the Paulines in Hungary was founded under the sign of the Holy Cross. In Kraków the Pauline Monastery on the ``Skałka'' (a little rock) is a short distance from Wawel Castle. In both places, the mystery of the Cross remains an object of deep contemplation and a source of spiritual strength. St. Hedwig gives us a shining example of that special love of the Holy Cross.}
\medbreak

\subsection{Invocations to St. Hedwig the Queen}

\begin{litany}St. Queen Hedwig, \responce{Pray for us.}\par 
St. Hedwig, faithful daughter of God, \par
St. Hedwig, always united with God, \par
St. Hedwig, model of a royal wife, \par
St. Hedwig, caring for your poor subjects, \par
St. Hedwig, founder of churches and monasteries, \par
St. Hedwig, apostle of Christian Faith, \par
St. Hedwig, restorer of the Academy of Kraków, \par
St. Hedwig, protector of young students, \par
St. Hedwig, propagator of peace, \par
St. Hedwig, tried by calumnies, \par
St. Hedwig, patroness of Poland, \par
\end{litany}
\begin{litany}That we may profess our faith courageously, \responce{Intercede to God for us.}\par 
That we may live a virtuous life, \par
That the Polish Nation may remain united in Christ, \par
That our families may flourish with Christian virtues, \par
That we may cherish a special love and devotion to the Most Holy Virgin Mary, \par
That we may bear patiently our sufferings and discomforts of life, \par
That our lives may be guided by the will of God, \par
\end{litany}
\medbreak
\begin{samepage}
\prayer{L.}{\textbf{Let us pray:}\\ O God, you are the life and the glory of the poor. You made St. Hedwig, Queen, a zealous propagator of Christian faith and love; grant us through her intercession we may become apostles of your truth and goodness. Through Christ our Lord.}
\prayer{All}{Amen.}
\end{samepage}
\medbreak

\subsection{Prayer to St. Hedwig the Queen}

Our Beloved Queen Hedwig, you have made a great personal sacrifice in order to bring many souls to salvation and have become the Apostle of Lithuania. You propagated the Catholic Faith and built churches of divine worship. You dried up many tears of sorrows during your life here on earth and hereafter. Many sick recovered their health through your care and prayer. You ruled your subjects with goodness, meekness and mercy.
\medbreak

Therefore, we ask you today to protect and defend all families and particularly our religious family from all dissensions and troubles. Pray for the conversion of poor sinners and implore for them the grace of eternal salvation.
\medbreak

We trust in you, our Queen Hedwig and we hope that you will never fail us. As in the past, your spirit intercedes to God for us today. In heaven you contemplate the glorified face of Jesus. Pray to him on our behalf and tell him, ``Lord who will return their tears?'' The Lord will listen to you now as He once listened to you in the Wawel Cathedral. In your compassion, look down on your people, often discouraged and oppressed. Bring them health of body and consolation of the spirit. For all your favours bestowed upon them, they will bless you forever. St. Hedwig, pray for us!
\medbreak
\newpage
\begin{appendices}
\newpageheader{Prayers for Vocations}{Prayers for Vocations}

\chapter{Prayers for vocations}

\instruction{The Constitutions of our Order impose on all the members an obligation to foster and pray for religious vocations as the Lord has requested, ``the harvest is rich, but the labourers are few, so ask the Lord of the harvest to send labourers to his harvest (Mt 9:27). Encouraged, therefore, by Christ himself, ``let all members intercede before God that he might constantly renew our Monastic Community and enrich it with new members'' (art. 63). Also, the General Directory reminds us of our duty to pray for an increase of good vocations, ``each Pauline Community must pray for many solid vocations to the Order, especially during the noon adoration of the Most Blessed Sacrament. Moreover, Superiors are responsible for arranging in each of our churches a designated week and specified days of prayer for vocations'' (norm 72).}
\medbreak

\section{Prayer of Pope Pius VI}

O Jesus Christ the Divine Shepherd, you have sent your Apostles to be the shepherds of souls. Draw to yourself the warm and noble minds of young and generous men and make them your followers and servants of all.
\medbreak

Grant that they may share your desire for the universal salvation for which you renew your Sacrifice on our altars everywhere. You are the Lord of all, who lives forever. Therefore, open before them new horizons that they may see and hear an unspeakable prayer for the light of truth and warmth of love of all their brothers and sisters today. Give them the courage to say ``yes'' to your invitation to continue your mission on earth in building up your Mystical Body, the Church. Make them the light of the world and the salt of the earth.
\medbreak

We pray also that you may inspire many pure and magnanimous hearts of women and pour into them the desire for Evangelical perfection and dedication to serve the Church in their sisters and brothers in need. Amen.
\medbreak

\section{Prayer for religious vocations}

O Good Jesus, our Lord and Saviour! Remember the words you have spoken; ``Ask the lord of the harvest to send labourers to the harvest''. With an unshakeable trust, through the intercession of your Immaculate Mother and our Mother, the Queen of Jasna Gora and through the intercession of St. Paul the First Hermit, our Father, we pray that you may look down kindly upon our religious family and bless us abundantly. May we live in the spirit of our Constitutions and ever increase in your grace before God and men. May it grow in number with new sons, whom you call for your greater glory.
\medbreak

Speak to them again today, ``Come follow me'' Make them worthy of your confidence and make them well disposed to carry on the mission of your Church. May they shine forth with personal sanctity, knowledge and untiring love in spreading your Kingdom on earth. To you along with the Father and the Holy Spirit be glory and honour for ever and ever. Amen.
\medbreak

\section{Prayer for fervent priests and religious}

O God of Mercy! You desire the salvation of all people. For this, your only begotten Son came to us as the Good Shepherd to search out and rescue those who were lost. You have given your Son his first disciples and priests to carry on this work of redemption in his Name until the end of time. To us however, you have recommended we ask you, the Lord of the Harvest, to send labourers to your vineyard. Therefore, we pray most fervently, that you may give to the Holy Church of your Son many fervent, saintly priests and religious. Choose and call them yourself from among your people and that no one undesirable to you may remain with us and that no one called by you be unwelcome. Guide them with the grace of the Holy Spirit as you once prepared your Apostles and the Disciples of Christ.
\medbreak

May they find joy in the contemplation of your Word. Pour into their hearts your warm love towards you, fidelity and filial affection for the Holy Church and St. Peter, the Rock, on which you have built your Holy Church.
\medbreak

May they work tirelessly in your vineyard always with prudent and reasonable fervour. Give them a spirit of gentleness and compassion for the weak and erring. May they expect the results of their labour only from you, in humility, prayer and mortification. May they look for nothing else, except your glory and the salvation of souls. We ask all this through your only-begotten Son, Jesus Christ, the Eternal High Priest and the Shepherd of our souls, who lives and reigns with you in the unity of the Holy Spirit, One God for ever and ever. Amen.
\medbreak

\section{Prayer to Our Lady of Jasna Gora for priestly vocations}

Our Lady of Jasna Gora, Most Holy Mother of Jesus Christ, our High Priest! Your Son has made you the Mother of all priests when he entrusted his beloved disciple John to you as your son.
\medbreak

Mother of Priests! We ask you to watch over those your Son chooses to be his representatives and sharers in his priestly dignity. Pray for them, that they never lose this priestly vocation but remain faithful to the end, working fervently for the souls entrusted to their pastoral care.
\medbreak

Mother of Priests! Pray for our fathers and mothers that they may offer their Sons with joy to serve your Son as priests.
\medbreak

Mother of Priests! Implore for us the grace of understanding and appreciation for all the blessings which God bestows upon us through his holy priests. Amen.
\medbreak

\section{Prayer to the Holy Spirit for good priestly vocations}

O God the Holy Spirit! You are the Supreme Teacher of all the faithful. On the Day of Pentecost, you descended upon the Apostles gathered in the Cenacle, in the form of fiery tongues, that they would become worthy dispensers of your divine mysteries. We pray that you may bestow the fullness of your gifts on all priests who are administering the mission of Jesus Christ and are your instruments in guiding the faithful to salvation.
\medbreak

Grant them the gifts of wisdom, skill, knowledge, counsel, courage and piety. Fill them with the holy fear of God. Invest them with your abundant graces, that they may teach us in word and example and lead us safely to eternal life, where you receive glory with God the Father and his Son Jesus Christ, forever and ever. Amen.
\medbreak

\section{Prayer for worthy priests and religious}

O God the Shepherd and Teacher of the faithful! To preserve the deposit of your Church and to ensure its salvific mission, your have instituted the holy Priesthood. You have told your Apostles, ``The harvest is rich, but the labourers are few. Ask the Lord of the Harvest to send more labourers into his harvest.''
\medbreak

Behold, we approach you with our heartfelt request and we beg you to send more priests and religious to your Holy Church.
\medbreak

O Lord, give them the courage to say ``yes'' to your calling. Protect them from all dangers of the world. Invest them with the spirit of counsel, strength, wisdom and piety. Fill them with the spirit of your holy fear, that blessed with the precious gift of priesthood or religious profession, they may teach us how to walk on the road of your commandments and bring us to union with you, who lives and reigns for ever and ever. Amen.
\medbreak

\section{Prayer of Pope John Paul II}

O Jesus the Good Shepherd! Accept our humble adoration and thanksgiving for all vocations which, through your Holy Spirit, you give unceasingly to your Church. Support the bishops, priests, missionaries and all consecrated persons. Grant them the grace to be an example of a true Evangelical life. Give them the strength of perseverance in their resolutions in preparation for their holy service. Multiply the workers of the Gospel to bring the Good News to all nations. Surround the children and youth of our families and of our religious communities with your protection. Give them the spirit of generosity and magnanimity to follow you, especially in our times. Give all the invited the energy to leave everything and to choose you, who are Love. Forgive the hesitations and infidelities of those you have already called.
\medbreak

O Christ our Lord! Graciously hear our petitions through the intercession of the Most Holy Virgin Mary, your Mother and the Queen of the Apostles. May she who believed and generously said ``yes'' to your invitation, be the cause of joy, and accompany those you have called with her maternal presence and example in the service for your Kingdom. Amen.
\medbreak

\section[Prayer for vocations]{Prayer for vocations \\ {\small\normalfont (Message of the Holy Father John Paul II \\ delivered at the XXXII World Day of Vocations).}}

Lord Jesus Christ, our Master! We stand before you and intently listen to your calling, ``Ask the Lord of the Harvest to send labourers to his harvest'' (Mt 9:38). Moved with these words, which you constantly direct to us, we pray that the young people of today may clearly discover and make the right choice for their state of life. Conscious of a great responsibility for an increase of vocations, we beg you with trust and insistency for this great gift for the Church. Every new vocation is a precious gift, obtained only by prayer and holiness of life.
\medbreak

You O Lord are the source of every vocation. You choose for each of us the way of life and give us a personal role to play in your Church. Open the hearts and minds of young men and women, that a fresh bloom of holy vocations may remind the world of your love, and that all may come to know Jesus Christ, the true Light who has come into the world to bring an unfailing hope of eternal life to all men.
\medbreak

Be praised Lord in all those, who have discerned their vocation, in accordance with your thought --- and lead them on the right path.
\medbreak

We thank you, God, for the pastoral care of the church for the youth, and bound with this care that of vocations. The young man, living the gift of youth, stood by you and asked:
\medbreak
\begin{itemize} \item What do I need to do, to gain eternal life? \item What do I need to do, for my life to have purpose? \item What is God's plan for my life? \item What is God's will for me? \end{itemize}\medbreak

Lord Jesus, your loving gaze gives youthful hearts the capacity to ask these essential questions. The experience of the meeting between the human heart and Your heart, filled with gentle goodness, gives the strength to seek the answers. Awareness that God loves us allows us with an open heart to receive your answer. Follow me!
\medbreak

In following you, youth reveals itself, full of the possibilty of reaching its fullness.
\medbreak

\rubric{*} You are the Way, the Truth and the Life
\medbreak

In following you, the youth discover the meaning of life as the gift of self and the beauty and truth of growing in Love.
\medbreak

\rubric{*} You are the Way, the Truth and the Life
\medbreak

\section{Prayer for those preparing for priesthood and religious life}

O Lord Jesus Christ! You have offered your own life for the salvation of all peoples and have become the Eternal Priest. We pray for those, who are preparing themselves for the priesthood and religious life. Surround them with the grace of your blessing. Shield them from all temptations and confirm them in holiness. Make them ever eager to know you better and love you more. We pray for the candidates, novices and seminarians, that they make steady progress in following you, and remain united with you always. For you live and reign for ever and ever. Amen.
\medbreak

\section[Prayer of St. John Paul II for the Intention of vocations to the Priesthood and Consecrated life]{Prayer of St. John Paul II for the Intention of vocations to the Priesthood and Consecrated life \\ {\small\normalfont Message for the XXXV World Day of Prayer for Vocations}}

\begin{litany}Spirit of eternal Love, \par
who proceeds from the Father and the Son, \par
we thank you for all the vocations \par
of apostles and saints which have enriched the Church. \par
Continue, we pray, this work of Yours. \par
Remember when, at Pentecost, \par
you descended on the Apostles gathered together in prayer \par
with Mary, the mother of Jesus, \par
and look upon your Church which today has \par
a particular need for holy priests, \par
of faithful and authoritative witnesses of your grace; \par
she needs consecrated men and women, \par
who show the joy of those who live only for the Father, \par
who make their own the mission and the offering of Christ, \par
who build up in charity the new world. \par
Holy Spirit, eternal Spring of joy and peace, \par
it is you who open the heart and the mind to the divine call; \par
it is you who make effective every impetus \par
towards good, towards truth, towards charity. \par
Your "inexpressible groans" \par
rise up to the Father from the heart of the Church, \par
which suffers and struggles for the Gospel. \par
Open the hearts and minds of young men and women, \par
so that a new flowering of holy vocations \par
may show forth the fidelity of your love, \par
and all may know Christ, \par
the true light come into the world \par
to offer to every human being \par
the sure hope of eternal life. Amen. \par
\end{litany}

\section[Prayer to the Most Holy Virgin Mary]{Prayer to the Most Holy Virgin Mary \\ {\small\normalfont Apostolic Exhortation on consecrated life ``Vita Consecrata''}}

Mary, image of the Church, the Bride without spot or wrinkle, which by imitating you ``preserves with virginal purity an integral faith, a firm hope and a sincere charity'', sustain consecrated persons on their journey toward the sole and eternal blessedness.
\medbreak

To you, Virgin of the Visitation, do we entrust them, that they may go forth to meet human needs, to bring help, but above all to bring Jesus. Teach them to proclaim the mighty things which the Lord accomplishes in the world, that all peoples may extol the greatness of his name. Support them in their work for the poor, the hungry, those without hope, the little ones and all who seek your Son with a sincere heart.
\medbreak

To you, our Mother, who desire the spiritual and apostolic renewal of your sons and daughters in a response of love and complete dedication to Christ, we address our confident prayer. You who did the will of the Father, ever ready in obedience, courageous in poverty and receptive in fruitful virginity, obtain from your divine Son that all who have received the gift of following him in the consecrated life may be enabled to bear witness to that gift by their transformed lives, as they joyfully make their way with all their brothers and sisters toward our heavenly homeland and the light which will never grow dim.
\medbreak

We ask you this that in everyone and in everything glory, adoration and love may be given to the Most High Lord of all things, who is the Father, the Son and the Holy Spirit. Amen.
\medbreak
\newpage
\includeimage{pauline.pdf}\pagestyle{empty}
\hspace{0pt}
\vfill

In the spirit of the charism of the Order, which harkens back to the tradition of the Egyptian desert hermits, the peaching of the Word of God, heard in solitude, confronted in the community of the brothers, becomes an answer to the earnest request of the pilgriming people to the hermits: ``Abba tell us a word!''.
\vfill
\textit{A Pious Pauline Novice, painted at Jasna Góra in the first half of the XVIII century.}
\hspace{0pt}
\newpage
\newpageheader{Meditation}{Meditation}

\chapter{Meditation}

The Pauline spirituality enrichens itself with the unceasing, centuries old, varieties of forms of prayer. However above all we need to choose the \textit{optiamam patrem} (the better part) of Mary the sister of Lazarus and Martha, which is meditation, the beginning of the spiritual life. Not without the spice of attic salt, did Fr Andrew Eggerer OSPPE write about our semi-contemplative life, in his first volume of his Pauline annals, published in the year 1663: ``so too did our fathers in the hermitages and then in the cenobiums of Hungary. Thus fulfilling the service of the busy Martha, at the same time they did not cease abiding with Magdalen at the feet of Jesus, giving themselves over to the contemplation of heavenly things. Let them remember of how they ought not to imitate those who today in monasteries live only for their stomachs, complaining about the yoke of the work of their own hands.''
\medbreak

Quite so, in those now distant years, the discerned inspiration of the Holy Spirit, found its place in our spirituality of meditations, which although it is a passive prayer, beginning in secret, bears its results visibly not only in the whole life of the monk, but also in the Order and finally in the Church. The almighty is a hidden God, and His Son led a hidden life for many years, only to shine later in the preaching of the Gospel and the numerous miracles which finally shone forth in his conquering of death, and in his resurrection.
\medbreak

Consecrated persons, likewise, experience an incomparable role in the hidden life of Jesus, in which they practice prayer, which is not for show, but seeks the closed doors of the cell in which God sees that which is hidden and repays that which we cannot fully express (cf. Mt 6:6). In the decree \textit{Perfectae Caritatis, }that is the decree on the adaptation and the renewal of Religious life, in paragraph 6, we read:
\medbreak
\begin{addmargin}[\parindent]{\parindent} 6. Let those who make profession of the evangelical counsels seek and love above all else God who has first loved us (cf. 1 John 4:10) and let them strive to foster in all circumstances a life hidden with Christ in God (cf. Col. 3:3). This love of God both excites and energizes that love of one's neighbour which contributes to the salvation of the world and the building up of the Church. This love, in addition, quickens and directs the actual practice of the evangelical counsels. Drawing therefore upon the authentic sources of Christian spirituality, members of religious communities should resolutely cultivate both the spirit and practice of prayer. In the first place they should have recourse daily to the Holy Scriptures in order that, by reading and meditating on Holy Writ, they may learn ``the surpassing worth of knowing Jesus Christ'' (Phil. 3:8). They should celebrate the sacred liturgy, especially the holy sacrifice of the Mass, with both lips and heart as the Church desires and so nourish their spiritual life from this richest of sources. So refreshed at the table of divine law and the sacred altar of God, they will love Christ's members as brothers, honour and love their pastors as sons should do, and living and thinking ever more in union with the Church, dedicate themselves wholly to its mission. \end{addmargin}\medbreak

The conciliar decree explains the manner of life of consecrated persons and defines them, as hidden in Christ, the foundations of such an approach we see in the every-day immersion in the Sacred Scriptures. Twice appears the phrase, ``above all'', highlighting the priority of this practice, which forms our order's identity: above all prayer, and first prayer amongst prayers, the seeking of Christ in the Sacred Scriptures! Secondly, the Council draws attention to the tight link between meditation and preaching with the Eucharist. To learn about Christ not only out of sheer curiosity or out of intellectual endeavour, but solely to create a relationship of Love, of friendship, which will always deepen, which will not cease but will overflow onto the People of God. This is the fundamental vocation of every consecrated person. For this to happen, meditation of the Word of God, must have a privileged place in our Monastic society.
\medbreak

St Thomas Aquinas in his famous \textit{Summa Theologiae}, wrote: \textit{necessaria est iugis oratio} --- the yoke of prayer is necessary --- everyday, like in a yoke we ought to meditate, because a repeated habitual action, not episodic, lastingly changes a man. A monk without meditation is like a mindless person, capable only of unsettling himself and others with egocentric ambitions or fears. Jesus in every spare moment, utilised His time in meditation. He had the ethos of reading the scrolls of the Torah and personal prayer in mediation, of which writes the evangelist Luke, whilst speaking of the habits of Jesus, that his is an unchanging practice of prayer (cf. Lk 4:16-17; 22:39-40).
\medbreak

The ethos of a monk ought to be the same as the practice of Jesus. Ετος in Greek is the same as our word ``habit,'' that is; a permanent disposition towards action, the practice of a particular thing or expressing it in the light of the gospel, a sweet yoke and a light burden. Although giving time every day to mediation may seem like a burden for our bodies, nevertheless it makes our lives spiritually light. It should become our saving passion, because only passion protects us from damming addictions. The habit of prayer and meditation on the Word of God, not only protects us against addictions, but also opens heaven to us, because it is exactly that, a new plain of existence, knowing that the one we were born into will suffer destruction (cf. Lk 3:21). The passion of meditation contains strength surpassing all human obsessions and opens him to a higher level of existence, that is the sharing in the very life of God. St John Chrysostom believed that no one could remain in sin, who ceaselessly harnesses himself in the yoke of prayer and meditation on the Scriptures. In our newest Pauline constitutions, article 27 speaks of Sacred Scripture as the source of our meditation, which we must cultivate, whole heartedly in our everyday lives, for it is the fundamental method of dialogue with God and directs us towards Mary, who from her earliest earthly years, kept and meditated in her heart His Words.
\medbreak

An exceptional distinction for the order was the arrival, at a particular time, of the miraculous icon of Jasna Góra into the heart of our order, as well as the discernment of the movement of the Holy Spirit directing our community to care for the People of God, particularly those in marian shrines. She herself, even before she became the Mother of God, the Virgin listening and meditating, directs the Paulines. It does not behove us to omit the central detail of the architecture of the church that puts before us the pattern of imitating her, and thus in such a way, like her, to accept and come to know Jesus Christ.
\medbreak

The Mother of Jesus is a model of faith, because she is unmatched in the listening to and meditating on the mysteries of the Word of God. The freedom of God and the freedom of man meet uniquely in meditation, opening forth the event of the annunciation in Nazareth. God turns to man with the Word saturated with Life, for life is most radically attributed to God, whereby the reception of Him is very much similar to conception. The Synoptic Evangelists recall the visit of the relatives of Jesus, who standing outside a house in which He was teaching, sought to meet with Him. Then He said these words, which for many seem troublesome, when at the same time they are an invitation to accept with a discerning heart His words and thanks to them fulfil the will of God: ``My mother and My brothers are those who, hear the word of God and keep it.'' (cf. Lk 8:21).
\medbreak

Adrienne von Speyr interpreted this response of Jesus, as a renunciation of the rights of Motherhood in regards to the Son of God in favour of all who like Mary whole-heartedly nourish themselves on His Words. St Augustine in his homily on the Birth of Christ, said, ``The Mother who carried Him in her womb, we carry also in our hearts (\ldots ) Let us not be infertile -- let our hearts be fertile in God.'' St Ambrose similarly courageously commented on this mysterious response of Jesus: ``He fulfilled the will of the Father, that you would become the Mother of Christ.'' Because our thoughts are not His thoughts (Is 55:8), it is vitally important to devote time to meditation, so that the words of God read in the Bible become incarnate in the received experience of the Monk. So that he may become ever more similar to the Mother of the Saviour, just as it took place in the life of Mary, in whom the Word became flesh, and the thoughts of her meditated prayers were fulfilled in the events of her fate.
\medbreak

It is not by accident, that just as the landmark events for a Pauline, the clothing in the habit and Solemn Vows take place on Solemnities on which the church reads the gospel pericope of the annunciation (8 XII ad 25 III). A radical otherness of the Word of God to the human intellect always causes the monk to firstly purify his thoughts and quieten his interior before beginning to meditate. Fr Gregory Terecy, (+1660) in his \textit{Conversio morum Paulinorum, }recommends that in preparation for meditation, the monk casts far away from himself all useless thoughts, phantasms, memories and all connotations bounded up with human aspirations and fears, which depart from the subject of the meditation. Thus meditation for the Pauline is a personal incarnation -- thanks to this personal prayer, the presence of the Holy Spirit makes our consecrated lives really spiritual, not only by profession, but also by practice. Meditation upon the Word of God fully develops faith, creating the best environment for the living out of the Eucharist, as the source of our saving life in a vivacious, conscious and loving way. In the Pauline constitutions from the year 1930 we find more than an encouragement not to omit meditation, but we also find the Prior being obligated to whole heartedly care for the meditations of the monks entrusted to his care. For as we read in chapter XII ``Our Order, thanks to holy meditation and contemplation was once flourishing.''(68) Certainly the roots of such an obligation to meditation reach to the constitution from the year 1644, where in article 10, we can read even stronger words, where the growth of our order is dependent on this very form of prayer. A Pauline, even if he could not participate in the communal morning meditation or the evening one, then he is obliged to fulfil personally his mediation, even if he must give up his free time to do so. The Sub-prior was even obliged to question the monks on the effects of their meditations. The care then for this prayer, through the ages, was regarded as the foundation of the very life and strength of the order itself. Likewise, every fall in the spiritual life always began with a falling away from this sublime form of prayer. In the Rule of Bartholomew of the five churches, from the year 1215, dedicated to the Hermits of Patacs, meditation and prayer took first place among the many activities of the monks.
\medbreak

The Monks of Jasna Góra in the XVII century rose at midnight, so that after Matins, they could do half an hour of meditation. This is written in the manuscripts of Jasna Góra and it convinces us that among the forms of prayer, the Paulines gave great weight to this form. This was not a meditation removed from their lives nor from the Bible, drifting into ideal and theoretical subjects, because the Pauline \textit{Liber Meditationis,} (1630) proposes meditation under its first section, entitled: \textit{De magnitudine malitiae peccati -- }Concerning the great malice of Sin. During the expanse of the centuries, the methods of mental prayer were greatly influenced by Carmelite and Jesuit currents; nevertheless, they were inclined to follow the exhortations of the previous council, encouraging the practice of \textit{lectio divina. }It is necessary to underline that the method used by every community and each one of us is adapted to our spiritual characters, modifying and making them more vivacious and robust to use.
\medbreak

At the source of \textit{lectio divina }lies the method that was systemised by the Cathurisan Guido II, the Angelic (+1188) in his \textit{opusculum} entitled \textit{Scala Claustralium: LECTIO} (reading) \textit{MEDITATIO} (meditation), \textit{ORATIO} (prayer), \textit{CONTEMPLATIO} (contemplation). Due to Pauline Spirituality not being limited to the work of contemplation but likewise ordered towards the apostolic work, shares its own won spiritual treasures, through working in the confessional, through preaching, through retreats, days of recollection and through spiritual direction. We can confidently adapt the Carthusian method to our spiritual needs with one additional point: \textit{PRONUNTIATIO} (Proclamation) in accordance with the words of Our Lord: ``What I tell you in the dark, utter in the light; and what you hear whispered, proclaim upon the housetops.'' (Mt 10:27).
\medbreak

The mother of our redeemer did likewise when after Gabriel's annunciation, addressed to her in the darkness of meditation, left for the far away land of Ajn Karim. Then under the roof of the house of Elizabeth she could fulfil the hymn, which every day in vespers reminds us of the necessity of the transmission of the faith to others. In connection with this, it behoves us, not to treat our meditation solely as some private prayer, rather as a preparation for a homily, recollection or some other form of evangelisation. Every page of sacred scripture is like a well, from which we need to draw water, to quench the soul's thirst for the Holy Spirit, just as the Samaritan widow at Sychar was thirsty for Love. Because nothing encourages more than an example, let us offer a short example of meditation, to illustrate its process. Due to many of us using readings from the Liturgy for mediation, it is good to remember that the liturgical readings do not contain the entirety of the Sacred Scriptures.
\medbreak
\newpage
\newpageheader{Meditation}{A Proposed Meditation}

\section{A Proposed Meditation}

Having in mind the good of those among us who are in formation, let us use the readings from the lectionary in tome V, from the 33\textsuperscript{rd}rd Week of Ordinary time, year II, Tuesday: Rev 3:1-6;14-22; Ps 15, 1b-3a. 3bc-4.5; Lk 19:1-10.
\medbreak

At the \textit{LECTIO} stage, in recollection, we familiarise ourselves with the contents of the texts by reading all three pericopes to find a central theme. It is good, on the day before, particularly in the evening before, to read the texts and become familiar with them, so that even before we sleep, we can begin a preliminary meditation on the texts. It is important to note particular details that catch our attention. This work ought to be done in solitude, in accordance with our motto -- Alone with God alone. Our primary focus could be a word, which appears in two or three liturgical readings, or a thought, or a dominating idea throughout the texts. We should not ignore the biblical contexts of the chosen fragments, which we have chosen. The Bible too, itself was written for catechetical purposes and does not lack meaningful neighbouring contents. For this proposed meditation, we shall use the concept of \textit{seeking God.} The reading from revelations encourages us to seek Jesus without fooling ourselves. That by his name, therefore the identity of the one we are seeking would become a reality, which expresses one dead spiritually coming to life. That one masking his poverty and greed would become spiritually wealthy. Frequently an impoverished soul seeks gratification in material opulence. How very high does one have to climb spiritually to fulfil the requirements asked by Psalm 15? Only one who isa honest and free from greed can climb the mountain of the Lord. In the background of these readings is the gospel of St Luke, which already appears as a concrete reality, but not without its mystical shades, becomes \textit{sui generis} the matrix of conversion.
\medbreak

Finding all the common themes of the readings -- in this case -- the search for God and not wealth is the first goal of the \textit{LECTIO} stage. In this stage, we need to mention an additional step: \textit{DIVINATIO}, that is a deepening of the texts, directed by the inspiration of the Holy Spirit. A great help in this stage would be something of a biblical lexicon, a dictionary of symbols, or a biblical concordance. An inestimable treasure here is the knowledge of the original languages. That is why, during our time in the seminary, the \textit{ratio studiorum, }foresees at least the learning of the Greek language. Using an interlinear version and dictionary not only allows us to find the correct translation of words, but also to find the place in the bible, in which the word occurs, and through a comparison find a broader and fuller meaning of the word. A beneficial practice is the writing down of our meditations, a \textit{LECTIO} \textit{SCRIPTA}, for this helps in recollection and, particularly in the morning, fights against distractions, likewise it allows us to practise the style of a future homily. We must not forget that zelous evangelisation is also part of our charism (Pauline Constitutions 3.4).
\medbreak

The third stage is \textit{MEDITATIO}, that is discursive meditation. This is the leaving behind the continuity of theme, it is an attempt to discern the inspiration of the Holy Spirit, the taking root of the found ideas and a weaving of them together into a complete whole. All the possible thoughts capable of being reflected on should be written down. The act of writing multiplies thoughts and orders them and this in turn causes the meditation to become clear and more easily finds its place in the memory, almost like a seal. The mother of our Lord, when the Angel illuminated the will of God, meditated in her heart, but the original text used the word διαλογιζομαι (dialogizomai) which means meditating with care (Mt 6:7-8), with the desire to understand. Moreover, even to attempt to comprehend that which is incomprehensible or begets fear (Lk 5:21-22), or likewise the internal struggle (Lk 24:38) and finally the ideal that appears from the very depths of the subconscious (Lk 2:35). All these biblical passages concern themselves with an intense discernment, which does not lack internal tension, attempting to understand with effort, which will lead to an understanding, not only of this Divine mystery, but also our very inner most self, the personal subconscious. \textit{MEDITATIO} then is the work of the intellect, at such a level, which is achieved thanks to the grace of faith in knowing Jesus Christ, in which the mystery of man becomes known to man himself.
\medbreak

The coincidental text concerning Zacchaeus, is associated liturgically with the reading from Revelations and with Psalm 15, and may lead us to our next discovery in this state of MEDITATIO. Psalm 15 for the reader may be similar to a glance in a mirror, because which one of us can climb the mountain of the Lord's holiness, searching for Him, without any blemish of sin or fault? The just, speaking not only the truth with the lips, but also with his heart, who of us does not lie, who never harmed anyone, or has not discovered in themselves some greed or some higher aspirations? Zacchaeus was a greedy liar, having committed many frauds and the many people he harmed he held in contempt. His name, Zacchaeus, may be translated as ``pure'' or ``innocent,'' but his life greatly contradicted his name. His personage ideally fits the accusations in the apocalyptic letters addressed to Laodicea and Sardis, which speak of him who has a name, that means that he is alive, but is in fact dead. He was rich materially, but poor in conscience. Nevertheless, what could save him was the irresistible desire to see Jesus, which he expressed whilst searching for him. ``He wanted to necessarily see Jesus!'' When this happened, Zacchaeus' life changed dramatically, from a greedy abuser he became a merciful alms giver. Of particular value is the key phrase ``He wanted to see Jesus.'' A literal translation of the original is ``he sought to see Jesus.'' This searching ζετεο (Zeteo) appears again in the same Gospel of Luke, at this moment, thanks to the Greek dictionary or a biblical concordance. Mary along with Joseph were looking for Jesus in the temple of Jerusalem. After finding him, Mary says: ``Behold, your father and I have been looking for you anxiously.'' (Lk 2:48) Looking in an English -- Greek dictionary, we see, that this word, ζετεο can likewise also mean to discern, therefore to meditate, that is mentally and spiritually attempting to comprehend something. So thus, Mary and Joseph not only were searching for Our Lord in the streets of Jerusalem, but they were also searching for him with their hearts, with their anxious hearts. This unsettling anxiety, accompanies every person, until they find Our Lord. They sought for him with their souls, their thoughts, and their emotions with their troubled thoughts. The essence of real mediation is not to confine oneself solely to one's half-hour morning or evening meditation, but to personally, most deeply seek to know and to discover the truth not only about our redeemer, but also about ourselves. Setting the bounds of our dialogue with Our Lord in such a way, allows us to adhere to Him and come to love Him more fully. Searching strenuously does not limit itself to only half an hour, but expands itself to seek understanding, in all available aids: beginning firstly from the Fathers of the Church and ending with exegetical commentaries, lexicons and dictionaries, so as to comprehend the scriptures most fully. Following such a comprehension and understanding comes the possibility of drawing closer to Our Lord and only then gives us a chance for love. All evil stems from a focusing on oneself and all good from focusing all ones strength on God. The stage of \textit{MEDITATIO} can last even a few days, because the consultation of sources and the search of a library may take us quite some time. We should not be restrained only by the morning's half hour, but should use every free moment we have to deepen the subject, which we are mediating upon.
\medbreak

\subsection*{Oratio}

The thorough intellectual study, using all possible sources and means, as mentioned above, opens our horizons for the feeling of wonder, sadness, compunction, an honest look at one's self and awakens a desire for a better life to the one we have led to this moment. The movements of our emotions, which have accompanied us in our mediation, require us certainly to turn towards God in prayer, in the most personal form of prayer, deep and sincere. The wider the confession, the stronger a bond appears between Our Lord and ourselves. Identifying with Zacchaeus, we climb the Sycamore, which earlier symbolised human pride and the will for power or greed, thanks to the desire for seeing the Lord, becomes and illustration of the ascent to the highest levels of the desire for God, which is capable of changing our entire lives. Prayer in the sighs and words of the spirit in the end transform into the contemplative gazing upon the Lord. Therefore, the next stage of \textit{Lectio Divina} is \textit{CONTEMPLATIO}.
\medbreak

\subsection*{Contemplatio}

In this moment of our mediation, we ought to fall silent in all our dialogues with ourselves and even with the Lord Himself. Placing ourselves before the ineffable Word of God, which did the impossible, became man, and communicated to us the greatest news, pushes man to silence filled with awe and praise. God always remains incomprehensible, even though He should always be known. St Ephraim the Syrian wrote that to hold that one has understood The Word is the greatest desecration, because it implies the conviction that the knower is greater than the object of the thing known. The stage \textit{CONTEMPLATIO} allows man to remind himself of the seriousness of silence in the presence of the revealed mystery, which will always remain a mystery, regardless of the understanding and insight gained from Mediation. This stage is good to do in front of the Blessed Sacrament during a time of adoration, which ought to be a natural crowning of a monastic day. Without this gazing into the countenance of Christ, without adoration of the Blessed Sacrament in the atmosphere of the complete morning mediation, \textit{As a consequence life itself loses meaning, the faces of brothers and sisters are obscured and it becomes impossible to }\textit{recognise}\textit{ the face of God in them. (Starting Afresh from Christ III. 25) }Moments of contemplation, which allows us to be raptured into the gaze of the face of the Lord, through the ascent of the sycamore of Zacheus, the whole of our consecrated life is renewed, as St John Paul II expressed on his homily on the 2\textsuperscript{nd}nd of February 2001.
\medbreak

\subsection*{Pronuntiatio}

The whole material, gathered, written down and ordered within the few days becomes a holy basis, from which we can firstly share with the community in the form of \textit{COLLATIO}, which springs from the eremitical tradition and later with the faithful in the church during the homily or a spiritual conference. In this way mediation, which begins with being alone with God alone, becomes an enrichment for the brothers and finally with the People of God. The instruction of the congregation for institutes of consecrated life, 2002, entitled ``To begin afresh from Christ,'' tells us that consecrated persons only when they are themselves nourished with the Word of God, having become new people, free people, evangelic people and thus so authentic servants in the work of evangelisation and strongly convincing, are able to bring the fruit of our prayers to others. (III, 24) The riches of experience and reflection and finally the treasure of being enraptured by the face of the Lord, after a well prayed meditation, finds its crowning moment in adoring contemplation in His presence in the Blessed Sacrament, begs us to share our experience with others, for we are debtors of the Church no less than Zacchaeus, who promised to repay his debts, in the indebting promise he gave to the Lord. The conciliar decree on the ministry and life of priests, \textit{Presbyterorum Ordinis,} already in its first chapter, states \textit{their ministry, which begins with the evangelical proclamation}, although the peak is the leading of someone to the full participation in the Eucharist united with Our Lord. The Council even more insistently reminds the priest of this duty in the fourth article of the second chapter where it states \textit{The People of God are joined together primarily by the word of the living God, and rightfully they expect this from their priests.} Regardless of this obligation directed to priest, which for better or for worse, eloquently or not so, is always and everywhere to preach the Gospel is not left in any doubt, as to the role of meditation upon the Word of God, which is the first step to evangelisation.
\medbreak

If Zacchaeus did not climb the sycamore, to see Jesus, a great many people would not have been repaid or enriched by him. If a monk does not consciously strive to meditate, which is like climbing the tree of coming to the knowledge of one true and good God, a great many people would not have the debt of their salvation repaid. Zacchaeus gave half his wealth to the poor and those harmed by him were quadruply repaid. Sharing the wealth of meditation allows the church to be rewarded and enriched with the treasure of the wisdom which is hidden in Christ. This is also the duty of consecrated life. In this very same point of \textit{Presbyterorum Ordinis} that priests are not less so obligated to share the truths of the Gospel of which they are ministers. The Latin text of this document, in this place, uses the word \textit{debitores} -- debtors! The ordained priest is a debtor of the church, similar to one who has in his possession a key to a treasury, which is possessed by the whole church. What good is a house full of gold, if no one was able to open its doors? The key to the treasury, which is the treasure of all, is possessed by the monk and particularly by the priest. If they do not share any of the treasure, they become the debtors of their neighbours as wrote St John Maria Vianney. In spirit of the charism of our order, which harkens to the tradition of the Egyptian desert hermits, the preaching of the Word of God, heard in solitude, confronted in the community of the brethren, becomes the answer to the insistent demand of those pilgriming to the hermits: ``Abba, say to me a word!''
\medbreak
\newpage
\newpageheader{Latin Prayers}{Prayers}

\chapter{Latin Prayers}

\section{Prayers}

In nomine Patris et Filii et Spiritus Sancti. Amen.
\medbreak

\textbf{Pater noster}, qui es in caelis, sanctificetur nomen tuum, adveniat regnum tuum, fiat voluntas tua, sicut in caelo et in terra. Panem nostrum cotidianum da nobis hodie, et dimitte nobis debita nostra, sicut et nos dimittimus debitoribus nostris et ne nos inducas in tentationem, sed libera nos a malo. Amen.
\medbreak

\textbf{Ave Maria}, gratia plena, Dominus tecum, benedicta tu in mulieribus et benedictus fructus ventris tui Iesus.\\ Sancta Maria, Mater Dei, ora pro nobis peccatoribus, nunc et in hora mortis nostrae . Amen.
\medbreak

\textbf{Gloria Patri} et Filio et Spiritui Sancto, sicut erat in principio et nunc et semper et in saecula saeculorum. Amen.
\medbreak

\textbf{Credo in Deum}, Patrem omnipotentem, Creatorem caeli et terrae, et in Iesum Christum, Filium eius unicum, Dominum nostrum: qui conceptus est de Spiritu sancto, natus ex Maria Virigine, passus sub Pontio Pilato, crucifixus, mortuus, et sepultus: descendit ad inferos; tertia die resurrexit a mortuis; ascendit ad caelos; sedet ad dexteram Dei Patris omnipotentis: inde venturus est iudicare vivos et mortuos. Credo in Spiritum Sanctum, sanctam Ecclesiam catholicam, Sanctorum communionem, remissionem peccatorum, carnis resurrectionem, vitam aeternam. Amen.
\medbreak
\newpage
\newpageheader{Latin Prayers}{Morning Prayers}

\section{Morning Prayers}

\begin{samepage}
\prayer{\Vbar.}{\textbf{In nomine Patris et Filii et Spiritus Sancti.}}
\prayer{\Rbar.}{Amen.}
\end{samepage}
\medbreak
\begin{samepage}
\prayer{\Vbar.}{\textbf{Aperi, Domine,}}
\prayer{\Rbar.}{os meum ad benedicendum nomen sanctum tuum: munda quoque cor meum ab omnibus vanis, perversis et alienis cogitationibus; intellectum illumina, affectum inflamma, ut digne, attente ac devote hoc Officium recitare valeam, et exaudiri merear ante conspectum divinae Maiestatis tuae. Per Christum, Dominum nostrum. Amen.}
\end{samepage}
\medbreak

\label{morninglahtmlangelus}

\subsection*{Ordinary Time}

\begin{samepage}
\prayer{\Vbar.}{\textbf{Angelus Domini nuntiavit Mariae.}}
\prayer{\Rbar.}{Et concepit de Spiritu Sancto.}
\end{samepage}
\medbreak

\begin{addmargin}[6em]{0em}\noindent\textit{Ave Maria, gratia plena, Dominus tecum, benedicta tu in mulieribus et benedictus fructus ventris tui Iesus. }\end{addmargin}
\medbreak

\begin{addmargin}[6em]{0em}\noindent\textit{Sancta Maria, Mater Dei, ora pro nobis peccatoribus, nunc et in hora mortis nostrae. Amen. }\end{addmargin}
\medbreak
\begin{samepage}
\prayer{\Vbar.}{\textbf{Ecce ancilla Domini.}}
\prayer{\Rbar.}{Fiat mihi secundum verbum tuum.}
\end{samepage}
\medbreak

\begin{addmargin}[6em]{0em}\noindent\textit{Ave Maria\ldots }\end{addmargin}
\medbreak
\begin{samepage}
\prayer{\Vbar.}{\textbf{Et Verbum caro factum est.}}
\prayer{\Rbar.}{Et habitavit in nobis.}
\end{samepage}
\medbreak

\begin{addmargin}[6em]{0em}\noindent\textit{Ave Maria\ldots }\end{addmargin}
\medbreak
\begin{samepage}
\prayer{\Vbar.}{\textbf{Ora pro nobis Sancta Dei Genitrix.}}
\prayer{\Rbar.}{Ut digni efficiamur promissionibus Christi.}
\end{samepage}
\medbreak
\begin{samepage}
\prayer{\Vbar.}{\textbf{Oremus:}\\Gratiam tuam, quaesumus Domine, mentibus nostris infunde ut qui Angelo nuntiante, Christi Filii tui incarnationem cognovimus, per passionem eius et crucem ad resurrectionis gloriam perducamur. Per Christum, Dominum nostrum.}
\prayer{\Rbar.}{Amen.}
\end{samepage}
\medbreak

\label{morninglahtmlreginacaeli}

\subsection*{Easter}

\begin{samepage}
\prayer{\Vbar.}{\textbf{Regina caeli}}
\prayer{\Rbar.}{laetare, alleluia,\\ quia quem meruisti portare, alleluia.\\ Resurrexit, sicut dixit, alleluia.\\ Ora pro nobis Deum, alleluia.\\}
\end{samepage}
\medbreak
\begin{samepage}
\prayer{\Vbar.}{\textbf{Gaude et laetare, Virgo Maria, alleluia,}}
\prayer{\Rbar.}{Quia surrexit Dominus vere, alleluia.}
\end{samepage}
\medbreak
\prayer{\Vbar.}{\textbf{Oremus:}\\Deus, qui per resurrectionem Filii tui, Domini Iesu Christi, mundum laetificare dignatus es, praesta, quaesumus; ut, per eius Genitricem Virginem Mariam, perpetuae capiamus gaudia vitae. Per eundem Christum, Dominum nostrum.}
\prayer{\Rbar.}{Amen.}
\hrlinered\medbreak
\begin{samepage}
\prayer{\Vbar.}{\textbf{Gloria Patri et Filio, et Spiritui Sancto,}}
\prayer{\Rbar.}{sicut erat in principio, et nunc et semper et in saecula saeculorum. Amen.}
\end{samepage}
\medbreak
\begin{samepage}
\prayer{\Vbar.}{\textbf{Requiem aeternam dona eis Domine.}}
\prayer{\Rbar.}{Et lux perpetua luceat eis.}
\smallbreak
\prayer{\Vbar.}{\textbf{Requiescant in pace.}}
\prayer{\Rbar.}{Amen.}
\end{samepage}
\medbreak

\label{morninglahtmlmaryordinary}

\subsection*{Ordinary Time}

\begin{samepage}
\prayer{\Vbar.}{\textbf{Gaude, Dei Genitrix,}}
\begin{litany}\prayer{\Rbar.}{Virgo Immaculata.\par Gaude quae ab Angelo gaudium suscepisti\par Gaude quae genuisti Aeterni Luminis Claritatem.\par Gaude, Mater, gaude Sancta Dei Genitrix.\par Tu sola Mater es intacta.\par Te laudat omnis creatura, Genitrix Lucis!\par Sis pro nobis, quaesumus, perpetua interventrix\par ad Dominum Iesum Christum.}\end{litany}
\end{samepage}
\medbreak
\begin{samepage}
\prayer{\Vbar.}{\textbf{In omni tribulatione et angustia nostra. Alleluia.}}
\prayer{\Rbar.}{Succurre nobis, piissima Virgo Maria! Alleluia.}
\end{samepage}
\medbreak

\label{morninglahtmlmarylent}

\subsection*{Lent}

\prayer{\Vbar.}{\textbf{Salve quae immensum incomprehensumqur sancto virgineo utero comprehendisti,}}
\begin{litany}\prayer{\Rbar.}{per quam sancta Trinitas glorificatur et adoratur;\par per quam pretiosa crux celebratur, et in universo orbe adoratur;\par per quam coelum exsultat,\par per quam angeli rt archangeli laetantur,\par per quam daemones fugantur,\par per quam tentator diabolus coelo decidit,\par per quam prolapsa creatura in coelum assumitur,\par per quam universa creatura idolorum vesania detenta, ad veritatis agnitionem pervenit;\par pcr quam sanctum baptisma obtingit credentibus\par per quam exsultationis oleum\par pcr quam toto terrarum orbe fundatae sunt Ecclesiae\par per quam gentes adducuntur ad poenitentiam: et quid plura dicam?\par per quam unigenitus Dei Filius iis qui in tenebris rl in umbra mortis sedebant, lux resplenduit;\par per quam prophetae praenuntiarunt,\par per quam apostoli salutem gentibus praedicarunt\par per quam mortui exsuscitantur,\par per quam reges regnant, per sanctam Trinitatem.}\end{litany}
\medbreak
\begin{flushright}\red{St. Cyril of Alexandria Bishop. Omilia IV. P.G. LXXVII, 991}\end{flushright}
\medbreak
\hrlinered\medbreak
\begin{samepage}
\prayer{\Vbar.}{\textbf{In omni tribulatione et angustia nostra.}}
\prayer{\Rbar.}{Succurre nobis, piissima Virgo Maria!}
\end{samepage}
\medbreak
\begin{samepage}
\prayer{\Vbar.}{\textbf{Oremus:}\\Domine Deus Omnipotens, qui ad principium huius diei nos pervenire fecisti; tua nos hodie salva virtute, ut in hac die ad nullam declinemus peccatum, sed semper ad tuam iustitiam faciendam nostra procedant eloquia, dirigantur cogitationes et opera. Per Christum, Dominum nostrum.}
\prayer{\Rbar.}{Amen.}
\end{samepage}
\medbreak
\begin{samepage}
\prayer{\Vbar.}{\textbf{Angele Dei,}}
\prayer{\Rbar.}{qui custos es mei, me tibi commissum pietate superna illumina, custodi, rege et guberna. Amen.}
\end{samepage}
\medbreak
\begin{samepage}
\prayer{\Vbar.}{\textbf{Veni, Sancte Spiritus,}}
\prayer{\Rbar.}{reple tuorum corda fidelium et tui amoris in eis ignem accende.}
\smallbreak
\prayer{\Vbar.}{\textbf{Emitte Spiritum tuum et creabuntur. \rubric{(}Alleluia\rubric{)}.}}
\prayer{\Rbar.}{Et renovabis faciem terrae. \rubric{(}Alleluia\rubric{)}.}
\end{samepage}
\medbreak
\begin{samepage}
\prayer{\Vbar.}{\textbf{Oremus:}}
\prayer{\Rbar.}{Deus, qui corda fidelium Sancti Spiritus illustratione docuisti, da nobis in eodem Spiritu recta sapere, et de eius semper consolatione gaudere. Per Christum, Dominum nostrum.}
\prayer{\Rbar.}{Amen.}
\end{samepage}
\medbreak
\begin{samepage}
\prayer{\Vbar.}{\textbf{Regina Eremitarum.}}
\prayer{\Rbar.}{Ora pro nobis.}
\smallbreak
\prayer{\Vbar.}{\textbf{Mater Ordinis nostri.}}
\prayer{\Rbar.}{Ora pro nobis.}
\end{samepage}
\medbreak
\newpage
\newpageheader{Latin Prayers}{Marian Antiphons}

\section{Marian Antiphons}

Sub tuum praesidium, confugimus Sancta Dei Genitrix, nostras deprecationes ne despicias in necessitatibus nostris, sed a periculis cunctis libera nos semper, Virgo gloriosa et benedicta. Domina nostra, Mediatrix nostra, Advocata nostra, Tuo Filio nos reconcilia, Tuo Filio nos recommenda, Tuo Filio nos repraesenta.
\medbreak
\hrline\medbreak
\begin{litany}\prayer{}{Maria, Mater gratiae,\par Mater misericordiae,\par Tu nos ab hoste protege\par Et mortis hora suscipe.}\end{litany}
\begin{litany}\prayer{}{Jesu, Tibi sit gloria,\par qui natus est de Virgine\par Cum Patre et Almo Spiritu\par In sempiterna saecula. Amen.}\end{litany}
\hrline\medbreak
\begin{litany}\prayer{}{Salve Regina, mater misericordiae.\par Vita dulcedo, et spes nostra salve.\par Ad te clamamus exsules filii Hevae.\par Ad te suspiramus, gementes et flentes\par in hac lacrimarum valle. Eia ergo Advocata nostra,\par illos tuos, misericordes oculos\par ad nos converte. Et Iesum,\par benedictum fructum ventris tui, nobis\par post hoc exsilium ostende.\par O clemens, o pia, o dulcis Virgo Maria.}\end{litany}

\subsection*{Gaude Dei}

\subsubsection*{Melodia I.}

\begin{center}
\noindent\includegraphics[width=\textwidth,height=\textheight,keepaspectratio]{modern/lilypond/01/lily-9d5bb884-1}
\noindent\includegraphics[width=\textwidth,height=\textheight,keepaspectratio]{modern/lilypond/01/lily-9d5bb884-2}
\noindent\includegraphics[width=\textwidth,height=\textheight,keepaspectratio]{modern/lilypond/01/lily-9d5bb884-3}
\noindent\includegraphics[width=\textwidth,height=\textheight,keepaspectratio]{modern/lilypond/01/lily-9d5bb884-4}
\noindent\includegraphics[width=\textwidth,height=\textheight,keepaspectratio]{modern/lilypond/01/lily-9d5bb884-5}
\noindent\includegraphics[width=\textwidth,height=\textheight,keepaspectratio]{modern/lilypond/01/lily-9d5bb884-6}
\noindent\includegraphics[width=\textwidth,height=\textheight,keepaspectratio]{modern/lilypond/01/lily-9d5bb884-7}
\noindent\includegraphics[width=\textwidth,height=\textheight,keepaspectratio]{modern/lilypond/01/lily-9d5bb884-8}
\end{center}

\subsection*{Gaude Dei \\ {\small\normalfont (From the Cantuale of Innocent Burchaciński from the year 1820.)}}

\subsubsection*{Melodia II.}

\begin{center}
\noindent\includegraphics[width=\textwidth,height=\textheight,keepaspectratio]{modern/lilypond/63/lily-7eb3617c-1}
\noindent\includegraphics[width=\textwidth,height=\textheight,keepaspectratio]{modern/lilypond/63/lily-7eb3617c-2}
\noindent\includegraphics[width=\textwidth,height=\textheight,keepaspectratio]{modern/lilypond/63/lily-7eb3617c-3}
\noindent\includegraphics[width=\textwidth,height=\textheight,keepaspectratio]{modern/lilypond/63/lily-7eb3617c-4}
\noindent\includegraphics[width=\textwidth,height=\textheight,keepaspectratio]{modern/lilypond/63/lily-7eb3617c-5}
\noindent\includegraphics[width=\textwidth,height=\textheight,keepaspectratio]{modern/lilypond/63/lily-7eb3617c-6}
\noindent\includegraphics[width=\textwidth,height=\textheight,keepaspectratio]{modern/lilypond/63/lily-7eb3617c-7}
\noindent\includegraphics[width=\textwidth,height=\textheight,keepaspectratio]{modern/lilypond/63/lily-7eb3617c-8}
\noindent\includegraphics[width=\textwidth,height=\textheight,keepaspectratio]{modern/lilypond/63/lily-7eb3617c-9}
\noindent\includegraphics[width=\textwidth,height=\textheight,keepaspectratio]{modern/lilypond/63/lily-7eb3617c-10}
\noindent\includegraphics[width=\textwidth,height=\textheight,keepaspectratio]{modern/lilypond/63/lily-7eb3617c-11}
\end{center}

\subsection*{Omni Die}

\begin{center}
\noindent\includegraphics[width=\textwidth,height=\textheight,keepaspectratio]{modern/lilypond/82/lily-4c3d384d-1}
\noindent\includegraphics[width=\textwidth,height=\textheight,keepaspectratio]{modern/lilypond/82/lily-4c3d384d-2}
\noindent\includegraphics[width=\textwidth,height=\textheight,keepaspectratio]{modern/lilypond/82/lily-4c3d384d-3}
\noindent\includegraphics[width=\textwidth,height=\textheight,keepaspectratio]{modern/lilypond/82/lily-4c3d384d-4}
\noindent\includegraphics[width=\textwidth,height=\textheight,keepaspectratio]{modern/lilypond/82/lily-4c3d384d-5}
\noindent\includegraphics[width=\textwidth,height=\textheight,keepaspectratio]{modern/lilypond/82/lily-4c3d384d-6}
\end{center}

\subsection*{Maria Regina Mundi}

\begin{center}
\noindent\includegraphics[width=\textwidth,height=\textheight,keepaspectratio]{modern/lilypond/d9/lily-68cc2042-1}
\noindent\includegraphics[width=\textwidth,height=\textheight,keepaspectratio]{modern/lilypond/d9/lily-68cc2042-2}
\noindent\includegraphics[width=\textwidth,height=\textheight,keepaspectratio]{modern/lilypond/d9/lily-68cc2042-3}
\end{center}
\newpage
\newpageheader{Latin Prayers}{Litany to St. Paul the First Hermit}

\section{Litany to St. Paul the First Hermit}

\begin{litany}Kyrie, eleison. Christe, eleison. Kyrie eleison. \par
Christe, audi nos. Christe exaudi nos. \par
\end{litany}
\medbreak
\begin{litany}Pater de caelis, Deus, \responce{miserere nobis.}\par 
Fili, Redemptor mundi, Deus, \par
Spiritus Sancte,Deus, \par
Sancta Trinitas, unus Deus, \par
\end{litany}
\medbreak
\begin{litany}Sancta Maria, \responce{ora pro nobis.}\par 
Sancta Dei Genitrix, \par
Sancta Virgo virginum, \par
Sancte Michael, \par
Sancte Gabriel, \par
Sancte Raphael, \par
Sancte Angele Custos, \par
\end{litany}
\medbreak
\begin{litany}Omnes Sancti Angeli et Archangeli, \responce{orate pro nobis.}\par 
Omnes Sancti beatorum Spirituum Ordines, \par
\end{litany}
\medbreak
\begin{litany}Sancte Paule, Ordinis nostri Pater, \responce{ora pro nobis.}\par 
Sancte Paule, antesignane eremitarum, \par
Sancte Paule, pater anachoretarum, \par
Sancte Paule, vitae monasticae magister, \par
Sancte Paule, cultor solitudinis, \par
Sancte Paule, zelator perfectionis, \par
Sancte Paule, speculum pauperitatis, \par
Sancte Paule, humilitatis singulare ornamentum, \par
Sancte Paule, imitator vere S. Joannis Baptistae, \par
Sancte Paule, venerande senex, \par
Sancte Paule, serve fidelis Summi Regis, \par
Sancte Paule, nostrae militiae vexillum, \par
Sancte Paule, nostrae peregrinationis asylum, \par
Sancte Paule, horror et terror daemonum, \par
Sancte Paule, tentatorum refugium, \par
Sancte Paule, altissimae contemplationi dedite, \par
Sancte Paule, divino amore ardentissime aestuans, \par
Sancte Paule, a mundi turbinę sequestrate, \par
Sancte Paule, abstinentiae et inediae exemplar, \par
Sancte Paule, vas prophetici spiritus, \par
Sancte Paule, corvi ministerio vivens, \par
Sancte Paule, patientiae idea, \par
Sancte Paule, vir sine querela, \par
Sancte Paule, divi Antonii spectaculum, \par
Sancte Paule, angelicae vitae aemulator, \par
Sancte Paule, innocentiae conservator, \par
Sancte Paule, qui expansis manibus spiritum Deo reddidisti, \par
Sancte Paule, ministerio leonum post mortem honorate, \par
Sancte Paule, a magno Antonio tumulate, \par
Sancte Paule, ab Angelis i caelos delate, \par
Sancte Paule, Beatorum choris associate, \par
Sancte Paule, incola Ierusalem caelestis, \par
\end{litany}
\medbreak
\begin{litany}Ut nos votorum nostrorum memores eorundemquecompotes faciat Christus, \responce{impetra nobis gratiam.}\par 
Ut nos divitiarum et rerum mundanarum contemptores faciat Christus, \par
Ut nos mente et corpore castos faciat Christus, \par
Ut nos mortis et iudiciorum Dei perpetuo memores faciat Christus, \par
Ut nos veros morum tuorum et sanctitatis imitatores faciat Christus, \par
Ut abnegantes saecularia desideria sancte pieque in Ordine vivamus, \par
Ut numquam a vigore et rigore Regulae deficiamus, \par
Ut te duce impetus satanae, mundi carnisque illecebras retundamus, \par
Ut omnibus bonus odor Christi simus, \par
Ut Ordinem nostrum augere, promovere, virtutibusque suis Christus exornare dignetur, \par
Ut fratres Ordinis nostri defunctos ad aeterna gaudia perducat, \par
\end{litany}
\medbreak
\begin{litany}Per intercessionem beatissimae Virginis Matris tuae, \responce{Christe, te rogamus audi nos.}\par 
Per merita Sanctorum omnium et electorum tuorum, \par
\end{litany}
\medbreak
\noindent\begin{minipage}{\textwidth}
\begin{litany}Agnus Dei, qui tollis peccata mundi, \responce{parce nobis, Domine.}\par 
\end{litany}
\begin{litany}Agnus Dei, qui tollis peccata mundi, \responce{exaudi nos, Domine.}\par 
\end{litany}
\begin{litany}Agnus Dei, qui tollis peccata mundi, \responce{miserere nobis.}\par 
\end{litany}
\end{minipage}
\medbreak
\begin{samepage}
\prayer{\Vbar.}{\textbf{Ora pro nobis Sancte Paule prime Eremita.}}
\prayer{\Rbar.}{Ut digni efficiamur promissionibus Christi.}
\end{samepage}
\medbreak
\begin{samepage}
\prayer{\Vbar.}{\textbf{Oremus:}\\ Deus, qui beatum Paulum primum Eremitam mundi turbine segregatum Angelorum coetibus mirabiliter sociasti, ab omni nos, quaesumus, mundi purga contagio et caelestibus donis sanctifica. Per Christum, Dominum nostrum.}
\prayer{\Rbar.}{Amen.}
\end{samepage}
\medbreak
\newpage
\newpageheader{Latin Prayers}{Prayers in the Refectory}

\section{Prayers in the Refectory}

\instruction{Before Meals:}
\medbreak
\begin{samepage}
\prayer{\Rbar.}{\textbf{Benedic Domine nos et haec omnia dona, que de tua largitate sumus sumpturi. Per Christum, Dominum nostrum.}}
\prayer{\Vbar.}{Amen.}
\end{samepage}
\medbreak

\instruction{After Meals:}
\medbreak
\begin{samepage}
\prayer{\Rbar.}{\textbf{Gratias agimus Tibi, omnipotens Deus, pro universis beneficiis tuis: Qui vivis et regnas in saecula saeculorum.}}
\prayer{\Vbar.}{Amen.}
\end{samepage}
\medbreak

\instruction{or:}
\medbreak
\begin{samepage}
\prayer{\Rbar.}{\textbf{Retribuere dignare Domine, omnibus nobis bona facientibus propter nomen tuum vitam aeternam.}}
\prayer{\Vbar.}{Amen.}
\end{samepage}
\medbreak
\newpage
\newpageheader{Latin Prayers}{Blessing by a Bishop}

\section{Blessing by a Bishop}

\begin{samepage}
\prayer{\Rbar.}{\textbf{Sit nomen Domini benedictum.}}
\prayer{\Vbar.}{Ex hoc nunc et usque in saeculum.}
\end{samepage}
\medbreak
\begin{samepage}
\prayer{\Rbar.}{\textbf{Adiutorium nostrum in nomine Domini.}}
\prayer{\Vbar.}{Qui fecit caelum et terram.}
\end{samepage}
\medbreak
\begin{samepage}
\prayer{\Rbar.}{\textbf{Benedicat vos omnipotens Deus, Pater, et Filius, et Spiritus Sanctus.}}
\prayer{\Vbar.}{Amen.}
\end{samepage}
\medbreak
\newpageheader{Latin Prayers}{Blessing by a Priest}

\section{Blessing by a Priest}

\begin{samepage}
\prayer{\Rbar.}{\textbf{Dominus vobiscum.}}
\prayer{\Vbar.}{Et cum spiritu tuo.}
\smallbreak
\prayer{\Rbar.}{\textbf{Benedicat vos omnipotens Deus, Pater, et Filius, et Spiritus Sanctus.}}
\prayer{\Vbar.}{Amen.}
\end{samepage}
\medbreak
\newpage
\end{appendices}
